\documentclass[book, paper=a4paper, 11pt]{jlreq}
\usepackage[dvipdfmx]{graphicx,xcolor,hyperref}%graphics
\usepackage{float}
\usepackage{amsmath}
\usepackage{amsthm}
\usepackage{verbatim}
\usepackage{latexsym}
\usepackage{amssymb}
\usepackage{dynkin-diagrams}
\usepackage{ccicons}
\usepackage{listings}
\usepackage[basic, box,gate,oldgate, ic, optics,physics]{circ}
\usepackage{ulem}
\hypersetup{% hyperrefオプションリスト
   setpagesize=false,
   bookmarksnumbered=true,
   bookmarksopen=true,
   colorlinks=true,
   linkcolor=blue,
   citecolor=purple,
} 
\usepackage{mflogo}
%\usepackage{emath}
\usepackage{MePoTeX}

\input xy
\xyoption{all}

\def\emptyline{\vspace{12pt}}
\newcommand*{\term}[2][]{\textcolor{red}{#2}%
	\ifx\relax#1\relax\index{#2}\else\index{#1@#2}\fi}

\usepackage{makeidx}
\makeindex

\newtheorem{PP}{授業の目標}
\newtheorem{Thm}{定理}
\newtheorem{Prop}{命題}
\newtheorem{Q}{演習}
\newtheorem{Quiz}{問題}
\theoremstyle{definition}
\newtheorem{Def}{定義}

\theoremstyle{remark}
\newtheorem{ex}{例}
\newtheorem{exQ}{例題}

\DeclareMathOperator{\Ker}{Ker}
\DeclareMathOperator{\Image}{Im}
\DeclareMathOperator{\stab}{stab}
\DeclareMathOperator{\sign}{sign}
\DeclareMathOperator{\id}{id}
\DeclareMathOperator{\Is}{Is}
\DeclareMathOperator{\tr}{tr}
\DeclareMathOperator{\vol}{vol}
\DeclareMathOperator{\const}{const}
%\DeclareMathOperator{\GL}{GL}
%\DeclareMathOperator{\SL}{SL}
%\DeclareMathOperator{\PSL}{PSL}
%\DeclareMathOperator{\Ot}{O}
%\DeclareMathOperator{\SO}{SO}
%\DeclareMathOperator{\U}{U}
%\DeclareMathOperator{\SU}{SU}
%\DeclareMathOperator{\Sp}{Sp}
\DeclareMathOperator{\IM}{Im}
\DeclareMathOperator{\RE}{Re}
\DeclareMathOperator{\Res}{Res}
\DeclareMathOperator{\arcSin}{Sin^{-1}}
\DeclareMathOperator{\noT}{NOT}
\DeclareMathOperator{\anD}{AND}
\DeclareMathOperator{\oR}{OR}
\DeclareMathOperator{\nand}{NAND}
\DeclareMathOperator{\nor}{NOR}
\DeclareMathOperator{\xor}{XOR}
\DeclareMathOperator{\idm}{ID}

\def\emptyline{\vspace{12pt}}
\def\smallemptyline{\vspace{6pt}}
\setlength{\belowcaptionskip}{0mm}

\title{情報とプログラミング\par \normalsize{---歴史と数学からの入門---}}
%\title{Introduction to Information and Programming\par \normalsize{---An Introduction from History and Mathematics---}}
\author{snow00two}

\begin{document}
\maketitle

\frontmatter

\chapter{イントロダクション}

この本は理系全般，特に数学に興味のあるの大学生・社会人に向けて
情報学とプログラミングの入門となることをまとめた．
特に情報学とプログラミングの歴史に関係したことと
そこに表れる数学に関連したこととから説き起こし，
すでに知っていると思われる高校の情報Iのレベルより
更に高いレベルへの入門を目指した．

このような本をまとめた背景を述べたい．
学生に接しているとCLI=Command Line Interfaceを触ったことがなく，
開発環境をローカルに構築したことがない，VSCodeすら触ったことの無い
学生が多数いることがわかる．
またコロナ禍以降，スマートフォン中毒者もよく見かけるようになった．
スマートフォンは時間を意味なく消費し，
目や耳や首や肩などを始めとして体の不調を引き起こし，
脳の発達を邪魔し，アルゴリズムが世界観をも壊している．
更に，これからは生成AIに依存・代替される人間が増えることが予想されている．

これらに対応するため，これからの学生・社会人には
\begin{itemize}
   \item
      スマートフォンとインターネットと生成AIに適切な距離を置く
   \item
      CLIを使い(\textcolor{red}{生成AIは既に使いこなしている}ので，
      逆に生成AIを使いこなすためには必須である)こなす
   \item
      必要に応じて自力で開発環境を構築し，ツールがつかえるようになる
   \item
      幾つかのプログラミング言語でプログラミングができる
   \item
      生成AIにMarkdown などを上手く使って適切に構造化した指示が出来る
   \item
      生成AIにより拡散された偽情報を見抜ける
\end{itemize}
ことが大切になると思われる．

そこでこの本では
CLIの演習を行い，
デジタルデータの仕組みを述べ，取り扱いを説明する．
次に論理回路やチューリングマシンなど計算論の基礎への説明を与え，
インターネットやWEBの仕組みやセキュリティーについて説明する．
最後にはデジタルデトックスについても論じている．
付録には具体的で身近なプログラミング言語
(例えばExcel/VBA, \LaTeX, markdown, JavaScript, Shell Script, git)
のプログラミング演習となるような資料を幾つか上げたので，
成るべく多くのものにチャレンジし，技術を深掘りしてほしい．
また，演習問題も多数あげたので自分の理解の確認に使って貰えると有り難い．
授業のレポートで友達のものやネットに上がっているものを単にコピペしても
理解も技術も深まらないように，生成AIの答えを貼り付けても理解も技術も深まりません．
それらを参考にしても構わないが，大凡のアイデアを理解したら，
自分なりにプログラミングの本質を理解し，アイデアを検討しなおし，
適当な関数にまとめるとか，クラスにまとめるとかして下さい．
コードが読みやすいようにソースを適切にまとめ，分かり易い変数の名前を選ぶ
とか，適切なコメントを入れるとかしてください．
そうすれば理解も技術も深めて自分のものにすることができる．
%生成AIについてもできる限り取り上げている．
全体を合わせると，
上に上げた項目に対処できる力に一歩でも近づけるものが増えると期待している．

なおこの本では MacOSの英語表示を使って説明をしている．
Windows でも大体は同じことが成り立つが，
慣れないと対応が分かりづらい部分もあると思う．
ネット等で調べたり生成AIに聞いて各自対応してほしい．

この本の図のソースファイルや関連した動画やアプリを
\href{https://snow00two.github.io/mySite/}{私のサイト}に用意したので
それらを動かしながら理解を深めてもらいたい．

\tableofcontents

\mainmatter

%%\documentclass[a4paper,11pt]{jlreq}
%\documentclass[a4paper,12pt]{amsart}
%\usepackage[dvipdfmx]{graphicx,xcolor}%graphics
%\usepackage{float}
%\usepackage{amsmath}
%\usepackage{verbatim}
%\usepackage{latexsym}
%\usepackage{amssymb}
%\usepackage{qrcode}
%%\usepackage{emath}
%\usepackage{listings}
%\usepackage{ccicons}
%\usepackage{mflogo}
%%\usepackage{fancybox}%enclose by box
%\usepackage[dvipdfmx]{hyperref}
%\usepackage{verbatim}
%\hypersetup{% hyperrefオプションリスト
%   setpagesize=false,
%   bookmarksnumbered=true,%
%   bookmarksopen=true,%
%   colorlinks=true,%
%   linkcolor=magenta,
%   citecolor=blue
%} 
%\newtheorem*{exQ}{レポート}
%\theoremstyle{remark}
%\newtheorem*{Q}{演習}
%
%\renewcommand{\figurename}{図}
%\renewcommand{\proofname}{証明}
%\renewcommand{\tablename}{表}
%
%\def\emptyline{\vspace{11pt}}
%\def\smallemptyline{\vspace{6pt}}
%\begin{document}
%\begin{flushright}
%   2026/04/14
%   %2026/04/04
%   %2026/03/29
%   %2026/03/11
%   %2026/02/21
%   %2026/02/10
%   %2025/12/22
%   
%   九大・数理 大津幸男
%\end{flushright}
%
%\begin{center}
%   {\Large 情報科指導法II---01}
%   \href{https://creativecommons.org/licenses/by-nc-nd/4.0/}{\ccbyncnd}
%   \par{Command Line Interface}
%\end{center}
%
%\section{授業目標と計画}
%この授業は毎回の講義と演習から成る．
%学期前半の授業では
%各授業を2つに分け，その前半では情報科で知っておくべき基礎的なトピックを講義形式で解説する．
%後半ではそれに関する環境構築に関する技術的なことやプログラミングやプロンプトエンジニアリングを演習形式で体験してもらう．
%%参考文献として \cite{Mo1}, \cite{Mo2}, \cite{K}, \cite{F}, \cite{KP} などを挙げておく．
%
%学期後半の残りの授業でも
%各授業を2つに分け，
%前半では生成AI の原理・プログラミングとプロンプトエンジニアリングとの関係について講義し，
%後半ではデジタルや生成AIと個人及び社会との関係に付いて，
%様々なCLツールと生成AIを使用しながら検討し，
%それを元にディベート形式で議論を深める演習を行う．
%
%この講義の重要な部分は演習なので出席をとります．
%出席して演習を行ってください．
%
%\section{情報科の参考資料}
%以下の文科省の発表した情報科の参考資料をこの授業でも参考にするので，
%ダウンロードして時間があるとき目を通しておいてください
%(このpdf ではマゼンタの文字の部分がHyperLink)．
%
%\begin{itemize}
%   \item
%      \href{https://www.mext.go.jp/a_menu/shotou/new-cs/1407074.htm}
%      {高校学習指導要領}
%   
%   \item
%      情報科編（上のリンクから辿れる）
%   
%   \item
%      \href{https://www.mext.go.jp/a_menu/shotou/zyouhou/detail/mext_00017.html}
%      {「情報I」教員研修用教材}
%   
%   \item
%      \href{https://www.mext.go.jp/a_menu/shotou/zyouhou/detail/mext_00018.html}
%      {「情報II」教員研修用教材}
%   
%   \item 
%      \href{https://www.mext.go.jp/b_menu/shingi/chukyo/chukyo3/118/mext_00002.html}{次期指導要領についての最近の文科省の動向}
%\end{itemize}
%
%この講義の準備のために
%\cite{F}, \cite{Fu}, \cite{K}, \cite{OM},  \cite{Sa}, \cite{Yag}, \cite{IT}\footnote{ITパスポートは情報処理推進機構IPAが実施する試験, 他にも基本情報技術者試験, 応用情報技術者試験がある．}
%などを参考にさせてもらった．
%日米の高校生・大学生・社会人に求められる情報リテラシーを知るのに有用だと思うので，
%手に入れる必要は無いが，いくつか手に取って眺めてみることをお奨めしたい．
%
%
\chapter{CLI入門}
\section{コンピューターの歴史}
コンピューターの歴史を簡単に振り反ってみる：
\begin{itemize}
   \item
       19世紀半ばにバベッジ(Babbage)による階差機関・解析機関の構想があった．
       計画は現代のコンピューターに近かったが，
       開発には成功しなかった．
       \href{https://www.youtube.com/watch?v=0anIyVGeWOI}
       {近年実現された階差機関}
 
   \item
       19世紀後半にはアナログコンピューターが登場した．
       例へば1930年頃のバネバー・ブッシュ(Vannevar Bush) の微分解析機がある．
   \item
      19世紀末にホレリス\footnote{のちのIBM( International Business Machines )の創業者の一人}(Hollerith)
      により開発されたパンチカードを使った集計で
      合衆国の国勢調査が行われた．
   \item
      1936年チューリング(A.~Turing)がチューリングマシンの概念を発表し，
      計算可能性について考察を行った(\cite{Turing})．

   \item
      最初の電子計算機ENIACがモークリー(J.Mauchly)とエッカート(J.P.Eckert)により1943〜1946に開発された．
      
   \item
      1951年フォン・ノイマン(von Neumann)による史上初の
      ストアドプログラム(stored program)方式の
      コンピューターEDVACが開発された．
      現在のコンピューターは全てこのノイマン型の子孫である．
      そのころの入力法は
      \href{https://www.youtube.com/shorts/DveA0EuI5qc}{紙テープ}や
      \href{https://www.youtube.com/watch?v=0065zKzeCWI}{パンチカード}だった．
   
   \item
      最初は機械語でプログラムされていたが．
      \href{https://www.youtube.com/watch?v=xdFfUmUJZ-w} 
      %\href{https://www.youtube.com/watch?v=e2LiY4ou3VI&t=200s}
       {グレース・ホッパー}(Grace Hopper)によりコンパイラーが開発され，
      アセンブラでコーディングされるようになり，
      更にFORTRANやCOBOLなどの高級言語が生まれた
      \footnote{FORTRAN(formula translator)は
      IBM, John Warner Backus により1957,
       COBOL(Common Business Oriented Language)は 
       G.~Hopper, CODASYL(データシステムズ言語協議会)により
       1959に開発された(\cite{W}). 共に現役で使われている言語である．}．
   
   \item 
      Operating systemの発達：
      最初は
      オペレーターと呼ばれる人が
      順番にまとめてプログラムを読み込んで実行する
      バッチ処理(batch processing)の形で行われていた．
      その後1964年に
      オペレーターの代わりをプログラム自身で行う
      Time sharing system の
      \href{https://www.youtube.com/watch?v=Q07PhW5sCEk&list=PLT8Ms9NT2xpb6yBvtVpBXzUoIMAhDu0G9}
      {CTSS}(Compatible Time Sharing System) が生まれ，
      テレタイプやキーボードからモニターを見ながら
      コンピューターに命令するというインタープリターを使った
      使用が盛んになり，そこからそのプログラムは
      OS (Operating System)と呼ばれるようになった．
      数々の商用OS(例
      \href{https://www.youtube.com/watch?v=npgvV_-Nh60}{IBM System360})
      が発売され
      Multics の失敗ののち，
      特に
      
      UNIX\footnote{UNIXはX/Open Company Ltd.~の登録商標. ただし，カーニハンは Unix と綴ることを勧めている(\cite{K}).}とC言語と shell の誕生：
      \begin{itemize}
         \item
            1969\footnote{この年インターネットの最初の実験と
            アポロ11号による最初の月着陸が行われた．} %\href{https://www.youtube.com/watch?v=CknIJHtxxsw}
            AT\&TのBell 研究所の\href{https://www.youtube.com/watch?v=HADp3emVABg&t=302s}{ケン・トンプソン}(Ken Thompson)
             によりPDP-7上でアセンブラでUNIXの開発が始まった
         \item
            1970 PDP-11 に移植(図~\ref{PDP11})
         \item
            1973 \href{https://www.youtube.com/watch?v=RaotHMNZPaI}
            {デニス・リッチー}(Dennis Ritchie)は
            ケン・トンプソンの開発したB言語をもとにC言語を開発し，
            ケン・トンプソンと共にC言語でUnixを書き直した．
            始めての高級言語によるOSの開発だった．これにより
            OSの移植性が格段に高まった．
            
          \item
            1991 リーナス・トーバルズ(Linus B.~Torvalds) が Linux を公開し，
            今でも開発が続いている．
            皆さんが使っているアンドロイド携帯や
            クラウドのコンピューターの多くはLinuxベースである\footnote{Linuxの重要な判断の最終決定はトーバルズが行っている．そのためトーバルズに何かがあったとき，コンクラーベが行われることになっている}．
            

      \end{itemize}
      
       \begin{figure}[htb]
         \centering
         \includegraphics[width=8cm]{~/Documents/GitHub/introInfo/ch1/images/Screen_Shot_2020-09-19_at_7.16.21_AM_big.png}
         \caption{Ken Thompson (sitting) and Dennis Ritchie working together at a PDP-11, タイプライターのようなものが\href{https://www.youtube.com/watch?v=I_O-KlHhkII}{テレタイプ}(Wikipedia より)}\label{PDP11}
      \end{figure}
 
   
   \item 1970年代GUI(Graphical User Interface)が開発された(\S \ref{GUI})．
   
   
   \item
      個人的には嫌だが，BCI(ブレイン・コンピューター・インタフェース）
      の時代が来るかも知れない．

\end{itemize}


\section{GUIの歴史}\label{GUI}
\begin{itemize}
   \item 
   第二次大戦後の1945年バネバー・ブッシュは``As we may think'' でMEMEX の構想を述べた．
   
   \item 
      1963年サザーランドがスケッチパッドを開発
   \item 
      1960年代中期にエンゲルバートがマウスを開発
   \item 
      1974 年Xerox のPARC (Plo Alto Research Center) でALTOが作られ，
      アラン・ケイ(Alan Kay)が1977年に
      \href{https://www.youtube.com/watch?v=BNmPQQwE50U}{ダイナブック}
      の構想を発表した．
     最初の \href{https://www.youtube.com/watch?v=H1GJMoZqLdo}
     {オブジェクト指向言語}の一つのSmallTalkやMulti-Windowシステムも開発された．
   
   \item
      \href{https://www.youtube.com/watch?v=pQocN_c2uLI}{PARC の訪問}に刺激をうけたスティーブ・ジョブス(Steven Jobs)による
      Lisa 開発と1983年発売，失敗
   
   \item 
      Lisa 開発から外れたスティーブ・ジョブスによる
      Macintoshが
      \href{https://www.youtube.com/watch?v=ErwS24cBZPc}{1984年}\footnote{「1984年」は現代社会を予言したジョージ・オーウェルのディストピア小説\cite{Orwell}．
         このリンク先はリドリー・スコット監督が1983年に発表したマッキントシュのコマーシャルである\cite{Jobs}. 
      }
      発売され
      商業的に成功し現在に繋がる．
   
\end{itemize}
コンピューターの発展の詳しい歴史については \cite{Jobs}, \cite{interface}, \cite{history}, \cite{UNIX}, \cite{Kay}, \cite{tools} などを参照してほしい．YouTube などで検索すると興味深い動画も多数見つかる．
インターネットの発達と生成AIの発達については後日解説する．

\section{CLI}
現在ではGUI で色々なことができるようになり，
実際ほとんどの人がGUI を使っていると思うが，
Command Line Interface(CLI, CUI =Character User Interface と呼ばれることもある) でないと出来ないことも多い．
何が出来て何が出来ないかはCLIを使ってみたことがないと分からない
と思うが，\textcolor{red}{生成AIもCLIを使いこなしている}ので
生成AIが何をしているのか分かるぐらいまではCLIを使いこなしてほしい．

ここでは CLI のOSの例としてUnixを取り上げ
(MacOSはUnixの兄弟FreeBSDを基に作られていて，
Linux もUnixの親戚に当たる), 
まずUnixのファイルシステムについて紹介する． 
shell の使い方については次のセクションで説明する．
\begin{itemize}
   \item
      ディレクトリー：% DOM
      \begin{figure}[htb]
         \centering
         \includegraphics{~/Documents/GitHub/introInfo/ch1/images/figure.1}
         \caption{ファイルシステムのツーリー構造}\label{tree}
      \end{figure}
   \begin{figure}[htb]
         \centering
         \includegraphics[width=12cm]{~/Documents/GitHub/introInfo/ch1/images/column.pdf}
         \caption{Column 表示でのツーリー構造の確認}\label{column}
      \end{figure}
            
         ディレクトリーやファイルは図~\ref{tree}のようなツリー構造に配置されている．
         特にパソコンのハードディスクの一番上の部分をルートといい \verb/// と表す．
         全てのディレクトリーやファイルはルートから一意的にたどっていく
         ことが出来る．
         ファイルやディレクトリが\verb/// から見てどこにあるかを示すのを
         \term[]{絶対path}という．
         例えば
         図~\ref{tree}のprogramming.texは
         
         /Users/sonw00two/Documents/classRoom/General/
         
         Programming/src/programming.tex
         
         \noindent
         のように表される．
         後で出てくる .. は一つ上のディレクトリーをあらわす\term[]{相対 path}の例

      \item
      不可視ファイル：
         ディレクトリ /usr  は
         普通には見えない設定になっているので
         知らないとアクセスできない．
         Finder のメニューバーの Go のメニューから Go to Folder を選ぶと
         四角いウィンドウが開くので /usr と入力し，goをクリックで移動できる．
         Finderで cmd shift G とタイプしてもよい．
         
   \item
         拡張子：%とpermission 
         ファイル名には hoge.tex の .tex のような
         \term[]{拡張子}(extension)があるが，
         多くの場合アイコンのみ表示され見えないようになっている．
         開発のときは様々なファイルを扱うので，
         拡張子を表示しないと効率的なファイルの管理が難しい．
         また，メールの添付ファイルのアイコンは偽造できるが
         拡張子により直ぐに峻別できるので，
         セキュリティーの観点からも表示するのが望ましい．
   
   \item        
         関係したファイルは一つのディレクトリーに，
         関連していないファイルは異なるディレクトリーにいれ，
         関連のある，あるいは類似のディレクトリーは纏めて一つの大きなデレクトリー
         に入れるなど入れ子構造にして（それがツリー構造の定義である）
         管理するのが基本である．
         また，簡単にそのディレクトリーを説明するテキストファイル ReadMe.txt, 
         README.mdをそのディレクトリーに置くことも良い習慣である．
         

%   \item
%
%
%shell 
%
%root 
%
%ログインディレクトリー
%作業ディレクトリ
%
%ファイル管理
%原則
%Login

\end{itemize}

\section{CLI 実践演習}
\begin{center}
   \href{https://www.youtube.com/watch?v=6XvmhE1J9PY}{オバマ元大統領からのメッセージ}
\end{center}

以下のアプリのインストール・色々なサービスへ登録してもらい，コマンドラインの使い方を説明する．
ただし，CLI では破滅的なことも出来るので，
バックアップは自分でこまめに取っておくようにして下さい．
\textcolor{red}{
   プログラミングの演習で生じた被害への補償は一切行えません．
   この免責事項を認めた人だけがこの本を読み続けることが出来ます．
}

  \begin{figure}[htb]
         \centering
         \includegraphics[width=12cm]{~/Documents/GitHub/introInfo/ch1/images/vscode.pdf}
         \caption{VSCode, 下が terminal}\label{vscodeTerm}
   \end{figure}

\begin{Q}
\begin{enumerate}
   \item
      VSCodeのインストール：
       
                   フリーのエディタというより統合開発環境
            (IDE = Integrated Development Environment)Visual Studio Code
            （略称VSCode, \href{https://code.visualstudio.com}{ここからダウンロード}
            可能，図 \ref{vscodeTerm}参)．
            様々なコンピューター言語のための（自動補完，構文解析のための）
            インテリセンス，リンター，スニペット，ユニットテスト
            \footnote{詳しい説明はネット等で調べてください.}等の機能拡張がある．
            マイクロソフトのVisual Studio  シリーズのエディタのみをフリーで公開した
            もので，他にも似たものが色々あったが，最近はデファクトスタンダード．
            JavaScriptで作られたアプリの一つで
            Chromium+JavaScript+Node.js+Electron 
            により作られている(図 \ref{vscode})．
     
      \begin{figure}[htb]
               \centering
               \includegraphics[width=6cm]{~/Documents/GitHub/introInfo/ch1/images/ScreenCapture2.pdf}
               \caption{VSCodeのクレジット}\label{vscode}
   \end{figure}
   
   \item
         terminal /shell:
       
       Macや Linux 等では既にインストールされている．
       Windows ではPowershell あるいは
       WSL(Windows Subsystem for Linux)をインストロールする必要があるだろう．
       
      これから必要になるtool等もCLIと上でも述べたtoolなどを使って
       インストールする必要のあるものが多い．
       Mac の標準の terminal は
      
       /Applications/Utilities/Terminal.app
   
       \noindent
       にある．
      Mac なら (Home)brew もインストールしておくと都合が良い．
      \href{https://brew.sh/ja/}{ここからダウンロード}
       
       不可視なディレクトリやファイルも沢山あるので
       適宜コマンドで確認する必要がある．
      VSCode の下の部分にもterminal が付いている（図 \ref{vscodeTerm}）．

%     \item
%      Web 技術：JavaScript ( TypeScript )
%      
%      \begin{enumerate}
%         \item  http とコマンドラインによるbrowserもどきの作成
%         
%         \item 
%              Chrome のDeveloper tool, Chromeには最初からDeveloper menu がある.
%             Developer ニューから JavaScript Console を開くと
%             Console が開き，デバッグが可能．
%         
%         \item  
%            JXA : JavaScriptが Mac のスクリプト言語として採用された：
%         
%            Script Editor・ターミナルから osascript・automator で使える．
%         
%         \item 
%            Single page application: JavaScript+Node.js+Electron etc.+Chromium でmacOS, 
%            Windows, スマートホンのいずれにも対応したデスクトップアプリケーションが
%            一度に作れる． 例VSCode 
%         
         
%         \item
%            p5.js ：
%            JavaScript で記述し，2D, 3Dの映像・動画の作成可能．
%            Processing というJavaのライブラリーをJavaScript に移植したもの
%            (\cite{P5})．
%      \end{enumerate}
   \item
      version 管理の基本ソフト git とネットサーバー GitHub. 
      多人数で開発を行うにはこれを使いこなさないといけないが，
      個人開発でも有用．
      GitHubで VSCode の設定ファイルの同期も出来る．
      大学等でMicrosoft365に契約していると，そのアドレスでGitHubに
       \href{ https://github.com/}{登録}し
      Student Developer Pack を使うように設定できるかもしれない．

   \item
      大学等でMicrosoft365に契約していると，
      ログインすると Copilot（chatGPTを office suite 向きにアレンジしたもの）が
      使用料無しで使えるようになっているかもしれない．
      GitHub のStudent Developer Pack に入っていると，
      VScode の拡張機能タブで「GitHub Copilot」を検索してインストールすると，
       使用料無しでVScodeからGitHub Copilot pro が使えるようになるかもしれない．
%   chatGPT etc
%      \begin{enumerate}
%      \item
%         機械学習
%      \item
%         プロンプトエンジニアリング(\cite{OT})．
%      \end{enumerate}

   \item
      JavaScript は  Chrome などのBrowser で既に使えるようになっているが，
      セキュリティのためディスクへのアクセスなどは制限されているので，
      強力な実行環境となる node.js とその管理ツール npm も必要になる．
      \href{https://nodejs.org/ja/download}{ここからダウンロード}.
      これからは
      \href{https://github.blog/news-insights/octoverse/octoverse-a-new-developer-joins-github-every-second-as-ai-leads-typescript-to-1/}
      {TypeScript}にも慣れる必要があるだろう．
            
   \item
      Python3 とその管理ツール pip や  (Ana)conda  も必要
      \href{https://www.python.org}{ここからダウンロード}
   
\end{enumerate}
\end{Q}


\begin{Q}
   拡張子を表示するようにパソコンを設定して下さい．
   \begin{figure}[h]
         \centering
         \includegraphics[width=8cm]{~/Documents/GitHub/introInfo/ch1/images/ScreenCapture1.pdf}
         \caption{Macでのエクステンションの表示}\label{extension}
      \end{figure}
\end{Q}

\begin{Q}
   VSCode の設定を，必要なら検索かCopilotに聞いて，
   GitHubで同期するように設定して下さい．
\end{Q}

\begin{Q}
   VSCode でGitHub Copilot を立ち上げ，
   （もし答えを既によく理解していれば簡単だが，
   その分野の知識がないと難しい）
   以下のような質問をいくつか行って下さい．
   \begin{enumerate}
      \item
         大学の代数の数学者として以下の質問に
        数学科の学生に分かるように答えて下さい．
         標数$p$の有限体が素体$\mathbf Z/(p)$のガロア拡大であることを示し，
         そのガロア群がどうなるか教えて下さい．
         
      \item
         大学の幾何の数学者として以下の質問に
         数学科の学生に分かるように答えて下さい．
         実射影平面$\mathbf R\mathbf P^2$($2$次元の実射影空間)の定義を述べ，
         それが閉多様体であることを示し，
         その$\mathbf Z$係数ホモロジーがどうなるか教えて下さい．
   \item
         大学の情報の専門家としてcommand line interface に不慣れな学生に，
         git, GitHub のcommand line interfaceによる管理の仕方を
         説明して下さい．
         
   \item
         最近 LLMや生成AIを勉強していて，
         それ自体はよく出来ていると思いますが，
         この方向でAGIがあと数年で出来るとも思えません．
         しかし生成AIを開発している会社のCEO達は「あと数年で
         AGI が出来そうだ」などと言っています．
         人工知能と経営学のそれぞれの専門家の立場から，
         なぜそのような言動を繰り返すのか答えてください．
         
      \item
         大学の経済学の専門家としてトランプ大統領にも分かるように，
         輸入した物に関税をかけると物価が上がることを説明して下さい．
         
   \end{enumerate}
   その生成AI の回答を読んで，
   （自分の知識の無い分野のその答の）内容が理解できたか確認して下さい．
\end{Q}


      コマンドラインインターフェースを使いこなさないとプログラミングは難しい．
        以下に簡単なコマンドの使用例をあげる．terminal(端末) を立ち上げるとプロンプト
        が表示される(図~\ref{terminal})．
        私の端末では$\cdots$\verb/% /となっているのがプロンプトで，
        そのままキーボードからコマンド（文字）を打ち，リターンキーを押すと
        そのコマンドが実行され，次の行に結果が表示される．
        \begin{figure}[htb]
         \centering
         \includegraphics[width=12cm]{~/Documents/GitHub/introInfo/ch1/images/terminal.pdf}
         \caption{Terminal}\label{terminal}
      \end{figure}
       これがコマンドラインインタープリター shellである．
       以下では \verb/% / でプロンプトを表す．
        
        以下\# 以降はコメント
        %以下，表示がインデントされすぎるのでインデントを減らす．
   \begin{lstlisting}
   % cd #  ホームディレクトリへ移動
   \end{lstlisting}
   \begin{lstlisting}
   % cd .. #  一つ上のディレクトリへ移動 
   \end{lstlisting}
   \begin{lstlisting}
   % pwd  #  現在のディレクトリの名前の表示
   \end{lstlisting}
   /Users
   \begin{lstlisting}
   % ls   #  現在のディレクトリにあるディレクトリやファイル表示
   \end{lstlisting}
   Shared   sonw00two
   \begin{lstlisting}
   % ls -a  #  不可視なものも含めたディレクトリやファイル表示
      localized   Shared   sonw00two
   % cd sonw00two  # このディレクトリのディレクトリ sonw00two へ移動
   \end{lstlisting}\begin{lstlisting}
   % mkdir tmp  # ディレクトリ tmp の作成
   \end{lstlisting}
   \begin{lstlisting}
   % cd tmp  # ディレクトリ tmp へ移動
   \end{lstlisting}
   \begin{lstlisting}
   % ps aux >ps.txt  # 全てのプロセスを ps.txt に書く
   \end{lstlisting}
   \begin{lstlisting}
   % more ps.txt  # ps.txt  を表示．q で表示を閉じる
   \end{lstlisting}
   \begin{lstlisting}
   % mv ps.txt ps1.txt  #  ps.txt の名前をps1.txt へ
   \end{lstlisting}
   \begin{lstlisting}
   % cp ps1.txt ps2.txt  # ps1.txt  をps2.txt にコピー
   \end{lstlisting}
   \begin{lstlisting}
   % rm *  # このディレクトリの任意のファイルを削除．*   はこのディレクトリの任意のファイル名を表すいわゆるワイルドカード．また，ファイルを全部消さないとディレクトリは削除できない．
   \end{lstlisting}
   \begin{lstlisting}z
   % cd ..  
   \end{lstlisting}
   \begin{lstlisting}
   % rmdir tmp  # ディレクトリ tmp の削除 % rm -R tmp で中のファイルを消さなくても削除可能ではあるが慣れるまでは確認してから削除を勧める．

   \end{lstlisting}
\emptyline

shellには全てをキーインしなくても，ファイル名などをある程度打ってタブキーを打つとあとの部分を補完してくれる機能がある．他の開発環境でも異なるキーかもしれないが大抵補間機能がついている．

   
%   \item
%      正規表現：パターンマッチングの表現．
%      どれでも良いのでよく使う言語の正規表現をマスターすることをお奨めする．
%      すると少しの変更で他の言語でも使えるようになる．
\begin{Q}
   コマンドラインを使い，ディレクトリーを移動したり，
   ディレクトリーやファイルの作成や削除を行って下さい．
\end{Q}

授業の中で出すレポートの幾つかを解いて
提出してください．最終成績はその出来で判断します．
例えば
\begin{exQ}
   自分のポートフォリオサイトを作って自分の
   （プログラムや動画や論文などの）作品を公開して下さい．
   そのURLを送って貰えるとその出来で成績を付けます．
   ちなみにこの講義のために作った
   \href{https://snow00two.github.io/mySite/}{私のポートフォリオ}です．
\end{exQ}

\begin{exQ}
   この講義を準備していて，
   講義の準備をするのが知識を深めるのに一番良いことに改めて気づきました．
   そこで情報I の担当になったとして適当な回数の講義案を\LaTeX で準備してください．
   その講義案を提出して貰えればそれをレポートとして受け付けます．
\end{exQ}


%\begin{flushright}
%   \qrset{link, height=1.5cm}
%  \qrcode{https://snow00two.github.io/mySite/}
%\end{flushright}

%
%\begin{thebibliography}{99}
%   \bibitem{Jobs}
%      アイザックソン, スティーブ・ジョブスI, II, 講談社, 2011.
%      
%   \bibitem{interface}
%      Apple Computer, 
%      Apple Human Interface Guidelines: 
%      The Apple Desktop Interface, Addison-Wesley, 1987,
%      日本語版, 1989.
%
%   \bibitem{F}
%      Feynman, Lectures on computation, Perseus, 1996, 
%      日本語版：ファインマン計算機科学, 岩波書店, 1999.
%
%   \bibitem{Fu}
%      藤原進之介, 
%      学校で習っていなくても読んで理解できる藤原進之介のゼロから始める情報 I, 
%      KADOKAWA, 2023. 
%  \bibitem{history}
%      ゴールドスタイン, 復刊計算機の歴史, 共立, 2016.
%
%   \bibitem{K}  
%      カーニハン, 
%      教養としてのコンピューターサイエンス講義 第二版, 日経BP, 2020.
%  
%     \bibitem{UNIX}
%      カーニハン,
%      カーニハンのUNIX回顧録, 丸善, 2025.
% 
%    \bibitem{Kay}
%      アラン・ケイ, アラン・ケイ, ASCII, 1992.
%      
%    \bibitem{W}クジラ飛行机, 
%         プログラミング言語大全, 技術評論社, 2025.
%
%   \bibitem{Orwell}
%      H.G.~Orwell, Nineteen Eighty-Four, 1949.
%
%    \bibitem{OM}
%       奥村 晴彦, 森本 尚之,
%       ［改訂第5版］基礎からわかる情報リテラシー, 技術評論社, 2023. 
%       
%   \bibitem{tools}
%      ラインゴールド, 思考のための道具, パーソナルメディア, 1987.
%         
%   \bibitem{Sa}
%      坂村健, 高等学校 情報 I , 数研出版, 2021.   
%   
%   \bibitem{T}
%      山口泰, 東京大学教養学部テキスト情報 第三版, 
%      東京大学出版会, 2025.
%   
%      
%    \bibitem{Turing}
%       A.~Turing, On Computable Numbers, with an Application to the Entscheidungsproblem, Proceedings of the London Mathematical Society 42, 1936.
%       
%  \bibitem{IT}
%      渡辺さき［令和8年度］ITパスポート 超効率の教科書＋よく出る問題集, インプレス, 2026.
%
%  
%\end{thebibliography}
%
%
%\end{document}


%\documentclass[a4paper,12pt]{amsart}
%\usepackage[dvipdfmx]{graphicx,xcolor}%graphics
%\usepackage{float}
%\usepackage{amsmath}
%\usepackage{verbatim}
%\usepackage{latexsym}
%\usepackage{amssymb}
%%\usepackage{emath}
%\usepackage{listings}
%\usepackage{ccicons}
%\usepackage{mflogo}
%\usepackage[utf8]{inputenc}
%\usepackage[T1]{fontenc}
%\usepackage{otf} 
%%\usepackage{fancybox}%enclose by box
%\usepackage[dvipdfmx]{hyperref}
%\usepackage{verbatim}
%\hypersetup{% hyperrefオプションリスト
%   setpagesize=false,
%   bookmarksnumbered=true,%
%   bookmarksopen=true,%
%   colorlinks=true,%
%   linkcolor=red,
%   citecolor=blue
%} 
%\renewcommand{\figurename}{図}
%\renewcommand{\tablename}{表}
%\renewcommand{\proofname}{証明}
%
%\newtheorem*{exQ}{レポート}
%\theoremstyle{remark}
%\newtheorem*{Q}{演習}
%\def\emptyline{\vspace{12pt}}
%\def\smallemptyline{\vspace{6pt}}
%\begin{document}
%\begin{flushright}
%   2026/04/28
%   %2026/04/21
%   %2026/04/03
%   %2026/03/11
%   %2026/02/24
%   %2025/12/22
%   
%   九大・数理 大津幸男
%\end{flushright}
%
%\begin{center}
%   {\Large 情報科指導法II---02}
%   \href{https://creativecommons.org/licenses/by-nc-nd/4.0/}{\ccbyncnd}
%   \par{デジタルデータ 1}
%\end{center}

\chapter{デジタルデータ}
\section{二進法}
\begin{itemize}
   \item
   辞書を引くと「digit とは指, あるいはアラビア数字の0〜9 のこと」とあり，
   その形容詞，あるいは(鍵盤楽器の)キー が digital で 対立語が analog とある．
   このことからdigitalは$10$進法と関係しているように見える．しかし，
   
   \item
   $0$, $1$を中心に置いた
   \term[]{$2$進法}(binary system)がデジタルの最小単位である． 
   ブール代数では$x^2=x$をみたす$x$が大事で，
   どの体でもこのような数は$0$, $1$しかないことから$2$進法が表れる
   という説明を受けたことがある\footnote{Bool\cite{Bo}のコメントによる}．
%  \href{https://www.youtube.com/watch?v=He_usBp673w}
   {電源ボタン}も
   $0$と$1$を並べてデザインされたもの（図~\ref{power}）
   である\footnote{$\, |$は電線が繋がった様子を模式化したものという説明もある．}．
   \begin{figure}[htb]
      \centering
      \includegraphics[width=7cm, trim=0 500 1000 0, clip]{~/Documents/GitHub/introInfo/ch2/images/ScannedDocuments.pdf}
      \caption{私の研究室のプリンタの電源ボタン}\label{power}
   \end{figure}
      
   $2$進法では
    \[
       1011=2^3+0\cdot 2^2+2^1+2^0=11
    \]
      のよう数字が表される．
   $0$, $1$を\term[]{ビット}{bit}(小片,  または，binary digit の略語) という．
   $N$ビットで表せる数は$2^N$個である． 
   代数的には剰余環$\mathbf Z/2^N\mathbf Z$と見なすことが自然だが，
   これについては暗号の回に解説する．
   
   また，負の整数(補数)や小数も$2$進法で表され，
   コンピューター上では特殊な方法で処理される．
   入試などではよく出されるようだが今回は省略する．

   \item
   \term[]{$16$進法}{hexadecimal system}では$0$〜$15$を
    \[
       0,  1,  2,  3,  4,  5,  6,   7,  8,  9,  \text{A},    \text{B},    \text{C},    \text{D},   \text{E},    \text{F}
    \]
   で表す．
   $16$進法を使っていることを明記するため 0x1D (0x1d のように小文字を使う場合もある)のように書くことがある．
   
   \item
   $8$ビットを$1$\term[]{バイト}{Byte}という．$2^8=256=16\times 16$なので
   $16$進法の二つの数で表される．
   例えば
   \[ 
     \text{0x3A}= 3\times 16 + 10 =58
   \]
   のようになる．ただし，歴史的にはByteが$8$ビットに決まるまでは紆余曲折があったそうである：
   \href{https://www.itmedia.co.jp/news/articles/2202/03/news151.html}
   {なぜByteは$8$ビットに決まったのか？}
    
 \end{itemize}
 
 上のようなコンピューターの特性からコンピューター上でデータを表すにはそのデータを二進法で表す（近似する）必要がある．
 これを\term[]{符号化}という．以下データ毎にどのように符号化されているかを解説したい． 大きく二つに分けて，
文字のならびからなるデータを\term[]{テキスト}{text}といい，
その他のデータをまとめて\term[]{バイナリー}{binary}という．
もちろん，
数自体も大事なデータであるが，ここではテキストの一部と見なす．
 
\section{テキストデータ}
 \begin{itemize}
   \item
   %Colume table
   \begin{table}[h]
      \centering
      \begin{tabular}{l|llllllll} 
      \noalign{\hrule height 1pt}
      &0&1&2&3&4&5&6&7\\ \hline \hline
      0&\textcolor{green}{NUL}&DLE&SP&0&@&P&`& p\\ \hline
      1&SOH&DC1& ! &1&A&Q&a& q\\ \hline
      2&STX&DC2& '' &2&B&R& b& r\\ \hline
      3&ETX&DC3&\# &3&C&S& c& s\\ \hline
      4&EOT&DC4&\$&4&D&T& d& t\\ \hline
      5&ENQ&NAK& \%&5&E&U& e& u\\ \hline
      6&ACK&SYN&\&  &6&F&V& f& v\\ \hline
      7&\textcolor{green}{BEL}&ETB&' &7&G&W& g& w\\ \hline
      8&\textcolor{green}{BS}&CAN&(&8&H&X& h& x\\ \hline
      9&\textcolor{green}{HT}&EM&)&9&I&Y& i& y\\ \hline
      A&\textcolor{green}{LF}&SUB&*&:&J&Z& j& z\\ \hline
      B&VT&\textcolor{green}{ESC}&+&;&K&[& k&\{\\ \hline
      C&FF&FS&, &$<$&L&\textbackslash & l& $|$ \\ \hline
      D&\textcolor{green}{CR}&GS&- &=&M& ] & m& \} \\ \hline
      E&SO&RS&. &$>$&N&\textasciicircum & n& 〜 \\ \hline
      F&SI&US&/&? &O & \_ & o & \textcolor{green}{DEL}\\ \hline
        \noalign{\hrule height 1pt}
      \end{tabular}
      \smallemptyline
      \caption{ASCII---上の行の数字＋列の$16$進数}
   \end{table}
%Width table
%\begin{table}[h]
%  \begin{tabular}{l|llllllllllllllll} 
%  \noalign{\hrule height 1pt}
%  &0&1&2&3&4&5&6&7&8&9&A&B&C&D&E&F\\ \hline \hline
%  0&NUL&SOH&STX&ETX&EOT&ENQ&ACK&BEL&BS&HT&LF&VT&FF&CR&SO&SI
%  \\ \hline
%  1&DLE&DC1&DC2&DC3&DC4&NAK&SYN&ETB&CAN&EM&SUB&ESC&FS&GS&
%  RS&US\\ \hline
%  2&SPA& ! & '' &\# &\$ &\% & \& &' &(&)&*&+&, &- &. &/\\ \hline
%  3&0&1&2&3&4&5&6&7&8&9&:&;&$<$&=&$>$&? \\ \hline
%  4&@&A&B&C&D&E&F&G&H&I&J&K&L&M&N&O \\ \hline
%  5&P&Q&R&S&T&U&V&W&X&Y&Z&[&\textbackslash & ] &\textasciicircum & \_ \\ \hline
%  6&`&a& b& c& d& e& f& g& h& i& j& k& l& m& n& o \\ \hline
%  7& p& q& r& s& t& u& v& w& x& y& z&\{& $|$ & \} & \textasciitilde & DEL\\
%    \noalign{\hrule height 1pt}
%  \end{tabular}
%\end{table}
ASCII (American national Standard-Code for Information Interchange) 
   米国情報交換用標準コード, 33個の制御文字とスペースを含むキーボードの
   特殊文字・数字・アルファベットの大文字小文字の文字を0から127
   までの$2^7$個の整数で表したもの．1963年に標準化された．
   例えば 0x7A は z を表す．
   二進数では$01111010$となり一番左の桁は常に$0$に成っていることも後で重要になる．
   また，制御文字の中で黒のものは現在では使われていない．
   
  
   ASCII の文字が写真のキーボーボードにすべて表れている
   (ただし，Tab, Control, Shift, Opt(Alt), Cmd, Fn などの修飾キーと Return は除く)
   ことに注意．
   
   \begin{figure}[htb]
      \centering
      \includegraphics[width=12cm]{~/Documents/GitHub/introInfo/ch2/images/keyboard.pdf}
      \caption{keyboard(HHKB)--QWERTY配列}\label{HHKB}
   \end{figure}

%  \begin{figure}[htb]
%     \centering
%     \includegraphics[width=12cm]{~/Documents/GitHub/introInfo/ch2/images/keyboardSet.pdf}
%     \caption{keyboard--QWERTY配列}\label{HHKB}
%  \end{figure}
   
   \TeX でお馴染みかも知れないが\verb/ \/ とASCII に含まれない ￥
   が文字化けする問題がある．
   Mac でもこれが困った問題を引き起こすことがあるので注意が必要．

   ここに載せたキーボードは
   左上のアルファベットの並びがQWERTY であることから
   QWERTY 配列と呼ばれている．

   tab （インデント）ついて：tab には機械タブ\verb/   /
   (これではよく分からないので\verb/\t/と書くことがある )と
   空白何個かで作るタブ\verb/   /の二種類がある．
   プログラミンではコードのツリー構造を人間に分かり易くし可読性を高めるため
   タブを入れることが勧められるが，
   機械タブと空白何個かで作るタブを区別するもの（コンパイラー）があるので
   注意が必要．美的にも二つのタブが混ざると綺麗でない．
   
   \emptyline 
     キーボードが出てきたのでついでに：     
     一生ものの外付けキーボード（特殊文字の位置固定）を手に入れて（図~\ref{HHKB}のキーボードだとテンキーは別途購入する必要がある）：
   \begin{enumerate}
      \item
         touch-type\footnote{以前はブラインドタッチと呼ばれていたが差別的な響きがあったことからこう呼ばれることが増えた．}をマスターし，
         コマンドラインインターフェースに慣れる．
         \begin{enumerate}
            \item
               F と J のキーの上には突起があり，
               人差し指の触覚だけで位置が確認できるようになっている．
            \item
               ラップトップパソコンだとついキーが見えてしまうので，
               タッチタイピングの練習のためにも
               外付けキーボードが必要である(私はやったことが無いが，
               見えないように手の上に布を置くという方法もある)．
            \item
               ラップトップのキーが外付けキーボードと異なるときは，
               調整する機能拡張(私はKarabineを使っている)が必要になる．
         \end{enumerate}

      \item
         よく使うアプリのよく使うメニューアイテムのショートカットを覚え，
         スピードを落とすマウスの使用量を減らす．
      
      \item
         キーキャップリムーバーを手に入れ，キーボードの洗浄を定期的に行う．     
   \begin{figure}[htb]
      \centering
      \includegraphics[width=8cm]{~/Documents/GitHub/introInfo/ch2/images/keyremover.pdf}
      \caption{Keycap Remover}\label{}
   \end{figure}

   \end{enumerate}
   ことをお奨めする．

\begin{Q}
   \, 
   
\begin{verbatim}
   % echo ABCDEFGHIJK > test     
   % od -c -t x1 test
   0000000    A   B   C   D   E   F   G   H   I   J   K  \n    
     41  42  43  44  45  46  47  48  49  4a  4b  0a    
   0000014
\end{verbatim}
   
   最初の行の最初の\verb/% /はprompt,  次の\verb/ >/ は redirect, 
   \verb/ >>/ にするとファイルがあるときはそれに追加される．
   \verb/od/ は octal dump の略で\verb/ - /以下はoption 
   
   \emptyline
   経験者からのアドバイス：
   演習で解答し（他でもコンピューターの知識，あるいは，コンピューターに限らず何でも調べ）たら，
   メモ帳などに説明を含む簡単なメモを残すようにして下さい．
   結局は時間の節約になります．
   \begin{figure}[h]
      \centering
      \includegraphics[width=12cm, trim=10 50 10 0, clip]{~/Documents/GitHub/introInfo/ch2/images/memo.pdf}
      \caption{私のMacのメモ(Notes)}\label{memo}
   \end{figure}
\end{Q}

\begin{exQ}
   タッチタイピングを練習し，出来るようになったら
   その様子が確認出来るような動画をつくり，
   私が確認出来るようにそのURLを送って下さい．
\end{exQ}



   \item
   日本語の文字コードについてはひらがな・カタカナ・漢字などの文字の種類が多いことから
   2Byte($=2^{16}\approx  64,000$) の数字を使う二バイト文字を使う必要があった．
   しかもその方法にJIS, SJIS, EUC など異なったものがあり文字化けが問題であった．
   更に，ネットの発達などにより世界中の文字への対応が必要となり， 
   \href{https://home.unicode.org/}{Unicode}
   \footnote{このサイトを訪れるとその壮大さに圧倒される．}が制定された．
   \begin{itemize}
      \item
      Unicode は数字に過ぎない．最近では約11万字以上を表すことができる．
      U+ の後に16進数で表し，U+000000〜U+10FFFFだけ有るが冗長性が大きいので，
      表現の体系が幾つかある．最近のスタンダードはUTF-8である，
      他にもUTF-16, UTF-32などがある．ただし，U+は入力しなくていい．
      
      \item
      例えばU+0041 はAを表し，
      ASCII を含んだ拡張になっていることが分かる．
      更にUTF-8では
      例えばU+0041を0x41の1byte で表すような，
      1byteの最初の2進数が0の文字をASCIIと判定するアルゴリズムを
      用いている．つまり，
      ASCII はASCII のままで使用できるようになっており
      既存のコンピューター資源をそのまま使うことが出来る．
      これがUTF-8が普及した理由である\cite{Y}．
       
      \item
      改行の示し方もUNIX, Mac ではLF (line feed character, line finder 改行文字)で，
      Windows ではCR+LFである．
      そこで改行の違うUTF-8が存在するので注意が必要である．
   \end{itemize}

   \begin{Q}
   CLIで文字化け対応 
\begin{verbatim}
   % brew install nkf  
   # nkf がインストールされていないとき
   # WSLでも多分 % sudo apt install nkf でOK
   % nkf -s letter.tex > test.tex
   % nkf -w test.tex > test2.tex
\end{verbatim}   %apt=Advanced Package Tool
   これらを普通のエディターで文字化けしないように開いて下さい．
   \end{Q}
   \begin{Q}
   Unicode のローカル（自分のPC）での入力法について検索して（あるいは生成AIに聞いて），
   ファイルに直接UTF-8で文字入力出来るようになってください．
   
   例えば事務から送られてくるメールでよく文字化けしている
   %\UTF{2460}と\UTF{3230}%
   %\ajMaru{1}と\ajKakkoYobi{1}
   のUnicode を調べて
   文字化けなく出力出来るようになって下さい．
   \end{Q}
   \begin{Q}
      Windowsなどで未だにSJISを使っている人は早めにUTF-8に移行することを
      お奨めします．やり方を調べて実行してください．
   \end{Q}


%  \item
%  日・時間データ
   
   \item
   文字のならびからなる文を\term[]{文字列}といい，
   プログラミングの中では変数や数や関数と区別するため  
   \verb/'/ や   \verb/"/ で囲んで表すことが多い．
   例えば \verb/"This is a statement."/ のように表す．
   \begin{Q}
   文字の数字の真の数字への変換：例えば shell で1+2を計算すると
   
   \begin{verbatim}
   % 1+2
   zsh: command not found: 1+2 #当たり前か？
   % echo 1+2 >text.txt 
   \end{verbatim}
   \begin{verbatim}

   % od -c -t x1 text.txt 
   0000000    1   +   2  \n                                                
            31  2b  32  0a   #少し考えると1+2 は文字列
   0000004
   % echo $((1+2)) #調べるとこうすれば良いことが分かる
   3
   % python #もっとも今だとこのあたりが普通か？
   >>> 1+2
   3
   >>> # control Dで終了
\end{verbatim}
   自分の使っている(興味のある)言語で文字の数字の数字への変換，
   逆の変換のやり方を調べてください．
\end{Q}

\end{itemize}
\section{画像データ}
 バイナリーデータは画像データや音声データ，プログラムの実行ファイルなどの区別がある．まずは画像データ：

\begin{itemize}     
   \item
   \begin{figure}[htb]
      \centering
      \begin{minipage}{0.4\columnwidth}
      \centering
      \includegraphics{~/Documents/GitHub/introInfo/ch2/images/figure.1}
     \caption{加法混色}\label{addColor}
      \end{minipage} 
      \begin{minipage}{0.4\columnwidth}
      \centering
      \includegraphics{~/Documents/GitHub/introInfo/ch2/images/figure.2}
     \caption{減法混色}\label{subColor}
      \end{minipage} 
   \end{figure}
   
   まず画像の物理的特性．
   テレビやパソコンのディスプレーで画像を作る単位をpixel(画素)という(図~\ref{Liquid Crystal}\footnote{
      \href{https://commons.wikimedia.org/w/index.php?curid=710606}
         {WikiPediaより By Ravedavel}
         \href{https://creativecommons.org/licenses/by/1.0/?ref=openverse}{\ccbysa}
   })．
   例えば私のモニターが$1200\times 1920$pixelの表示のときは
   $2,304,000\approx 2.3\cdot 10^6$個のpixelを使っていることになる．
   各pixelの色は \term[]{RGB}, つまり，Red Green Blueの略で
   加法混色(図~\ref{addColor}) することで
   色々な色を表現できることが分かっている．
   これらは$[0,1]$の中の実数の三つ組みで表せるが，実際は
   離散化（量子化とも言う）し，
   それぞれの色を1バイト($2^8=256$)に分けて近似する．
   全体で3バイトで
   \[
      2^8\cdot 2^8\cdot 2^8=2^{24}\approx 16\cdot 10^6
   \]
   種類の色を表すことができる（この仕様をフルカラーという）．それらがpixelの数だけあるので
   コンピュータの一画面を表すのに必要なデータは
   \[
      16\cdot 10^6\times2.3\cdot 10^6
      \approx 36.8\cdot 10^{12}
      \approx 40 \text{兆個}
   \]
   あることになる．
      \begin{figure}[htb]
      \centering   
      \begin{minipage}{0.5\columnwidth}
         \centering
         \includegraphics[width=4cm]{~/Documents/GitHub/introInfo/ch2/images/1024px-Liquid_Crystal_Display_Macro_Example_zoom.jpg}
         \caption{ pixel
         }
         \label{Liquid Crystal}
      \end{minipage} 
      \begin{minipage}{0.4\columnwidth}
         \centering
         \includegraphics{~/Documents/GitHub/introInfo/ch2/images/figure.8}
        \caption{色相環}\label{cycColor}
      \end{minipage} 
   \end{figure}

   色を表すのには直接 $0$〜$255$ を三つ並べてRGBを
   \[ 
      [125,45,255]\quad\text{or}\quad
      (125,45,255)\quad\text{or}\quad 
      \text{color}(125,45,255)
   \]
   のように書いたり，
   \term[]{カラーコード}といって\verb/#/の後に16進法の数字
   を6つ(省略形で3つ)並べて
   \[
      \verb/#/\text{a43aff}
   \]
   のように表すことがある．また，red, blue, mazentaなどように
   既に定義された名前をつかえることがある
   （どのように表すかは使っている言語による）．
   \begin{Q}
      VSCode ではカラーコードを入力するとサンプルが表示される機能がある(図~\ref{colorCode})ので実際に確かめてください．
      \begin{figure}[h]
      \centering
      \includegraphics[width=12cm]{~/Documents/GitHub/introInfo/ch2/images/color.jpg}
      \caption{VSCodeでカラーコード表示}\label{colorCode}
   \end{figure}
   \end{Q}

   
   印刷用にはCMYK(Cyan Magenta Yellow blacK)という減法混色が用いられる(図~\ref{subColor})．
   
    \item        
      色については美術や心理学の観点から，
      \begin{itemize}
         \item 
         色相(hue)：赤や青などの色の違い(色相環 図~\ref{cycColor} 参)．
         \item 
          彩度(saturation)：色の鮮やかさ
         \item
         明度(luminosity, brightness)：色の明るさ
      \end{itemize}
      により説明される．更に
      \begin{itemize}
         \item 
            補色（図~\ref{cycColor}の反対側にある色）
         \item
            暖色・寒色
         \item
            有彩色・無彩色
         \item
            トーン
         \item
            ベースカラー，メインカラー，アクセントカラーの比率
      \end{itemize}
      などの概念を知って，
      それを使いこなす練習をすれば
      実際のデザイン・プログラミングにも役立つだろう
      (例えば \cite{R}, \cite{TK}, \cite{Uti}, \cite{W},  参照)．

   \item
      画像を動かすことで動画が得られる．個々の画像をframeという．
      テレビ放送では30fps (=frames per second)つまり，
      1秒間に30 frame 表示することになっている．
      したがって1秒間に$1200$兆個のデータを使っていることになる．
   
   \item
      画像データをコンピュータで表すには
      点と直線や曲線の組で表すベクタ表現，
      図形を画素の集合として表すラスタ表現
      の二種類の方法がある．
     ベクタ表現は拡大縮小が自由であるが
      ラスタ表現は拡大すると見づらくなる．
      \begin{figure}[htb]
         \centering
         \begin{minipage}{0.4\columnwidth}
         \centering
         \includegraphics{~/Documents/GitHub/introInfo/ch2/images/figure.3}
         \caption{ベクタ表現}\label{vector}
         \end{minipage} 
         \begin{minipage}{0.4\columnwidth}
         \centering
         \includegraphics{~/Documents/GitHub/introInfo/ch2/images/figure.4}
         \caption{ラスタ表現}\label{rast}
         \end{minipage} 
      \end{figure}
      
   \item
      文字符号を具体的に画像データとして表示した物が
      フォントである．
      フォントにはベクタ表示のアウトラインフォントと
      ラスタ表示のビットマップフォントがある．
      近年ではアウトラインフォントが主流となっている．
      現代では様々な書体のフォントが使用可能である．
      
      ただし，□（豆腐）という形で，
      フォントがない文字が表されることがある
      （いたずら好きの友達がヒエログリフで送ってきた可能性もないわけではないが，
     pdfで文字の埋め込みに失敗した可能性の方が高い）．
      これを起こさないためにGoogle による
      NoTo (=No Tofu)いうもプロジェクトが進んでいる．

\end{itemize}

%\begin{thebibliography}{9}
%   \bibitem{B}
%      Bool, {\em Laws of thought}, 1848
%   
%
%    \bibitem{R}   
%       ガー・レイノルズ, プレゼンテーションZen 第3版, 丸善出版, 2021.
%    \bibitem{TK}
%    高橋 佑磨, 片山 なつ, 
%    伝わるデザインの基本 増補改訂3版 
%    よい資料を作るためのレイアウトのルール, 技術評論社, 2021.
%
%    \bibitem{U}
%       内田 広由紀, 7日間でマスターする 配色基礎講座 
%       (デザイン ビギナー シリーズ), 2000
%    \bibitem{W}
%       Robin Williams,  ノンデザイナーズ・デザインブック [第4版], 
%       マイナビ出版, 2016
%   \bibitem{Y}
%      矢野 啓介, 
%      [改訂新版］プログラマのための文字コード技術入門, 技術評論社, 2018.    
%\end{thebibliography}
%
%\end{document}


\section{音声データ}
   音は物理的には空気中の圧力の差が秒速$340$mぐらいで伝播していく縦波であり，
   それを静止あるいは比較的低速で移動する動物が耳で捉えたものである．
   よって，
   音楽データは時間$t$についての圧力の変化$f(t)$によって表せる．
   まず$f(t)$は単なる連続関数で(場合によっては不連続点のあるものも扱うので)
   そのままではコンピュータでは扱えないことを注意する．
   
   ヒトの聴覚は 20Hz(ヘルツ)から2万Hzまでといわれているので
   ある程度高い音は聞こえないと考え，
   フィルターで除外する．
   その周波数を$W$とし，そのため生じる最小周期を$T=\frac{2\pi}{W}$と置く．
   得られた関数をもう一度$f(t)$と置き，
%   つまり，
%   この$f(t)$のフーリエ変換$\hat f(\omega)$と置くとき
%$\omega\ge W$で$\hat f(\omega)=0$と仮定する．すると$\hat f(\omega)$は
%$[-W,W]$上の関数となり，$\mathbf R$上に周期的に拡張することで
%$S^1$上の関数と見なせる．%      のでその逆変換で
%      $f(t)$が復元できるが，それはFourier級数と解釈出来る．
$n$を整数として 
     \[
        a_n=f(n T)
     \]
   を作る
   (図~\ref{sampling})．$W$を\term[]{サンプリング周波数}，
   $\{a_n\}$を作ることを \term[]{サンプリング}という．
   \[
      f_n(t)=\frac{\sin\pi(n-\frac t{T})}{\pi(n-\frac t{T})}
   \]
   を使って
   \[
      \bar f(t) := \sum_{n=-\infty}^{\infty}a_nf_n(t)
   \]
   を作ると，$f$が$\bar f$と一致するというのが
   シャノン-染矢の\term[]{サンプリング定理}である．
       
   現実には$a_n$の値は二進数に量子化され，
      そのビットを\term[]{量子化ビット数}という．
   音は高周波数成分を含むことと量子化により
   ある程度の誤差は避けられない．
    
   \begin{figure}[htb]
   \centering
   \begin{minipage}{0.4\columnwidth}
      \centering
      \includegraphics{~/Documents/GitHub/introInfo/ch2/images/figure.5}
      \caption{音波と標本化}\label{sampling}
   \end{minipage} 
   \begin{minipage}{0.4\columnwidth}
      \centering
      \includegraphics{~/Documents/GitHub/introInfo/ch2/images/figure.6}
      \caption{$f_0$(t)}\label{}
   \end{minipage} 
   \end{figure}

%   \emptyline
   
   具体的にCDの場合の規格を見てみると，
   サンプリング周波数は$44.1$kHzで
   量子化ビット数は$16$bit である．
   したがって，$1$秒間で必要なデータ量は
   \[
      64000 \times 44100\approx 2.8 \times 10^{11}
   \]
   ステレオだとこの$2$倍である．
   \emptyline
   
   DTM(=Desk Top Music)のためには音楽の理論や心理学についても
   必要になると思うがここでは取り上げない．
   このセクションの名前は音声データであるが，
   人の声は意味を伝える特別な音であるので後にもう一度取り上げる．
  
\section{実行ファイル}
コンピュータのバイナリーの中で重要なものの一つはプログラム(の実行ファイル)そのもののである．C言語のファイル
\begin{verbatim}
   #include <stdio.h>
   
   int main(){
   	printf("Hello, World!\n");
   }
\end{verbatim}
をhello.c という名前で保存．
冒頭の\verb/<stdio.h>/は標準入出力を使うために必要なライブラリーであり，それを最初に読み込んでいる．
そのディレクトリで
\begin{verbatim}
   % ls
   hello.c         README.md
   % gcc hello.c 
   % ls
   a.out           hello.c         README.md
   % ./a.out
   Hello, World!
   % od -c -t x1 a.out | more
\end{verbatim}
最後の more ではスペースキーでページが更新され，Bで前のページに移り，Qで終了する．
gcc は gnu\footnote{GNUはGUN's Not Unix の頭文字でグニューと読む．R.~Stallman が始めた．} c compilerの略である．

\begin{Q}
上の作業を自分で行い，Hello, World! が何処にあるかを探してください．
\end{Q}
現在では(システム開発者以外は)C言語を知らなくてもそれほど困らないと思うが，
参考文献としてはやはり\cite{KR}を挙げておく．
また，これからは Rust や Go などを学んだ方が役に立つかもしれない\footnote{プログラミング言語の将来予測は極めて難しいことが知られている．}．
\emptyline 

Shell Script について：
\begin{verbatim}
   #!/bin/sh
   
   echo "Hello World!"
\end{verbatim}
を hello という名前で保存．
\verb/#!//bin/sh はShell Script の書き出しの決まり文句である\footnote{私が昔読んだ\cite{KP}ではここが必要とは書れていないが，\cite{Ya} では書くのは当然と書いてある．}．
そのディレクトリで
\begin{verbatim}
   %ls -l
   total 8
   -rw-r--r--  1 snow00two  staff  33 Apr 30 13:28 hello
   # このままだと実行可能でない
   % chmod +x hello   # 実行可能に変換
   % ./hello
   Hello World!
\end{verbatim}
helloはテキストファイルであるが，実行ファイルでもあるので，これもバイナリーと呼ばれることがある．

\begin{Q}
上の作業を自分で行ってください．
\end{Q}


\section{データの取扱い}
次にデータの取り扱いについて述べる：
 \begin{itemize}
   \item
      データの授受の時よく使われるテキストの形式のことを
      \term[]{構造化データ}という．
      いずれも基本ASCIIファイルであり，
      (文明が崩壊しない限り)未来永劫どのようなエディタでも開き
      編集することが出来るものである．
      \begin{enumerate}
         \item
            CSV =Comma Separated Values： 
            Excel のデータを文字列で書き出し
            $,$(Comma)で区切ったテキスト (図~\ref{csv})．
            $,$をtabに代えた(TSV =Tab Separated Values) もあるが，
            .csv の拡張子なのに実はtabで区切られたものも多い．
            
            また，Windows のcsvファイルにはファイルの先頭には，
            BOM (=Byte Order Mark, ぼむ)という
            3バイトのデータが着いていることがあり，
            Mac・UNIX環境では悪さをすることがある．
            BOMの取り除き方については検索して探してほしい．

         \begin{figure}[htb]
         \centering
         \includegraphics[width=12cm]{~/Documents/GitHub/introInfo/ch2/images/ScreenCapture1.pdf}
         \caption{CSV:\href{https://www.nstac.go.jp/use/literacy/ssdse/past/}{SSDSE-F-2023}}\label{csv}
         \end{figure}
         %TSV= Tab Separated Values は$,$の代わりにTab( HT=0x09)で区切ったもの．
              \begin{figure}[htb]
            \centering
            \begin{minipage}{0.45\columnwidth}
               \centering
               \includegraphics[width=6cm]{~/Documents/GitHub/introInfo/ch2/images/ScreenCapture2.pdf}
               \caption{JSON}\label{JSON}
            \end{minipage} 
            \begin{minipage}{0.45\columnwidth}
               \centering
               \includegraphics[width=6cm]{~/Documents/GitHub/introInfo/ch2/images/ScreenCapture3.pdf}
               \caption{YAML}\label{YAML}
            \end{minipage} 
         \end{figure}

      \item
         JSON = JavaScript Object Notation：元々はJavaScriptの連想配列から
         作ったもの\footnote{JSONは2001年クロックフォードによりJavaScript
         の中から「発見」されたといわれている．}だが最近色々な所で使われる(図~\ref{JSON})．
         npm では package.json という
         ソフトウェアの依存性を記述したファイルがよく使われている．
         
      \item   
         YAML = YAML Ain’t Markup Language：Pythonに似た構造化データ．
         JSONよりコンパクトに記述できる(図~\ref{YAML})．
         ソフトの設定などに使われる．
     \item 
        XML  = eXtendable Markup Language： 
        HTML に似たデータ記述言語で一時期はやったが，
        最近はあまり使われない．しかし，
        パソコンの内部では今でもよく使われている．
              
   \end{enumerate}
   \emptyline

これらのアルゴリズムはコマンドラインのコマンドの形や，
   プログラミング言語のライブラリーの形で提供されている．
   
   \begin{Q}
      yaml, json, xml 間の変換：yq\, を使う
      
\begin{verbatim}
   % cat math.json 
[
   {
      "lastname": "円", 
      "firstname": "周率"
   },
   {
      "lastname": "黄金", 
      "firstname": "率"
   },
   {
      "lastname": "虚数", 
      "firstname": "単位"
   }     
]
   % brew install yq #yqがインストールされていないとき
   % yq -P -oy math.json > math.yml 
   % cat math.yml
- lastname: 円
  firstname: 周率
- lastname: 黄金
  firstname: 率
- lastname: 虚数
  firstname: 単位
   % yq -o json math.yml  > math2.json
\end{verbatim}


   \end{Q}

   \item
      Base64：メールやHttpのヘッダなどでバイナリーをテキストに変換して送る方式．
      バイナリーを$6$ビットずつに分け，それを0〜63 の数字で表すとき
      0から25にA〜Z, 26から51にa〜z, 52から61に0〜9, 
      62, 63に\verb/-/  \verb/_ / を当てはめることで
      （最後の二つは異なるやり方もある）
      文字列に変換するやりかたである．
      \emptyline
      
      このアルゴリズムもコマンドラインのコマンドの形や
      プログラミング言語のライブラリーの形で提供されている．
\begin{Q}\, 

\begin{verbatim}
   % echo "abc" > test
   % tar -czf test.tar.gz test #適当なバイナリーを作成
   % base64 -i test.tar.gz
   H4sIAIj+uGkAA+3PMQ6AIAwFUGZP0SO0WOA8aGQ1Ebi/IMaRDY1J3/
   LTtMNv2mJSgyGiZYaazporUbf5xkCGjHGacbaApImcAhxdrMox+aNU2
   VPMvbtyFkJn3x6BJ3/CL+v0dQchhBDvOwEZLljpAAgAAA==
\end{verbatim}
\verb/#/ 画面に入らないので適当に改行
\end{Q}

  
 
   \item
     任意の長さのデータから固定長のデータを取り出す関数を
     ハッシュ関数(hash funcition)といい，その値を
     ハッシュ値(hash value)という．
    特に\term[]{一方向ハッシュ関数}が重要である．
 
      ただこのままでは何を言っているのかさっぱり分からないと思うので，
      一方向ハッシュ関数に期待されてる性質をまとめておく．
     \begin{itemize}
         \item
            どのような長さのデータでも出力は同じ長さ
         \item
            計算は易しく，大して時間が掛からない．
         \item
            耐衝突性：異なる入力には異なる出力を出す．
            もっともディリクレの原理から出力のビット数を超える個数の
            入力の中では必ず衝突が起こるので，
            少ない個数の異なる入力には異なる出力を返すという意味である．
            
         \item
            一方向性：ハッシュ値から入力の逆算が難しい．
      \end{itemize}
      \emptyline

      ハッシュ値は暗号とは異なることに注意する．
      また，そのアルゴリズムの設計も難しそうであることも分かる．
      \emptyline
      
      
      ハッシュ値の使い方について：
      ハッシュ値はコンピュータ・セキュリティの分野で広く日常的に使われている．幾つか例を挙げるが，これから先も関係した話題が出てくるだろう：
       \begin{itemize}
         \item 
            パスワードの保存：パスワード自体をそのままで保存するのは危険なので，
            パスワードのハッシュ値を保存して，パスワードを打つ毎にそのハッシュ値を
            比較してパスワードが正しいことを毎回確認することが行われている．
            \emptyline

            ただし，ありそうなパスワードのハッシュ値を事前に計算していおいて，
            そのパスワードのハッシュ値のリスト(\term[]{レインボーテーブル}という)
            と比較してパスワードを推測する\term[]{辞書攻撃}があるので，
             ユーザー毎に疑似乱数(ソルトという)を決めておき，
            パスワードにソルトを加えたデータのハッシュ値を使う方法も行われている．
         
         \item 
            ソフトウエアを配布するとき，
            配布元のサイトのミラーサイトを使うことはよく行われているが，
            それが本物であることを保証するため，そのソフトウエアのハッシュ値を計算して
            配布元のサイトに置いておき，ミラーサイトからソフトウエア手に入れた入手者が
            そのハッシュ値を計算し，配布元のサイトのハッシュ値と比較して
            それが本物であることを確認することが行われている．
       \item
          疑似乱数：先程もでてきたが一方向ハッシュ関数をつかうことで
          (ランダムに近く見える)疑似乱数の生成が行われる．

      \end{itemize}
      \emptyline

  \begin{table}[h]
      \centering
      \begin{tabular}{llll} 
           \noalign{\hrule height 1pt}
             一方向ハッシュ関数& 長さ & 発表年& コメント \\ \hline \hline
            MD5& 128bit& 1992  &危殆化 \\ \hline
            SHA-1& 160bit&  1995&危殆化\\ \hline
            SHA-256&256bit &  2001&\\ \hline
            SHA-512& 512bit&  2001&\\ \hline
            SHA-3& 可変&   2015&\cite{Y}に詳しい説明\\ \hline
           \noalign{\hrule height 1pt}
      \end{tabular}
    \smallemptyline     
      \caption{使われている(た)一方向ハッシュ関数}
   \end{table}

       現在使われている(た)一方向ハッシュ関数を挙げておく．
       ただし，コメントに書かれた\term[]{危殆化}とは
       既に 不十分であることが分かっていることを意味していて，
       互換性のために未だサポートされているが，
       これからの積極的な使用は推薦されていないことを意味している．
       \emptyline

         これらのアルゴリズムもコマンドラインのコマンドの形や
         プログラミング言語のライブラリーの形で提供されている．

\begin{Q}
\,

\begin{verbatim}
   % echo "abc" > test
   % sha256 test
\end{verbatim}
{\small SHA256 (test) = edeaaff3f1774ad2888673770c6d64097e391bc362d7d6fb34982ddf0efd18cb}

\noindent
\verb/#/ 画面に入らないので小さくした．
\end{Q}

    \item
        圧縮：データを保存したり送ったりするとき
        内容を変えないように圧縮して送れると
        保存量・通信量を減らせて都合が良い．
        元のデータ量と圧縮したデータ量の比を\term[]{圧縮率}という．
        \[
           \text{圧縮率}
           =\frac{\text{圧縮した後のデータ量}}
              {\text{圧縮する前のデータ量}}\times 100 \, \, \verb/%/
        \]
        
        圧縮をといたとき元のデータが完全に復元できるのが\term[]{可逆的圧縮}，
        データが完全に復元できず欠損が起こるのを\term[]{不可逆圧縮}という．
        とくに音声・映像データは不可逆圧縮で十分(で遅れがない方が大事)なことが多い．
        いくつか圧縮法の例を挙げる：
       
       \begin{itemize}
         \item
            ランレングス法：
              この紙（あるいはモニター画面）のような
              白と黒などの繰り返しの多いデータの時，
              連続する同じデータの数を数えて示すのが
               \term[]{ランレングス法}である．
              例えば
              \[
                 AAAAABBAAAAAAAABBBBBBBAAAAA
              \]
              のようなデータを$A5B2A8B7A5$と圧縮する．
              
         \item
           ホフマン法：例えば
\begin{gather*}
\text{おきゃがれこぼし　おきゃがれこぼし}\\
\text{ゆんべもこぼして　またこぼした}
\end{gather*}
という文(歌舞伎\href{https://www.youtube.com/watch?v=Fif6tRWZZVQ}{「外郎売り」}より, \cite{Kabuki}には先代の市川團十郎の舞台の動画もある．)を（空白を除いて）適切なかたまりに分割すると次のような頻度になる．
         \begin{table}[h]
            \centering
            \begin{tabular}{l|llllll} 
                 \noalign{\hrule height 1pt}
                     部分& こぼし & おきゃがれ& ゆんべも& て& また&た \\ \hline 
                     頻度 & 4 & 2& 1& 1&1&1\\ \hline
                     コード & 00 & 01& 100& 101&110&111\\ \hline
                 \noalign{\hrule height 1pt}
            \end{tabular}
            \smallemptyline     
            \caption{ホフマン法の例}
         \end{table}
         
         \noindent
          $6$個の部分があるので$3$bit 必要だが，
          頻度が多い$2$つの部分に$0$から始まる短い$2$進数$00$, $01$を与え，
          残りの$4$つに$1$ではじまる$3$bitの$2$進数を表のように与えると
           \[   
              010001001010010111000111
           \]  
           と置き換えることが出来る．
           文字の区切りがちゃんと判定できることを確かめてほしい．
           このように頻度に応じて短い数を当てるのが
           \term[]{ホフマン法}である．


      \end{itemize}
       %JPEG, ZIP で使われている．
      データの特徴を上手く使った方法を選ぶ必要がある．
      \begin{exQ}
         発声練習：「外郎売り」を練習し，発声出来るようになったら
         その様子が確認出来るような動画をつくり，
         私が確認出来るようにそのURLを送って下さい．
      \end{exQ}

\end{itemize}


%\section{補足：Fourier 変換}

%\begin{thebibliography}{9}
%%   \bibitem{B}Bool, {\em Laws of thought}, 1848
%%    \bibitem{Ko} Koblitz, A course in Number theory and Cryptography, GTM,  \href{https://link.springer.com/search?package=11649&facet-content-type=%22Book%22}{ダウンロード}
%   \bibitem{K}
%      木村秀紀, Fourier-Laplace 解析, 岩波書店, 1993.
%   \bibitem{KP}
%      Kernighan, Pike, 
%      Unix programming environment, Pearson, 1984
%      (訳 UNIX プログラミング環境, KADOKAWA, 2017).
%   \bibitem{KR}
%      Kernighan, Ritchie, 
%      The C programming  Language--2nd edition, 
%      Prentice Hall,1989. 
%      (訳プログラミング言語C 第2版 ANSI規格準拠, 共立出版, 1989.)
%   \bibitem{Ya}
%   山森丈範，[改訂第4版]シェルスクリプト 基本レファレンス, 技術評論社, 2024.
%   \bibitem{Y}
%      結城浩, 暗号技術入門 第３版, SBクリエイティブ, 2015.
%   \bibitem{U}
%      小学館DVDBook, 歌舞伎のいき, 第2巻「時代物・荒事編」, 2008.
%
%\end{thebibliography}


%\documentclass[a4paper,12pt]{amsart}
%\usepackage[dvipdfmx]{graphicx,xcolor}%graphics
%\usepackage{float}
%\usepackage{amsmath}
%\usepackage{verbatim}
%\usepackage{latexsym}
%\usepackage{amssymb}
%\usepackage{tabularray}
%%\usepackage{emath}
%%\usepackage{listings}
%\usepackage{ccicons}
%\usepackage{mflogo}
%\usepackage{wrapfig}
%\usepackage[basic, box,gate,oldgate, ic, optics,physics]{circ}
%
%%\usepackage{fancybox}%enclose by box
%\usepackage[dvipdfmx]{hyperref}
%\usepackage{verbatim}
%\hypersetup{% hyperrefオプションリスト
%   setpagesize=false,
%   bookmarksnumbered=true,%
%   bookmarksopen=true,%
%   colorlinks=true,%
%   linkcolor=red,
%   citecolor=blue
%} 
%
%\newtheorem*{exQ}{レポート}
%\theoremstyle{remark}
%\newtheorem*{Q}{演習}
%
%\renewcommand{\figurename}{図}
%\renewcommand{\proofname}{証明}
%\renewcommand{\tablename}{表}
%
%\DeclareMathOperator{\noT}{NOT}
%\DeclareMathOperator{\anD}{AND}
%\DeclareMathOperator{\oR}{OR}
%\DeclareMathOperator{\nand}{NAND}
%\DeclareMathOperator{\nor}{NOR}
%\DeclareMathOperator{\xor}{XOR}
%\DeclareMathOperator{\idm}{ID}
%\def\emptyline{\vspace{12pt}}
%\def\smallemptyline{\vspace{6pt}}
%\setlength{\belowcaptionskip}{0mm}
%
%\begin{document}
%\begin{flushright}
%   2026/05/19
%   %2026/03/11
%   %2026/01/12
%   
%   九大・数理 大津幸男
%\end{flushright}
%
%\begin{center}
%   {\Large 情報科指導法II---04}
%   \href{https://creativecommons.org/licenses/by-nc-nd/4.0/}{\ccbyncnd}
%   \par{コンピュータの原理 1}
%\end{center}

\chapter{コンピュータの原理}
\term[]{アルゴリズム}{algorithm}の語源はアラビアの数学者フワーリズミであるといわれる．
彼の数学書「インド人たちの数」が12世紀にラテン語に翻訳されたとき
「アル・フワーリズミは言った」
から文が始まっていたが，その後名前であることは忘れられ
インド・アラビア式の計算の方法\footnote{
インド・アラビア式の計算の歴史的な位置づけ：
古代インドに起源を持つ$0$を含む数字と位取り記法が，
アラビア数字に変わりながら
イベリア半島を通じてヨーロッパに入ってきたのも10世紀頃であった．
また，紙の生産は古代中国で始まったが，その製造法も
12世紀頃にアラビアを通じてヨーロッパに入ってきた．
それらは簿記などに使われ商業・金融の発展に繋がった\cite{T}．
その後
活版印刷が発明され，科学や工学や文化の発達，宗教の対立と戦争の拡大を招いた.}をアルゴリズムと呼ぶようになった 
(\cite{Mu})．

現在ではアルゴリズムは算法という意味で使われるが，
1930年代に数理論理学において論理の限界と計算可能性と関連して論じられるようになった．
フォン・ノイマンはこれらの研究のただ中にいたが，
それらをもとに1951年に世界初のノイマン型コンピュータ EDVAC を開発を成し遂げた．

ここでは論理学と電子回路から分かるコンピュータの仕組み，
更に時計付き回路とメモリーについて解説する．参考文献として\cite{F}, \cite{Mn}をあげておく. 

\section{論理回路}
\begin{itemize}
\item
$1$を真$0$を偽として表す．
コンピュータの内部では
ある一定の正の電圧$V$がかかった状態が$1$
電圧がかかっていない状態が$0$である．

\item
$A$, $B$ を$0$, $1$を取る変数とするとき, 
$A$の否定を  $\noT A$，
$A$かつ$B$を $A\text{ AND }B$,
$A$あるいは$B$を $A\text{ OR }B$と表す．
真理表で表すと
\begin{table}[htb]
   \centering
   \begin{minipage}{0.4\columnwidth}
      \centering
      \begin{tabular}{c|c}\noalign{\hrule height 1pt} 
        $ A$ &$\noT A$\\  \hline
         0&1\\
         1&0\\ \noalign{\hrule height 1pt} 
      \end{tabular}  
      \smallemptyline    
      \caption{$\noT$}   
   \end{minipage} 
   \begin{minipage}{0.4\columnwidth}
      \centering
      \begin{tabular}{cc|c}\noalign{\hrule height 1pt} 
         $A$ &$B$&$ A \anD B$\\  \hline
         0&0& 0\\
         0&1& 0\\
         1&0& 0\\
         1&1&1\\ \noalign{\hrule height 1pt} 
      \end{tabular}
     \smallemptyline     
      \caption{$\anD$}
   \end{minipage} 
\end{table}

\item
コンピュータの回路としては
$\anD$の否定の$\nand$,  
$\oR$の否定の$\nor$,
$\oR$で2つが共に$1$のときを$0$に代えた$\xor$(排他的$\oR$)
を考えると都合が良い．それらの真理表を書くと
\begin{table}[htb]
   \centering
       \begin{minipage}{0.4\columnwidth}
      \centering
      \begin{tabular}{cc|c}\noalign{\hrule height 1pt} 
         $A $&$B$& $A \oR B$\\  \hline
         0&0& 0\\
         0&1& 1\\
         1&0& 1\\
         1&1&1\\\noalign{\hrule height 1pt} 
      \end{tabular} 
      \smallemptyline 
      \caption{$\oR$}
   \end{minipage} 
   \begin{minipage}{0.4\columnwidth}
      \centering
      \begin{tabular}{cc|c}\noalign{\hrule height 1pt} 
         $A$ &$B$& $A\nand B$\\  \hline
         0&0& 1\\
         0&1& 1\\
         1&0& 1\\
         1&1&0\\ \noalign{\hrule height 1pt} 
      \end{tabular} 
      \smallemptyline   
      \caption{$\nand$}
   \end{minipage} 
\end{table}
%\newpage

\begin{table}[htb]
   \centering
   \begin{minipage}{0.4\columnwidth}
      \centering
      \begin{tabular}{cc|c}\noalign{\hrule height 1pt} 
         $A $&$B$& $A\nor B$\\  \hline
         0&0& 1\\
         0&1& 0\\
         1&0& 0\\
         1&1& 0\\ \noalign{\hrule height 1pt} 
      \end{tabular} 
      \smallemptyline 
      \caption{$\nor$}
   \end{minipage} 
   \begin{minipage}{0.4\columnwidth}
      \centering
      \begin{tabular}{cc|c}\noalign{\hrule height 1pt} 
         $A$ &$B$& $A\xor B$\\  \hline
         0&0& 0\\
         0&1& 1\\
         1&0& 1\\
         1&1&0\\ \noalign{\hrule height 1pt} 
      \end{tabular}   
      \smallemptyline  
      \caption{$\xor$}
   \end{minipage} 
\end{table}

\item
図式で考えた方が分かり易い場合があるので記号を準備する：
\begin{figure}[h]
   \centering
   \includegraphics{~/Documents/GitHub/introInfo/ch3/images/figure4.1}
   \caption{論理の図式}\label{}
\end{figure}

\item
数の$0$, 群での単位元，カテゴリー論での恒等射のように
%数学的な完全性のためには
何もしない論理素子$\idm$も考えると都合が良い．
それを以下のように表す．

\begin{table}[htb]
   \centering
   \begin{minipage}{0.4\columnwidth}
      \centering
      \begin{tabular}{c|c}\noalign{\hrule height 1pt} 
             \hline
         $A$ &$\idm A$ \\  \hline
         0&0\\
         1&1\\ \noalign{\hrule height 1pt} 
      \end{tabular}
      \smallemptyline 
      \caption{$\idm$}
   \end{minipage} 
   \begin{minipage}{0.4\columnwidth}
      \centering
      \begin{tabular}{cc|cc}\noalign{\hrule height 1pt} 
         $A$&$B$ &$S$&$C$ \\  \hline
         0&0&0&0\\
         0&1&1&0\\
         1&0&1&0\\
         1&1&0&1\\
        \noalign{\hrule height 1pt} 
      \end{tabular}
      \smallemptyline 
      \caption{二進数の和}\label{sum}
   \end{minipage} 
\end{table}

\item
一桁の二進数($0$, $1$のみからなる数)の和を計算し，
二進数で表す論理回路を設計してみる．
$A$, $B$が与えられた一桁の二進数とし，
和の一桁目を$S$と表し，位上がりを$C$と表すと
真理表は表~\ref{sum}のようになる．これから
$S$は$\xor$で$C$は$\anD$になることがわかる．
よって図~\ref{sumPic}
ようにすれば良い：
\begin{figure}[htb]
   \centering
   \begin{minipage}{0.4\columnwidth}
      \centering
      \includegraphics{~/Documents/GitHub/introInfo/ch3/images/figure4.2}
      \caption{}
   \end{minipage} 
   \begin{minipage}{0.4\columnwidth}
      \centering
      \includegraphics{~/Documents/GitHub/introInfo/ch3/images/figure4.3}
      \caption{一桁の二進数の和}\label{sumPic}
   \end{minipage} 
\end{figure}
\begin{Q}
   $\nand$ のみで$\noT$, $\anD$, $\oR$, $\nor$, $\xor$を表せ．
\end{Q}
\begin{exQ}
   二桁の二進数の二桁目の計算をして（一桁目の位上がりがあることに注意）二桁目と三桁目を二進数で表す論理回路を設計せよ．
  それを使って$n$桁の二進数の和を計算する論理回路を設計せよ．
\end{exQ}

\end{itemize}




\begin{itemize}
\item
論理回路のFET(Field Effect Transistor)などを使った物理的実現（論理ゲート）の例を挙げる：
\begin{figure}[htb]
   \centering
   \begin{minipage}{0.3\columnwidth}
      \centering
      \begin{circuit}0
      %NOT
         \nfet1 {?} G l % transistor
         \frompin nfet1S
         \- 1 d
         \gnd1
         \frompin nfet1D
         \- 1 u
         \.1
         \frompin .1
         \- 1 u
         \nl\R1 {$R$} u
         \- 1 u
         \cc\connection1  {$V$} c u
         \frompin nfet1G
         \- 1 l
         \cc\connection2  {$A$} c l
         \frompin .1 
         \- 3 r
         \cc\connection3  {OUT} c r
      \end{circuit}
      \caption{$\noT$}\label{NOT}
   \end{minipage} 
   \begin{minipage}{0.3\columnwidth}
         \centering
         \begin{circuit}0 
         %NAND
            \nfet1 {?} G l
            \frompin nfet1S
            \- 1 d
            \nfet2 {?} D l
            \frompin nfet2S
            \- 1 d
            \gnd1
            \frompin nfet1D
            \- 1 u
            \.1
            \frompin .1
            \- 1 u
            \nl\R1 {$R$} u
            \- 1 u
            \cc\connection1  {$V$} c u
            \frompin nfet1G
            \- 1 l
            \cc\connection2  {$A$} c l
            \frompin nfet2G
            \- 1 l
            \cc\connection3  {$B$} c l
            \frompin .1
            \- 3 r
            \cc\connection4  {OUT } c r
         \end{circuit}
      \caption{$\nand$}
   \end{minipage} 
   \begin{minipage}{0.3\columnwidth}
      \begin{circuit}0
         %NOR
         \nfet1 {?} G l % transistor
         \frompin nfet1S
         \- 1 d
         \gnd1
         \frompin nfet1D
         \- 1 u
         \- 2 r
         \.1
         \frompin .1
         \- 1 u
         \.2
         \frompin .2
         \- 1 u
         \nl\R1 {$R$} u
         \- 1 u
         \cc\connection1  {$V$} c u
         \frompin .1
         \- 2 r
         \- 1 d
         \nfet2 {} D r
         \frompin nfet2S 
         \- 1 d
         \gnd2
         \frompin nfet1G
         \- 1 l
         \cc\connection2  {$A$} c l
         \frompin nfet2G
         \- 1 r
         \cc\connection3  {$B$} c r
         \frompin .2 
         \- 3 r
         \cc\connection4  {OUT} c r
      \end{circuit}
      \caption{$\nor$}
   \end{minipage} 
\end{figure}

\noindent
ただし，細い四角は抵抗，$\Pi$みたいなものが FETである．

\item
例えば$\anD$ゲートは$\nand$ゲートと$\noT$ゲートの結合で実現される．$\nand$ゲートから$\noT$ゲートを作るのも簡単である．演習で見たように他の論理演算も$\nand$のみで表されるので集積回路は$\nand$ゲートのみから構成することが出来る．

\item
デコーダ：真理変数の組（これは二進数と同一視される）のデータから情報を取り出すには図~\ref{decode} (a) ような回路を使えば良い．
ただし，図の黒丸は(b)のような$\anD$を使った回路とする．

\begin{figure}[htb]
   \centering
   \includegraphics{~/Documents/GitHub/introInfo/ch3/images/figure4.4}
   \caption{二進数デコーダ}\label{decode}
\end{figure}

例として(c)のように$101$のデータを入力すると
赤の縦線に電圧がかかる．元々 横線には電圧がかかっているので$6$番目の線にのみ電圧がかかる．

\item
エンコーダ：逆に選ばれたワイヤーのデータを真理変数の組（二進数）
に変換するには，図~\ref{encode}のような一本のワイヤーを選ぶ装置を使えば良い．ただし，白抜きの四角は(b) のように$\oR$を使った回路とする．
例として(c)のように$7$番目の線に電圧をかけると$110$を取り出すことが出来る．

\begin{figure}[htb]
   \centering
   \includegraphics{~/Documents/GitHub/introInfo/ch3/images/figure4.5}
   \caption{二進数エンコーダ}\label{encode}
\end{figure}

\item
一般の論理回路：
一般の論理回路は, $0, 1$をとる$m$個の変数$X_i$ $(i=1,\dots, m)$
に$0, 1$をとる$n$個の変数$Y_j$ $(j=1, \dots, n$)を返す関数の組と考えられるので
\[
   Y_j = F_j(X_1, \cdots,X_i,\cdots,X_m)\quad j=1, \dots, n
\]
と表すことが出来る．
$\mathbf X=(X_1, \cdots,X_m)$, $\mathbf Y=(Y_1, \cdots,
Y_n)$, $\mathbf F=(F_1, \cdots,F_m)$と置く．
$ V^m=\{0,1\}^m$ とすると，
$\mathbf X\in V^m$, $\mathbf Y\in V^n$なので，$\mathbf F$は写像
 \[
    \mathbf F:V^m\to V^n
 \]
 とみなせる．
これは$m$バイトの二進数の集合とも見なせる（し，
$m$次元超立方体$[0,1]\times\dots\times[0,1]$の頂点の集合とも見なせる
(図~\ref{mathbf F} )）．
 つまり，$\mathbf F$は$m$バイトの二進数を$n$バイトの二進数に対応させる任意の写像とみなせる．
 
  \begin{figure}[htb]
   \centering
   \includegraphics{~/Documents/GitHub/introInfo/ch3/images/figure4.6}
   \caption{$\mathbf F:V^m\to V^n$}\label{mathbf F}
\end{figure}

この視点に立つと，二進数デコーダと二進数エンコーダを使ってこの写像$\mathbf F$を図 \ref{anyLogic} のように構成できることがわかる．
特に$\anD$と$\oR$, $\noT$だけで構成されることがわかる．
ただし，これは$\mathbf F$を実現する論理回路の存在を示したに過ぎず，
実際の回路の設計ではいかに効率よく資源を配分するかが大切になる．

 
\begin{figure}[htb]
   \centering
   \includegraphics{~/Documents/GitHub/introInfo/ch3/images/figure4.7}
   \caption{任意の論理回路の実現}\label{anyLogic}
\end{figure}
\end{itemize}


\section{フリップフロップ回路とメモリ}
\begin{itemize}
\item
\term[]{フリップフロップ回路}（双安定バイブレーターとも言う）とは図~\ref{flipflop}のような回路である．
$\bar Q$は$Q$の否定，$S$はセット，$R$はリセットで同時に真になることはない変数とする．
\begin{figure}[htb]
   \centering
   \includegraphics{~/Documents/GitHub/introInfo/ch3/images/figure4.8}
   \caption{フリップフロップ}\label{flipflop}
 \end{figure}
 \begin{table}[htb]
        \centering
         \begin{tabular}{cc|cc}\noalign{\hrule height 1pt} 
            $S$&$R$ &$Q$&$\bar Q$ \\  \hline
            0&0&1&0\\
            0&0&0&1\\
            1&0&1&0\\
            0&1&0&1\\
            \noalign{\hrule height 1pt} 
         \end{tabular}
         \smallemptyline 
         \caption{フリップフロップ}
\end{table}

\begin{figure}[htb]
   \centering
   \includegraphics{~/Documents/GitHub/introInfo/ch3/images/figure4.9}
   \caption{$S=0,R=0$の安定性}\label{stab}
\end{figure}

\begin{figure}[htb]
   \centering
   \includegraphics{~/Documents/GitHub/introInfo/ch3/images/figure4.10}
   \caption{$S=1\Rightarrow Q=1$%, $R=1\Rightarrow \bar Q=1$
   }\label{setReset}
\end{figure}

\item
この回路の応答を調べる．図~\ref{stab}から分かるように$R=S=0$のとき$Q$, $\bar Q$の状態を保つ，ここでは左右二つの図を描いたが上下を入れ換えれば同じ図なので以後一方のみ書くことにする．
図~\ref{setReset}から分かるように, 
右の図では回路は安定だが，
$S=1$ のときは左の図の状態は不安定で
安定な右の図の$Q=1$に遷移することが分かる．

\item
状態が移るとき時間がかかるので，それを考慮して
図~\ref{pulseWave}のような時計パルス波$\phi$で同期することを考える，
ただし$t$は時間である．$\phi$に$\noT$したものを$\bar\phi$と呼ぶ．
このパルス波を生成するには図~\ref{pulseCircuit} のような回路を考えれば良いことが知られている．抵抗$R$キャパシティ$C$のとき$0.7CR$秒間の周期でパルス波を生成する．
この回路も双(非安定)バイブレーターと呼ばれている．

\begin{figure}[htb]
   \centering
   \begin{minipage}{0.5\columnwidth}
      \centering
      \includegraphics{~/Documents/GitHub/introInfo/ch3/images/figure4.11}
      \caption{時計パルス}\label{pulseWave}
   \end{minipage} 
   \begin{minipage}{0.4\columnwidth}
      \begin{circuit}0
         \npn1 {?} B r
         \frompin npn1E
         \- 1 d
         \gnd1
         \frompin npn1B
         \- 1 r
         \- 4 u
         \.1
         \- 1 u
         \nl\R1 {$R$} u
         \- 1 u
         \.2
         \frompin npn1C
         \- 1 u
         \.3
         \- 2 u
         \nl\R2 {$R$} u
         \frompin R2t
         \vtopin .2
         \htopin .2
         \frompin .2
         \- 2 r
         \.4
         \- 2 r
         \.5
         \- 1 d
         \nl\R3 {$R$} d
         \- 2 d
         \.6
         \- 3 d
         \- 1 r
         \npn2 {?} B l
         \frompin npn2E
         \- 1 d
         \gnd2
         \frompin npn2C
         \- 2 u
         \.7
         \- 1 u
         \nl\R4 {$R$} u
         \frompin R4t
         \vtopin .5
         \htopin .5
         \frompin .3
         \- 1 r
         \nl\C1  {$C$}  r
         \frompin C1r
%         \vtopin .6
         \htopin .6
         \frompin .7
         \- 1 l
         \nl\C2  {$C$}  l
         \htopin .1
         \frompin .7
         \- 1 r
         \cc\connection1  {$\phi$} c r
         \frompin .4
         \- 1 u
         \cc\connection1  {V} c u
         \frompin .3
         \- 1 l
         \cc\connection1  {$\bar\phi$} c l
      \end{circuit}
      \caption{パルス波生成回路}\label{pulseCircuit}
   \end{minipage} 
\end{figure}


%以上の性質からフリップフロップ回路をメモリーとして使うことができる．そのためまず
\item
図~\ref{clockedFF}のように$\anD$回路を使ってフリップフロップ回路にパルス波を加えると時計付きフリップフロップ回路が得られる．以下右の略記号で表す．
\begin{figure}[htb]
   \centering
   \includegraphics{~/Documents/GitHub/introInfo/ch3/images/figure4.12}
   \caption{時計付きフリップフロップ回路}\label{clockedFF}
\end{figure}
\newpage

\item
この回路の応答を図~\ref{clockedFF2}で調べる．
\begin{figure}[htb]
   \centering
   \includegraphics{~/Documents/GitHub/introInfo/ch3/images/figure4.13}
   \caption{時計と回路の動き}\label{clockedFF2}
\end{figure}

\item
時計付きフリップフロップ回路を二つ並べてパルス波で同期を取ると
Dフリップフロップ（DFF）が得られる．
\begin{figure}[htb]
   \centering
   \includegraphics{~/Documents/GitHub/introInfo/ch3/images/figure4.14}
   \caption{Dフリップフロップ}
\end{figure}

\newpage
\item
この動作を調べて見るとDの意味もわかる 
\begin{figure}[htb]
   \centering
   \includegraphics{~/Documents/GitHub/introInfo/ch3/images/figure4.15}
   \caption{DFFの動き}
\end{figure}

\item
次にDFFを幾つか並べるとメモリを構成できる．
\begin{figure}[htb]
   \centering
   \includegraphics{~/Documents/GitHub/introInfo/ch3/images/figure4.16}
   \caption{DFFによるメモリ}
\end{figure}

\item
このDFFメモリにデータを格納することをシュミレートしてみる．
まず，全てが$0$になるように初期化する．
次に入力する二進数を$101$とする．
\begin{figure}[htb]
   \centering
   \includegraphics{~/Documents/GitHub/introInfo/ch3/images/figure4.17}
   \caption{D-FFによるメモリ格納}
\end{figure}

\begin{Q}
以下の回路でもフリップフロップ回路が実現できることを真偽表を書いて示せ．
\begin{figure}[htb]
   \centering
   \includegraphics{~/Documents/GitHub/introInfo/ch3/images/figure4.18}
   \caption{フリップフロップ回路}
\end{figure}
\end{Q}
\end{itemize}

%\begin{thebibliography}{9}
%   \bibitem{F}
%      Feynman, Lectures on computation, Perseus, 1996, 
%      日本語版：ファインマン計算機科学, 岩波書店, 1999.
%
%  
%   \bibitem{Mn} Minsky, Computation: Finite and Infinite Machines, 1967.
%   
%   \bibitem{Mu}三浦伸夫, 数学の歴史, 放送大学教育振興会, 2019.
%   
%    \bibitem{T}
%    田中靖浩, 会計の世界史, 日本経済新聞出版, 2018.
%\end{thebibliography}
%
%\end{document}
%
%\section{有限オートマトン}
%オートマトン
%
%\section{チューリングマシン}
%チューリングマシン，
%ユニバーサルチューリングマシン

%\begin{circuit}0
%\nfet1 {?} B l % transistor
%\frompin nfet1C % draw from collector
%\- 1 u % some wire
%\nl\A1 {$I_C$} u % amperemeter for current of collector
%\atpin nfet1B % continue drawing from base
%\- 1 l % some wire
%\R1 {510} l % resistor
%\- 1 l % some wire to the edge
%\centerto A1 % draw centered to amperemeter 1
%\nl\A2 {$I_B$} u % amperemeter 2
%\frompin A2b % link amperemeter 2 with resistor
%\vtopin R1l
%\frompin A1t
%\- 1 u
%\.1 % junction
%\frompin A2t % wire to amperemeter 2
%\vtopin .1
%\htopin .1
%\- 1 u
%\cc\connection1 {$U_b$} c u % driving voltage
%\frompin nfet1E
%\- 1 d
%\GND1 % ground
%\end{circuit}


\section{有限オートマトン}
前回説明した論理回路を抽象化したのが
\term[]{オートマトン}である(歴史的には多分オートマトンが先で論理回路が後であろう)．
\begin{itemize}
   \item
      まず例：
      $0$, $1$の列からなる入力の総和が偶数か奇数かを判定し
      順次そこまでが偶数か奇数を書き出すオートマトンは
      以下のように表される：
      \begin{figure}[htb]
         \centering
         \includegraphics{~/Documents/GitHub/introInfo/ch3/images/figure5.1}
         \caption{偶奇判定機械}
      \end{figure}
      
      〇で囲まれているのを状態という．
      この場合$0$と$1$が状態で，
      それぞれその時点での偶数と奇数の状態を表している．
      start は入力を受け付ける前の状態で
      この場合は$0$から始まることを表してる．
      矢印は入力が根元の数字のとき状態がどう移るかを表し，
      矢印の真ん中の部分の数字を出力するとする．
      この場合は$011010$のような入力を順次受け付けると，
      $0$ならば状態を変えず元の状態を書き出し，
      $1$ならば状態を替え異なる状態を書き出すことを表している．
      
   \item
      次の例：二進法の入力にその$2$個前の入力を出力する
      オートマトンを設計してみる．
      前$2$個の入力を
      $00$, $01$, $10$, $11$とする状態を考える．
      最初は入力が無いので$00$であると見なす．$0$と$1$が入力される度に左にそれが付け加わり，右側の$0$か$1$を出力すれば良い．
      従って，以下のようなものを考えれば良い．
      \begin{figure}[htb]
         \centering
         \includegraphics{~/Documents/GitHub/introInfo/ch3/images/figure5.3}
         \caption{$2$遅延機械}
      \end{figure}
      
      \begin{Q}
         二進法の入力にその$1$個前の入力を出力するオートマトンを設計せよ．
      \end{Q}
      
   \item
      オートマトンを数式で表してみる．
      $t=t_i=i\in \mathbf N$を時間のパラメータとし，
      $Q=\{q_1,\dots,q_L\}$を状態の集合, 
      $S=\{s_1,\dots,s_M\}$を入力の集合
      $U=\{u_1,\dots,u_N\}$を出力の集合とする．すると入力は
      \[
      	s:\mathbf N\to S
      \]
      と言う任意の関数で, 状態と出力の時間による変化は
      \[
      	q:\mathbf N\to Q, \quad u:\mathbf N\to U
      \]
      という関数で表される．
      \[
%      \begin{split}
%         F:Q\times S\to U\quad&\text{：出力の規則を示す関数}\\
%         G:Q\times S\to Q\quad&\text{：状態の遷移の規則を表す関数}\\
%         q(1)\in Q \quad&\text{：初期条件}
%      \end{split}
    \begin{split}
         F:Q\times S\to Q\quad&\text{：状態の遷移の規則を表す関数}\\
         G:Q\times S\to U\quad&\text{：出力の規則を示す関数}\\
         q(1)\in Q \quad&\text{：初期条件}
      \end{split}
     \]
      がオートマトンであり，時系列
      \[
         q(t_{i+1})=F(q(t_i),s(t_i)),\quad
         u(t_{i+1})=G(q(t_i),s(t_i))
      \]
      が出力と状態遷移であることが分かる．

   \item   
      このように定式化すると，
      この$F$, $G$を持つオートマトンが，前回紹介したクロック周波数で同期したメモリ（これは上の演習の$1$遅延機械でもよい）と
      論理回路により表現できることが分かる：
      まず，$S$, $U$, $Q$を二進法で表現しておき，$k$-ビット, $l$-ビット, $m$-ビットとし，その数のワイヤーでその値を表わしているとする．
      以下のようなクロック周波数$\phi$と$\bar\phi$ で同期された$2$つのメモリと$F$, $G$を実現する論理回路を組み合わせることでオートマトンが実現される．
      \begin{figure}[htb]
         \centering
         \includegraphics{~/Documents/GitHub/introInfo/ch3/images/figure5.4}
         \caption{一般のオートマトンを実現する論理回路}
      \end{figure}
   
\end{itemize}

\section{チューリングマシン}
\begin{itemize}
   \item
   \term[]{チューリングマシン}とは，
   箱状のセルを一列に刻んだテープのある一個のセルに
   オートマトンが付いたもの(図~\ref{TuringM})である．
   このオートマトンを\term[]{ヘッド}と言ったりする．
   テープの各セル上には最初に有限個の記号が書かれていて，
   空白を$0$と表すとある所から右，
   ある所から左はすべて$0$になっているものである．
   時間は離散的な時計に従ってすすみ，
   最初に指し示すセルとヘッドの状態は予め定まっている．
   ある時刻のオートマトンの状態と
   それの指しているセルに書かれている記号によって
   オートマトンは次の状態に移り，
   その箱に何かの記号を書き込み，
   右か左に一つだけ動き時間が進む．
   ある決まった状態に移ると機械は停止する．
   \begin{figure}[htb]
      \centering
      \includegraphics{~/Documents/GitHub/introInfo/ch3/images/figure5.5}
      \caption{チューリングマシン}\label{TuringM}
   \end{figure}
   
   \item   
   チューリングマシンの例：単なるオートマトンとは異なることを
   先程と同様に例をもって示したい．
   問題は次のような$($ と$)$の有限個の列が与えられたとき，
   それが$($ と$)$の対応として整合性があるかないかを判定することである．
   もちろん整合性があれば$($ と$)$の数は等しくなるので，
   $($ と$)$の数が等しくなければ整合性はないことがわかるが
   それでは十分ではない．また，任意個数の$($と$)$が表れるので
   有限状態のみでは判定できない．
   
   \begin{figure}[htb]
      \centering
      \includegraphics{~/Documents/GitHub/introInfo/ch3/images/figure5.6}
      \caption{初期状態}\label{}
   \end{figure}
   まず，テープの初期状態は調べたい$($ と$)$の組み合わせを一個一個
   セルに入れて両端を$E$ではさんでいるとする．
   そしてヘッドは左側の$E$を指しているとする．
   図に書いた$R$は右に，
   $L$は左にヘッドが動くことを表する．
  左から ペアの最も近い$($ と$)$を見つけまず$)$を
  関係ない記号$X$に書き換え左に動き$($を$X$に書き換え右に動く
  ことを実現するオートマトンは図~\ref{consist}のようになる．
  この入力だとまず一つ右に動き$($は無視してもう一つ右に動くと
  $)$があるのでそれを$X$に書き換え
   左に動き$($があるのでこれを$X$に書き換え右に動く．
    \begin{figure}[htb]
      \centering
      \includegraphics{~/Documents/GitHub/introInfo/ch3/images/figure5.7}
      \caption{括弧整合性判定機械のヘッド}\label{consist}
   \end{figure}
  
   終了の仕方について：
   まず$)$を$X$に書き換えた後左にすすんで何も見つからず
   $E$に到達した場合，整合性がないので$0$を出力して停止する．
   そうでなければ，
   対になる$($が全て見つかり，
   これを続けると右の$E$に到達することになる．
   そのあと左に移動して
   $($が残ってれば整合性がないので$0$を出力して停止する．
    $($が残っていなければ$E$に到達し整合性が確認され
    $1$を出力して停止する．
   
   \item
      オートマトンとリューリングマシンの違いは
      オートマトンの所で述べた$F:Q\times S\to Q$, 
   $G:Q\times S\to U$の他にテープの移動方向を与える関数
      \[ 
             D:Q\times S\to \{0,1\}
      \]
      を指定することである．$0$を左，$1$を右と決めておく．
    \item
        インデックスの付いたファイルを検索して，
      その内容を適当な場所にコピーする
      チューリングマシンを作成する．
      
      まず，テープのデータの格納状況を説明する．簡単のため
      インデックスとファイルの内容は共に$3$ビットの$2$進数で
      $X$と$X$(または$Y$)の間の前半$3$ビットがインデックスで，
      後半$3$ビットがファイルの内容である(図~\ref{initial})．
      左の$Y$がファイルの先頭で右の$Y$がファイルの終を表す．
      左側の端の$Y$と最も左の$X$との間の$3$ビットが探したい
      ファイルのインデックスであり，オートマトンは最初に
      最も左の$X$を指した状態にあるとする．
   \begin{figure}[htb]
      \centering
      \includegraphics{~/Documents/GitHub/introInfo/ch3/images/figure5.8}
      \caption{初期状態}\label{initial}
   \end{figure}
   
   オートマトンは次の図~\ref{file} のように設計される：
   \begin{figure}[htb]
      \centering
      \includegraphics{~/Documents/GitHub/introInfo/ch3/images/figure5.9}
      \caption{ファイル探索機械}\label{file}
   \end{figure}
   \begin{figure}[htb]
      \centering
      \includegraphics{~/Documents/GitHub/introInfo/ch3/images/figure5.10}
      \caption{探索中のテープの変遷}\label{tape1}
   \end{figure}

   \noindent 
   まずヘッドは左に進みながら数字の$0$を$A$へ$1$を$B$へ書き換え，
   $Y$に到達すると右に進み始める．
   まず最初の$A$を$0$へ$B$を$1$へ書き換え，
   次の$A$, $B$や$X$を無視して最初の数字まで右に進む．
   元の数字が$0$のときその数字が$1$，元の数字が$1$のときその数字が$0$ならば
   インデクスは一致していないので，次の右の$X$まで進み，
   そこから左に進みながら$0$を$A$へ$1$を$B$へ書き換え$Y$まで進み，
   次のインデックスの比較に移る．
   元の数字が$0$のとき最初の数字が$0$，元の数字が$1$のときその数字が$1$ならば，
   そこまではインデクスが一致しているので，その$0$を$A$へ$1$を$B$へ書き換え
   もう一度$Y$まで左に戻って，つぎの桁の数字の比較を行う．
   これを続け，完全に一致したとき，ヘッドは最も左の$X$まで動く．
   もし途中で一致しない場合は次の右の$X$まで進み，
   そこから左に進みながら$0$を$A$へ$1$を$B$へ書き換え
   $Y$まで進み，次の(上の$X$の左の)インデックスとの比較に移る．
  
  次にそのインデックスのデータを$Y$と最も左の$X$の間の部分にコピーする
  以下のようなオートマトンを考える：
  \begin{figure}[htb]
      \centering
      \includegraphics{~/Documents/GitHub/introInfo/ch3/images/figure5.11}
      \caption{コピー機械}\label{}
   \end{figure}
 
    \begin{figure}[htb]
      \centering
      \includegraphics{~/Documents/GitHub/introInfo/ch3/images/figure5.12}
      \caption{コピー中のテープの変遷}\label{tape2}
   \end{figure}

   検索が終わった後最も左の$X$とデータの間は全て$A$, $B$に置き換わっているので，
   右側の最初の数がデータの数字である．そこで$X$から右側に進み
   最初の数が$0$のとき$A$へ$1$のとき$B$へ書き換え
   左へ進み$Y$に戻る．
   $Y$から右に進み上の最初の数を元の数が$0$のとき$A$へ$1$のとき$B$へ書き換え，
   最も左の$X$まで右に進み，
   次の桁の数まで右に進み，
   次の桁の数をコピーしてそれを繰り返し終るまで続ける．
   ただし，このチューリングマシンでは
   最後は意味の無い次のファイルの最初の数をコピーして，
   $Y$と$X$の間にもうコピーする場所が無いので
   最も左の$X$で停止することになる．
 \end{itemize}

 
\section{ユニバーサルチューリングマシン}
全てのチューリングマシンの動作を完全にシミュレートできる
チューリングマシンを
\term[]{ユニバーサルチューリングマシン}という．
$T$がシミュレートされるチューリングマシンとし，
$U$をユニバーサルチューリングマシンとする．
   
簡単のため$T$の状態は$2$ビットとし，
$T$のテープが$Y$の左に再現されている．
そこのセルには$0$, $1$のみが書かれてるとするが
現在作業中のセルに$M$という記号が書かれているとする．

最も左の$X$の後にあるのは遷移状態のリストである：
状態とセルの値との組で$3$ビット，
次の状態の$2$ビットと書き込む値が$1$ビット，
ヘッドの移動向きを示す$1$ビットの組であり，
それら全てを最も左の$X$から左に$7$ビットずつ$X$で分けて
書いておき，最後に右の$Y$で終わるようにする．
左側の端の$Y$と最も左の$X$との間の$3$ビットが
$T$の状態とテープに書き込む数字を表している．

 \begin{figure}[htb]
   \centering
   \includegraphics{~/Documents/GitHub/introInfo/ch3/images/figure5.13}
   \caption{$T$のセルへの値の書き込みと移動と読み込み}\label{tape3}
\end{figure}

図~\ref{UT}にこのユニバーサルチューリングマシンの
ヘッドの構造を示す(\cite{Mn})．上の部分はファイル探索機械と同じであるが，
これはインデクスを$T$の状態とテープから読み込まれた数の組として，
次の遷移状態，つまり，
$T$の次の状態とテープへ書き込む数とヘッドの動く方向を探している
ことになる．
次の右の中部の部分はコピー機械と同じである．ただし
最後にコピーした数字はヘッドの動く方向を示すビットで
この場合は右を表している．

右の中部の下の部分では$X$から$M$まで左に進み，
まずヘッドの移動方向をそこに書き込む．
この場合は$B$である．
次に右に進み$A$を$0$, $B$を$1$に置き換えながら
最も左の$X$まで進む．

その$X$の左の$A$と$B$をもう一度$0$, $1$に戻すために，
最も右の数まで進み$A$を$0$に$B$を$1$に書き換えながら
最も左の$X$まで左に進む．
もう一つ左にすすみ最初の数字を読み込み, 記号$V$に置き換え，
$T$のテープの$M$の有ったところまで左に進む．
もし，そこに$A$があれば，そこに先程読み込んだ数字を書き込み
左に動く．
$B$があれば，そこに先程読み込んだ数字を書き込み右に動く．
そしてヘッドの動いた先の数字を読み込み，最も左の$X$まで進み
一つ左に進み$V$の書かれたセルに読み込んだ数字を書き込み，
最後に右に一つ$X$まで戻る．

これで一サイクルが終わり次のサイクルが始まる．
$T$が停止するまで続いていく．

\begin{figure}[htb]
   \centering
   \includegraphics{~/Documents/GitHub/introInfo/ch3/images/figure5.14}
   \caption{Minskyのユニバーサルチューリングマシン}\label{UT}
\end{figure}

このオートマトンは記号$8$個，状態$24$個からなっている．
\begin{Q}
   このユニバーサルチューリングマシンがちゃんと動くか
   確認して下さい．
\end{Q}


最後にチューリングについては伝記とその映画化(\cite{TuringM})がある．
数学も凄いが人生も凄い人なので興味のある人には鑑賞してもらいたい．

%\section{計算可能性}

%
%\begin{thebibliography}{9}
%   \bibitem{F}
%      Feynman, Lectures on computation, Perseus, 1996, 
%      日本語版：ファインマン計算機科学, 岩波書店, 1999.
%  
%   \bibitem{Mn} Minsky, Computation: Finite and Infinite Machines, 1967.
%   
%   \bibitem{Turing}
%       A.~Turing, On Computable Numbers, 
%       with an Application to the Entscheidungsproblem, 
%       Proceedings of the London Mathematical Society 42, 1936.
%
%
%   \bibitem{TuringM}
%      アンドルー・ホッジス, エニグマ アラン・チューリング伝 上下,
%      勁草書房, 2015.
%      映画化：イミテーション・ゲーム エニグマと天才数学者の秘密,
%      モルテン・ティルドゥム監督, 2014.
%\end{thebibliography}


%\documentclass[a4paper,12pt]{amsart}
%\usepackage[dvipdfmx]{graphicx,xcolor}%graphics
%\usepackage{float}
%\usepackage{amsmath}
%\usepackage{verbatim}
%\usepackage{latexsym}
%\usepackage{amssymb}
%%\usepackage{emath}
%\usepackage{listings}
%\usepackage{ccicons}
%\usepackage{mflogo}
%%\usepackage{fancybox}%enclose by box
%\usepackage[dvipdfmx]{hyperref}
%\usepackage{verbatim}
%\hypersetup{% hyperrefオプションリスト
%   setpagesize=false,
%   bookmarksnumbered=true,%
%   bookmarksopen=true,%
%   colorlinks=true,%
%   linkcolor=red,
%   citecolor=blue
%} 
%\newtheorem*{exQ}{レポート}
%\theoremstyle{remark}
%\newtheorem*{Q}{演習}
%
%\renewcommand{\figurename}{図}
%\renewcommand{\proofname}{証明}
%\renewcommand{\tablename}{表}
%
%\def\emptyline{\vspace{11pt}}
%\def\smallemptyline{\vspace{6pt}}\begin{document}
%\begin{flushright}
%   2026/08/27
%   %2026/06/02
%   %2026/04/16
%   %2026/03/30
%   %2026/01/12
%   
%   九大・数理 大津幸男
%\end{flushright}
%
%\begin{center}
%   {\Large 情報科指導法II---06}
%   \href{https://creativecommons.org/licenses/by-nc-nd/4.0/}{\ccbyncnd}
%   \par{インターネット1}
%\end{center}
%

\chapter{インターネット}
ネット net とは網(あみ)のことであり，network とは網全体，例えば，駅と線路と列車の全体，
テレビ局の全国放送網のことをさすが，
インターネット The Internet はコンピュータの全世界的な network のことである．

\section{インターネットの仕組み}
\subsection{コミュニケーションとメディアの歴史}
昔から狼煙・太鼓を使ったり，手紙によるコミュニケーションや
徒歩や馬による移動などが行われていた．
しかし，近代になって重要になった人の移動や人口の変化，
貨幣の循環を含むコミュニケーションやメディアの進展には
以下のようなものがある(世界史の復習に一冊あげるなら\cite{Yamakawa})：

\begin{itemize}
   \item 
      13〜15世紀のモンゴル帝国がユーラシア大陸に広がりと交易が発展した．
   \item 
     それがきっかけとなり14世紀に黒死病(ペスト)がヨーロッパに広がり，
     ヨーロッパでは3人に1人が死んだと言われている．
   \item 
      東ローマ帝国が1453年にオスマン帝国に滅ぼされ多くの文化人が西洋に亡命し，
      その人達の活動とイスラムとの接触を通して
      ギリシャ・ローマの文化が復活し
      ルネッサンスが始まるきっかけになった．
   \item 
      1450年代グーテンベルグによる活版印刷の開始．
  \item
     モンゴル帝国などから伝わった航海術が十字軍などを通して発展し，
     それを使った1492年\footnote{この年がスペインによるレコンキスタの完成の年であるのは偶然では無い．}
     のコロンブスによる新大陸の（ヨーロッパ人から見た）発見．
     それに続くヨーロッパからの南北アメリカ大陸への移民の流入．
     
  \item
     これ以降ヨーロッパから持ち込まれた伝染病で新大陸の多くの現地人\footnote{一説には10人中9人}が死亡した．
     それによるプランテーションでの労働力不足のために，
     17〜19世紀に奴隷貿易が行われ，
     アフリカの多くの黒人が奴隷としてアメリカ大陸に移住させられた．

  \item
      1545年の南米のポトシ銀山の発見のあと
      ヨーロッパへ銀が大量に流入し貨幣供給量が大幅に増加した．
      これが価格革命・商業革命をもたらした．
   \item 
      会社組織・簿記・銀行の発展．
      株式・債券・外国為替など様々な市場を含めた資本主義の発展．
   \item 
      新聞・雑誌等の発行によるマスコミの発達．
   \item 
      石炭を使った(第一次)産業革命．
   \item 
      蒸気機関車・自動車・道路・飛行機などの交通の発展．
   \item 
      写真・レコードの発明．
   \item 
      石油やガスの使用の増加\footnote{今回のイラン戦争で
      化学肥料も石油に依存していることが明らかになった．}．
      電(磁)気が発見され，十九世紀末には送電網が整備された(第二次産業革命)．
      以下の発明も電気があってのものである．
 
    \item 
      モールス S.~Morse による電信の発明．    

   \item 
      ベル G.~Bell による電話の発明．
      
   \item 
      マルコーニ G.~Marconi による無線の発明．
   
   \item 
      映画・ラジオ・テレビ放送の発展．
   
   \item 
      コンピュータの発明と普及．
 
    \item
      パソコンの普及．
        
   \item
      携帯電話(スマートホン)の発展．

   \item
      インターネットの発展．
\end{itemize}

\subsection{インターネットの歴史}

1960年代に米国国防総省の
ARPA(=Advanced Reseach Projects Agency, 国防高等研究計画局) で
ARPANET の計画が始まったが，次の二つの見方がある．

\begin{enumerate}
  \item 1960年代後半に，米国が核攻撃を受けてもグラフ状のネットの壊れた部分をパケット が自動的に迂回して繋がるような仕組みのARPANET の計画が始まった．
  \item 色々な種類のネット技術が分散している現実の世界でも，お互いに問題なく繋げられるネットワークを作りたいという理想に燃えた当時の技術者が internet の技術を開発した．
\end{enumerate}
ここではこれについての明言は避け\footnote{ \href{https://snow00two.github.io/mySite/}{私のポートフォリオ}にAIの見解や資料があります．}，
1969年にBerkley大学のコンピュータからStanford  大学のコンピュータへ
送信されたのがその最初の実験であったことのみ言及しておく．
その後1990年代に商用に広がり\term[]{インターネット}(the Internet)と呼ばれるようになった．

ただし，インターネットというネットが一つあるわけではなく，
家や会社や大学などのLAN =Local Area Network と大きな領域を繋ぐ
WAN=Wide Area Network の緩やかな塊がインターネットである．
物理的なネットの仕組みは色々バラエティーがあるが，
共通の決まりでパケットを送合い受けとるようになっている．

1秒間に何れだけの情報が送れるかというのがbps =bit per second 
という単位で計られる．通信容量ともいいい
通信の性能を測るため後で出てくる．

\subsection{LANや携帯電話}
\begin{itemize}
   \item  
         \begin{figure}[htb]
         \centering
         \includegraphics[width=10cm]{~/Documents/GitHub/introInfo/ch4/images/macaddress.pdf}
         \caption{MACアドレス}\label{macAdress}
      \end{figure}
      
      \begin{figure}[htb]
         \centering
         \includegraphics[width=8cm]{~/Documents/GitHub/introInfo/ch4/images/hub.pdf}
         \caption{ハブ}\label{}
      \end{figure}

      LANの具体的なものとして Ethernet\footnote{Etherとは相対性理論で否定されたエーテルのことである．もっとも否定されたのは古典論に従うEtherであって相対性理論的Etherは存在していても問題は無い．これも\href{https://snow00two.github.io/mySite/}{私のポートフォリオ}にAIの見解や資料があります．}
      がある．
      まず，全てのEthernetボードにMACアドレス(Media Access Control address)という一意的な12ケタの16進数が与えられている．私のmac の外付けEtherボードだと(図~\ref{macAdress})
      \[
         \text{24:1b:7a:f8:fb:43}
      \]
      になっている．
      
      あるエリアのパソコンやデバイスはイーサーケーブルとハブで結びついていて，
      エリア内のパソコンたちはそれを使って通信する．
      そのエリア内からルーターを通して外のWANに繋がっている．
      基本的な通信の基本はパケット通信で
      模式で表すと図~\ref{ethernet}のようなビット列である．
      ここでのアドレスとはMACアドレスのことである．
      初期のイーサネットの容量は3Mbps程度だったが
      最近では100Gbpsを越えるものも多い．
      \begin{figure}[htb]
         \centering
         \includegraphics{~/Documents/GitHub/introInfo/ch4/images/figure6.1}
         \caption{イーサーネットのパケットの構造}\label{ethernet}
      \end{figure}
      
      
      %PARC で開発された技術の一つである．
      
      
   \item       
      LANを無線で繋がるようにした無線LANを使うことも多い．
      IEEE 802.11のような規格で呼ばれるが，最近はWi-Fiということが多い．
      2.4〜2.5GHz, 5GHz 以上の周波数で動作する．
      無線LANは誰でも盗聴できるので 
      WPA2/WPA3 のような暗号化を用いることが多い．
      
       \begin{Q}
         自分の使っているEther器機(例えばWi-Fi)の
         Macアドレスを調べて下さい．
         近くの友達と比べて確かに異なることを確かめて下さい．
      \end{Q}
     
 
    \item       
        BlueToothなども2.4GHzの周波数でデバイスとパソコンの通信を行っている．
     
   \item       
      携帯電話は英語で Cellular Phone というが，
      Cell というのは平面上にある基地局を中心に平面を分割したものの一つ一つのことである．
      使用可能な周波数帯を隣りの局で混信しないように配置するが，
      ある程度遠くの局では混信しないので同じ周波数帯を再利用していい．
      更にその周波数帯を細かく分割して全ての携帯電話と基地局は連絡を取り合っている．
      携帯電話には国際移動体装置識別番号という15桁の識別番号が
      割り当てられていて，電源が入るとIDがブロードキャストされ
      携帯電話がどこかのセルにあると最寄りの基地局で
      その情報を把握して，電話会社が把握し
      必要に応じて電話回線やインターネットと繋ぐように
      運用されている．携帯がが動くに伴い基地局も変更される．
      無線LANがある所では安い方を使うようになっている．
      \begin{figure}[htb]
         \centering
         \includegraphics[trim=80 0  50 0, clip]{~/Documents/GitHub/introInfo/ch4/images/figure6.6}
         \caption{携帯電話基地とCellと携帯電話}\label{CellularPhone}
      \end{figure}
      
      \begin{table}[h]
         \centering
         \begin{tabular}{llll} 
              \noalign{\hrule height 1pt}
               世代&コメント&速度&年  \\ \hline \hline
               1G &アナログ通信&&  1980年代 終了\\ \hline 
               2G&デジタル通信& 最大64kbps &1990年代 終了\\ \hline
               3G &国際的な規格&最大3.6Mbps&2000年代 最近終了  \\ \hline
               4G &現在主流&最大1Gbps&2010年代\\ \hline
               5G &次世代&最大10Gbps& 現在\\ \hline
             \noalign{\hrule height 1pt}
         \end{tabular}
         \smallemptyline
         \caption{携帯電話の通信方法の変遷}
      \end{table}
\end{itemize}

\subsection{TCP/IP}
\begin{itemize}
   \item 
      \term[]{プロトコル}(protocol)とは
      もともと外交における議定書のことであるが，
      コンピュータのネットワークでは全体で守る規約のことである．
      特にインターネットではTCP/IPという規約で代表される
      ものである．
      これもパケット転送技術である．その構造を図で表すと図~\ref{levels}のようになる．

      \begin{figure}[htb]
         \centering
         \includegraphics{~/Documents/GitHub/introInfo/ch4/images/figure6.2}
         \caption{パケットの階層}\label{levels}
      \end{figure}
      
      \noindent
      物理層とは前のセクションで紹介したようなLANなどの通信の最も下層のレーヤーのことである．
      上の層が下の層のデータの中に埋め込まれているイメージで捉えてほしい．

   \item 
      \term[]{IP address}(Internet Protocol address)  は数字の並びで表した各コンピュータの名前である．
      IPV4では . で区切られた8bit の10進数($0\text{〜}255$)4個で表されている．よって，$2^{32}\approx 43$億ほどのアドレスしかない．私のマックだと(図~\ref{ipAdress})
      \[
         172.24.24.95
      \]
      \noindent
      になっている．
      
      \begin{figure}[htb]
         \centering
         \includegraphics[width=10cm]{~/Documents/GitHub/introInfo/ch4/images/ipaddress.jpg}
         \caption{IPアドレス}\label{ipAdress}
      \end{figure}
      
      
      当初はコンピューター一台毎にGlobal なIPアドレスを配布する予定であったが，
      私が使っているのはGlobal なIPアドレスではなく，
      プライベートアドレスで我々のローカルエリアでのみ使える
      (正確には同じアドレスが使い回されているが，そのローカルエリア以外からはアクセスできない)アドレスになっている．
      しかも，\term[]{DHCP}(Dynamic Host Configuration Protocol )
      という方法でアドレスを配っていて毎回使うアドレスは異なっている．
      通信のためには自分のIPアドレスと送りたいコンピュータのIPアドレスの組が必要になることにも注意する．
      
      IPV4を拡張したIPV6では16バイトの16進数を2バイトづつ:で区切って書く．例えば
      \[
         \text{2001:0db8:bd05:01d2:288a:1fc0:0001:10ee}
      \]
      これだと$2^{128}$個(日本語での十進法での単位が不明なので書かない)で十分な数がある．
      
      いずれにしろ人間は大きい数字を扱うのは得意でないのでyoutube.com のような名前でアドレスを表すことが多い．
      これを\term[]{ドメインネーム}という．
      ドメインネームをIPアドレスに変換するサービスを行うのが
      \term[]{DNS}(Domain Name System) である．
      実際に仕事を行っているサーバーをDNS Server という．

      \begin{figure}[htb]
         \centering
         \includegraphics{~/Documents/GitHub/introInfo/ch4/images/figure6.5}
         \caption{ルーティング}\label{}
      \end{figure}
      
      ルーター・ゲートウェイを通じてパケットは受け渡されてゆき
      最終的なターゲットのコンピューターに渡される．
      その順番が\term[]{ルーティング}である．
      ルーターやゲートウェイは頻繁に
      他のルーターやゲートウェイを通るとどれ位時間が掛かるかのリスト
      のようなものを常に更新している．
      それを元に次に何処に送れば最適かを判断して
      パケットを転送している．
   \item 
      IPパケットは
      \begin{figure}[htb]
         \centering
         \includegraphics{~/Documents/GitHub/introInfo/ch4/images/figure6.3}
         \caption{IPのパケットの構造}\label{}
      \end{figure}
      
      
      \noindent
      のようなものである．基本情報にはversion とかヘッダや全体の長さの情報が入っている．
      TTL はTime to live(生存時間)の略で$1$byteの大きさを持ち，
      初期条件が$40$程度に設定されていて
      gatewayを通過する度に$1$ずつ減っていき
      $0$になるとそのパケットは廃棄される．

   \item 
      \term[]{TCP}(Transition Control Protocol)は
      図~\ref{levels}に示したようにIPの上に存在するプロトコルである．
      各コンピュータ毎に \term[]{ポート}(Port)と呼ばれる$2$byte, 
      つまり$=2^{16}\approx  64,000$程の数字({port number})を
      指定できる．
      Portは日本語だと港\footnote{語学に興味のある人へ：
      port は港から派生して昔の戦艦の砲門のふたことを表すようになった．
      それから工学では外部器機やケーブルを接続するための
      穴のことを意味するようになった．}
      のことだが，
      会社の電話番号がIP アドレスでPortは会社の内線番号のような役割を
      はたしている．
      つまり一台のサーバーがこれを使って6.4万程度のコンピューターと
      同時に接続可能である．
      
         %socket  
   
      通信には送り手のIPアドレスとポートナンバーと
      受取り手のIPアドレスとポートナンバーの組が必要になる
      ことに注意してほしい．
      また，送りたいパケットが長いときは適宜分割して，
      その分割の何番目かを示すシークエンス番号が付いている．
      \begin{figure}[htb]
         \centering
         \includegraphics{~/Documents/GitHub/introInfo/ch4/images/figure6.4}
         \caption{TCPのパケットの構造}\label{}
      \end{figure}
      
      \term[]{UDP}(User Datagram Protocol)というTCPと同じレベルの
      プロトコルもあり，双方向のストリームを必要としないデータ交換に使われている．
      

   \item 
      図~\ref{levels}に示したようにTCP/UDPの上にのるのが
      \term[]{アプリケーション層}(Application layer)である．
      FTP, telnet, ssh, SMTP, POP3, HTTP, HTTPS などのようにインターネットで使うアプリケーションに対応したプロトコルである．
      例えばメールだとSMTP, POP3のようなプロトコルを使っている．
      FTP はデータを転送するためのプロトコルであり，
      HTTPはウェブを実現するプロトコルである(HTTPについては次回詳しく解説する)．
       
      特に最初の通信を確立するためにTCP/UDPのポートの規定値として使うのが
      \term[]{Wellknown port} である．その表を以下に挙げる：
      \begin{table}[h]
         \centering
         \begin{tabular}{lll} 
              \noalign{\hrule height 1pt}
                  Port No.&& プロトコル\\ \hline\hline
                  20 &TCP&FTP\\ \hline
                  22 & TCP/UDP&ssh \\ \hline
                  23 & TCP&telnet \\ \hline
                  25 &TCP/UDP&SMTP\\ \hline
                  52 &TCP/UDP& DNS\\ \hline
                  80 &TCP/UDP&HTTP\\ \hline
                  110 &TCP&POP3\\ \hline
                  443 &TCP/UDP&HTTPS\\
             \noalign{\hrule height 1pt}
         \end{tabular}
         \smallemptyline
         \caption{Wellknown port の例}
      \end{table}
      
      通常あるプロトコルの通信が確保されるとwell-known port以外の適当な
      ポートを指定し直して，それを使って以後の通信が行われる．
\end{itemize}

\subsection{Utilities}
UNIXの端末にはTCP/IP用の様々なツールが用意されている：
\begin{itemize}
   \item 
      netstat 
   \item    
      ifconfig
   \item 
      ping
   \item 
      traceroute 
    \item 
      port scanner: 標準ではインストールされていないがNmap などが使える．
\end{itemize}


%\begin{thebibliography}{99}
%   \bibitem{K}  
%      カーニハン, 
%      教養としてのコンピューターサイエンス講義 第二版, 日経BP, 2020.
%   \bibitem{T}  
%     井上 直也, 村山 公保, 
%     マスタリングTCP/IP—入門編—(第6版), オーム社, 2019.
%   \bibitem{TS}  
%      齋藤孝道, 
%      マスタリングTCP/IP 情報セキュリティ編 (第2版), オーム社, 2022.
%
%   \bibitem{TR}  
%      山川秀人, 水谷昭博, 
%      マスタリングTCP/IP ルーティング編, オーム社, 2007.
%
%   \bibitem{UNIX}
%      カーニハン,
%      カーニハンのUNIX回顧録, 丸善, 2025.
%\end{thebibliography}
%\end{document}

\section{Webの仕組み}
\subsection{Webの始まり}
現在インターネットで最も使われているのが
\term[]{WWW}(World Wide Web), 略してW3 とか Web と呼ばれるプロトコルである．
1980年にCERN (Centre Europ\' een de Recherche Nucl\' eaire)のティム・バーナーズ=リー(\cite{LT})が個人的な開発を始め， 
1989年に「Information Management: A Proposal」という提案書をまとめた．そこに書かれているのがWWWである．

WebはもともとURL・HTTP・HTMLというプロトコルの三つ組みから成っていた．

\begin{itemize}
   \item
      \term[]{URL}は Uniquely Resource Locater の略で
      情報の位置を示すもので基本構造は以下のようである：
      
       \href{https://info.cern.ch/hypertext/WWW/TheProject.html}{https://info.cern.ch/hypertext/WWW/TheProject.html}  
       
       最初のhttps が\term[]{スキーム}という接続方式を示す所で，
       次のinfo.cern.ch\footnote{ch はSwitzerlandのドメイン名}がサーバー名，残りがパスである．
       因みに上のURLはCERNに公表されている
       最初のWWWのページのアーカイブのURL(へのリンク)である．
       
   \item
      \term[]{HTTP}は Hyper Text Transfer Protocol の略で
      WWWの情報の受け渡し方法を決めている．
      後でもう少し詳しく述べる．
      スキームは(この時点ではhttpsは存在せず)
      httpのほかftp, fileなども柔軟に選ぶことが出来た．
   
   \item
      \term[]{HTML}は Hyper Text Markup Language の略で
      情報の表現方法を決めている．
      SGMLという既に存在したマークアップ言語
      (mark up とは印刷の用語で，原稿の余白に
      印刷に必要なフォントやその大きさ，割り付けなどを書き込むこと)
      を拡張したページの記述のための言語である．
      \term[]{ハイパーテキスト}(Hyper Text)の実装を含んでいた．
      
      Hyper Text自体の公表はテッドネルソンのザナドゥなどが
      先であったがインターネットでの実装はHTMLのほうが早かった．      
\end{itemize}


    \begin{figure}[htb]
	\centering
	\includegraphics{~/Documents/GitHub/introInfo/ch4/images/figure7.1}
	\caption{Client Server}\label{}
   \end{figure}


WWWの実装にはインターネット上に情報を持っている
\term[]{サーバー}(Server)と
それをローカルで表示する\term[]{クライアント}(Client)
としての\term[]{ブラウザー}(Browser)
が必要だったが，
ティム・バーナーズ=リーはそのようなサーバー(CERN httpd)と
ブラウザー(WorldWideWeb)も公表した．
また初期でも画像もHyper Textを用いて表示することが出来た．
最後に，彼はWebの著作権を全て放棄しているが，
それもWWWの成功の一因であると言われている．


\subsection{Webの発展}
手元にある短いindex.htmlというhtml のソースファイルを見てみる：
\begin{verbatim}
   <!DOCTYPE html>
   <html>
     <head>
       <meta charset="utf-8">
       <meta name="viewport" width=device-width, 
          initial-scale=1.0, maximum-scale=1.0, user-scalable=0>
       <style> body {padding: 0; margin: 0;} </style>
       <script src="p5.js"></script>
       <!-- <script src="p5.sound.min.js"></script> -->
       <script src="sketch.js"></script>
     </head>
     <body>
       <a href="https://p5js.org/">p5.js</a>
     </body>
   </html>
\end{verbatim}

基本\verb/<hoge>/ と\verb/<\/hoge\verb/>/で要素を囲む形で
構造を示しているのが分かる．このようなものを\term[]{タグ}(tag)という．
最初の行がこのファイルがhtmlであることを示し，
\verb/</html\verb/>/と\verb/<\/html\verb/>/で囲まれた部分がhtmlである．
それが\verb/<head>/と\verb/<\/head\verb/>/に囲まれたヘッダ部分と，
\verb/<body>/と\verb/<\/body\verb/>/に囲まれた本文に分かれている．
ヘッダ部分には文字コード，スタイルや\verb/<script>/と\verb/<\/script\verb/>/に囲まれた部分にJavaScriptの(そのディレクトリーに置かれた)ファイルが指定されている．
また本文の\verb/<a>/と\verb/<\/a\verb/>/には(hrefの後に)リンクが指定されている，
ただし，a は anchor の略である．
ツリー構造が分かるように適当にインデントが入れてあることも分かると思う．
また，index.htmlというのはそのディレクトリにアクセスするとき最初に読み込まれるファイルの標準的な名前である．
\emptyline

\begin{Q}
   以前紹介したmarkdown のファイル IntroToMarkdown.md のある
   ディレクトリーで
\begin{verbatim}
   % pandoc -f markdown -t html -s -o IntroMD.html 
         IntroToMarkdown.md
   % open IntroMD.html
\end{verbatim}
   とするとhtmlに変換されたページがブラウザーに表示されるので，
   確かめてください
   (300行ほどのファイルでそのうち200行以上が後で紹介するCSSの部分である)．
   ただし，\verb/% /はプロンプト．
\end{Q}

\begin{Q}
   上のCERNへのリンクをクリックして最初のWWWのページを訪れてください．
   ChormeだとViewのメニュー中のDeveloperのサブメニューのなかのView source
   を選ぶとソースファイルを見ることが出来るのでやってみてください．
   他のサブメニューのアイテムも選んでみて何が出来るかためしてみてください．
\end{Q}
    \begin{figure}[htb]
	\centering
	\includegraphics[width=12cm]{~/Documents/GitHub/introInfo/ch4/images/Screenshot4.png}
	\caption{世界初のWWWのホームページのソース}\label{}
   \end{figure}

\begin{Q}
   一般のサイト, 例えば\href{https://www.yahoo.co.jp/}{yahoo}の
   ホームページのソースファイルを上と同じ方法で確認してください．
   見たら分かると思いますが，最近のWebサイトのソースファイルは
   サーバーのコンピュータが自動生成したものが多くなっています．
\end{Q}

WWWは広く受け入れられていったが，直ぐに様々な拡張が成されていった．
\begin{itemize}
   \item
      初期の頃は
      HTMLから見た目に関する部分を分離したのが
      \term[]{CSS}(Cascaded Style Sheet) である．
      美術的なセンスが必要なところなので難しい所もあるが，
      簡易的に使えるものも色々出回っている．
   
   \item
      \term[]{HTTPS} は暗号化などHTTPに
      セキュリティーの強化を図ったプロトコルである．
   \item
    \begin{figure}[htb]
	\centering
	\includegraphics[width=12cm]{~/Documents/GitHub/introInfo/ch4/images/Screenshot2.png}
	\caption{クエリと返答}\label{query}
   \end{figure}
   
      \term[]{クエリ}(query)：amazonで「国債」を検索すると図~\ref{query}
      のような画面になる．
      このときブラウザーのアドレスバーに表示されているURLを抜き出すと
\begin{verbatim}
   https://www.amazon.co.jp/s?k=
   %E5%9B%BD%E5%82%B5&crid=2WZVBYMGFR6YT&
   sprefix=%E5%9B%BD%E5%82%B5%2Caps%2C179&
   ref=nb_sb_noss_1
\end{verbatim}

\noindent この？以下に表れる部分がクエリである．
このようなURLを使って質問を送り，サーバーが質問に対する答えのホームページを返している．
\verb/%/ 
で始まる16進数二つの並びがUFT-8であり，\term[]{パーセントエンコーディング}という．
確認してみると

\begin{verbatim}
   % echo -n "国債"|od -An -tx1
   e5  9b  bd  e5  82  b5                                        
\end{verbatim}
で$=$以下の部分と一致していることがわかる．
同様の方法で様々なデータのやりとりを行っている．

    \item
       \term[]{クッキー}(Cookie): Webのログイン情報や使用に関する情報を
       Web Browserの中に保存したもの．
       ChromeだとViewのメニュー中のDeveloperのサブメニューのなかの
       Developer toolsを選び，Application のタブのペインの中の
       Storage のCookieを選ぶと保存されたCookieを確認出来る．
    \begin{figure}[htb]
	\centering
	\includegraphics[width=12cm]{~/Documents/GitHub/introInfo/ch4/images/Screenshot3.png}
	\caption{https://www.google.com/のCookie}\label{Cookie}
   \end{figure}
   
      あとでセキュリティーの部分でもコメントするが，
      Cookieには重要な情報も保存されていて，
      それがネット上で売買されているので
      定期的に消去することをお奨めする．

   \item
      HTMLに動きを与える目的で作られたプログラミング言語が 
      \term[]{JavaScript} である．最初はLiveScriptと呼ばれたが，
      当時注目されていたJavaに関連してJavaScriptと名前を替えた．
      ただし，Javaとは直接関係のない言語である．

    \item    
     文書の構造をツリーで表したものをDOM(Document Object Model)という．
     図~\ref{dom}に後に使うhello.html のDOMを示す．
    これを使って JavaScriptは様々なページの内容にアクセスする
    ことが出来る．

    \begin{figure}[htb]
	\centering
	\includegraphics{~/Documents/GitHub/introInfo/ch4/images/figure7.2}
	\caption{hello.html のDOM}\label{dom}
   \end{figure}

 
      \begin{Q}
         クエリの実験を自分のパソコンで再現してください．
      \end{Q}
      \begin{Q}
         Cookieの実験を自分のパソコンで再現してください．
      \end{Q}

\end{itemize}


\subsection{Webのサーバー・クライアントの演習}
ここではWEBのクライアント・サーバーのコマンドラインでの演習を行う:
特に nc  = net cat と
curl = client for URL を使用して，WEBのクライアント・サーバーを体験してもらいたい(\cite{Ko}参)．
\begin{itemize}
   \item
      まず，サーバーを立ち上げる：作業ディレクトリーに移動しサーバー側の端末で
\begin{verbatim}
   % nc -l 8080 #return #Ubntuでは >nc -q 1 -l 8080
\end{verbatim}
      とする．8080がTCPのポートナンバーである．

      クライアント側の別の端末を立ち上げ
\begin{verbatim}
   % curl http://localhost:8080/ #return
\end{verbatim}
      とすると，サーバ側の端末に次のような表示が表れる
\begin{verbatim}
   % nc -l 8080 
   GET / HTTP/1.1
   Host: localhost:8080
   User-Agent: curl/8.1.2
   Accept: */*

\end{verbatim}
      そこで，サーバー側で手で
\begin{verbatim}
   HTTP/1.1 200 OK #return
   #return
   Hello, Web application! #return
\end{verbatim}
      と入力すると，クライアント側で
\begin{verbatim}
   % curl http://localhost:8080/ 
   Hello, Web application!
\end{verbatim}
      と表示されるはずである．もし，HTTP/1.1 200 OKも表示したいときは
\begin{verbatim}
   % curl  -v  http://localhost:8080/ 
\end{verbatim}
とすればいい．やってみると以下のような出力が得られた：
\begin{verbatim} 
   *   Trying 127.0.0.1:8080...
   * Connected to localhost (127.0.0.1) port 8080 (#0)
   > GET / HTTP/1.1
   > Host: localhost:8080
   > User-Agent: curl/8.1.2
   > Accept: */*
   > 
   < HTTP/1.1 200 OK
   * no chunk, no close, no size. Assume close to signal end
   < 
   Hello, Web application!
   * Closing connection 0
\end{verbatim}

   \item
      毎回手入力するのは大変なので，1回ctrol Dでncを終了して，
      もう一度
\begin{verbatim}
    % {echo 'HTTP/1.1 200 OK';echo ; 
    echo 'Hello, Web application!';}|nc -l 8080 
    #echo の前の改行不要
 \end{verbatim}
       とすると，自動で先程の応答を送ることが出来る．
 
    \item
      次に返事するファイルをhtmlに変更する．まず
      作業ディレクトリーがtmp のとき，その中に server/ と client/
      ディレクトリーをつくる     
\begin{verbatim}
   % mkdir server client
\end{verbatim}
      server/に以下のような hello.html を作成する
\begin{verbatim}
   <!DOCTYPE html>
   <html lang="ja">
       <head>
           <meta charset="UTF-8">
           <title>Test Page</title>
       </head>
       <body>
           <h1>Test Contents</h1>
           <p id="greed">Hello, Web application!</p>
       </body>
   </html>
\end{verbatim}
      サーバー側で
\begin{verbatim}
    % {echo 'HTTP/1.1 200 OK';echo ;cat hello.html} |
    nc -l 8080  #|nc の前の改行不要
\end{verbatim}
      クライアント側で
\begin{verbatim}
    % curl -s -o hello.html http://localhost:8080/
    % open hello.html
\end{verbatim}
      とすると
      図~\ref{Browser} のようにブラウザーに表示されたであろう．
    \begin{figure}[htb]
	\centering
	\includegraphics[width=12cm]{~/Documents/GitHub/introInfo/ch4/images/Screenshot1.png}
	\caption{TestのBrowserでの表示}\label{Browser}
   \end{figure}

   \item
      最後にブラウザーで直接表示するためには，
      (nc が既に終了しているので)サーバー側で同じことをして
\begin{verbatim}
% !!
\end{verbatim}
      ブラウザーのアドレスバーにhttp://localhost:8080/と入れると
      同じ結果が表示されるであろう．
      このときサーバー側には以下のようなログが表示される．
\begin{verbatim}
   GET / HTTP/1.1
   Host: localhost:8080
   Connection: keep-alive
   sec-ch-ua: "Chromium";v="148", "Google Chrome";v="148",
    "Not/A)Brand";v="99"
   sec-ch-ua-mobile: ?0
   sec-ch-ua-platform: "macOS"
   Upgrade-Insecure-Requests: 1
   User-Agent: Mozilla/5.0 (Macintosh; Intel Mac OS X 10_15_7) 
   AppleWebKit/537.36 (KHTML, like Gecko) Chrome/
   148.0.0.0 Safari/537.36
   Accept: text/html,application/xhtml+xml,application/xml;
   q=0.9,image/avif,image/webp,image/apng,*/*;q=0.8,
   application/
   signed-exchange;v=b3;q=0.7
   Sec-Fetch-Site: none
   Sec-Fetch-Mode: navigate
   Sec-Fetch-User: ?1
   Sec-Fetch-Dest: document
   Accept-Encoding: gzip, deflate, br, zstd
   Accept-Language: en-US,en;q=0.9

\end{verbatim}
 
       \noindent
       これらがHTTPプロトコルの応答になる．
    
      \begin{Q}
         この実験を自分のパソコンで再現してください．
         最後の実験でブラウザーでソースファイルを開き，
         hello.html と同じものであることを確認してください．
      \end{Q}

\end{itemize}


%
%JavaScript については解説する時間も無いので
%独立して読める%\href{https://qu365-my.sharepoint.com/:b:/g/personal/otsu_yukio_526_m_kyushu-u_ac_jp/IQAzL8zjA-ZRR6pcO1oq3UeKAd2SKBpVYA7TWGz80IFkqvg?e=NhinY5}
%{解説}をMoodleのホームページに上げておくので
%興味のある人は読んで欲しい．
%私のポートフォリオにもJavaScript で買いたアプリを公開している．
%%\section{one page application(OPA)}



%\begin{thebibliography}{99}
%   \bibitem{L} 
%      ティム・バーナーズ=リー, Webの創成 — 
%      World Wide Webはいかにして生まれどこに向かうのか, 
%      毎日コミュニケーションズ, 2001.
%   \bibitem{K}
%      小森 裕介,［ 改訂新版］プロになるためのWeb技術入門, 
%      技術評論社, 2024.
%\end{thebibliography}



%\documentclass[a4paper,12pt]{amsart}
%\usepackage[dvipdfmx]{graphicx,xcolor}%graphics
%\usepackage{float}
%\usepackage{amsmath}
%\usepackage{verbatim}
%\usepackage{latexsym}
%\usepackage{amssymb}
%%\usepackage{emath}
%\usepackage{listings}
%\usepackage{ccicons}
%\usepackage{mflogo}
%%\usepackage{fancybox}%enclose by box
%\usepackage[dvipdfmx]{hyperref}
%\usepackage{verbatim}
%\usepackage{ulem}
%
%\hypersetup{% hyperrefオプションリスト
%   setpagesize=false,
%   bookmarksnumbered=true,%
%   bookmarksopen=true,%
%   colorlinks=true,%
%   linkcolor=red,
%   citecolor=blue
%} 
%
%\newtheorem*{Thm}{定理}
%\newtheorem*{Prop}{命題}
%
%\newtheorem*{exQ}{レポート}
%
%\theoremstyle{remark}
%\newtheorem*{ex}{例}
%\newtheorem*{Q}{演習}
%\newtheorem*{Quiz}{問題}
%
%\renewcommand{\figurename}{図}
%\renewcommand{\proofname}{証明}
%\renewcommand{\tablename}{表}
%\setlength{\belowcaptionskip}{0mm}
%
%\DeclareMathOperator{\Ker}{Ker}
%\DeclareMathOperator{\Image}{Im}
%
%\def\emptyline{\vspace{11pt}}
%\def\smallemptyline{\vspace{6pt}}
%\begin{document}
%\begin{flushright}
%   2026/08/20
%   %2026/06/30
%   %2026/03/11   
%   
%   九大・数理 大津幸男
%\end{flushright}
%
%
%\begin{center}
%   {\Large 情報科指導法II---09}
%   \href{https://creativecommons.org/licenses/by-nc-nd/4.0/}{\ccbyncnd}
%   \par{セキュリティー1}
%\end{center}

\chapter{セキュリティー}
セキュリティーには数学的・技術的側面と人間・社会・経営的側面があるが，
まず数学的・技術的側面をとりあげ，その後人間・社会・経営的側面を取り上げる．

\section{暗号の数理}
\subsection{$\mathbf Z$の剰余類}
まず，暗号や認証で使われる初等整数論の復習を行う．
どうしても代数でよく使われる言葉が出てきてしまうので
\cite{Ho}, \cite{Ta}, \cite{Ko} などを参考にしながら
読み解いもらいたい．
この解説が代数へのイントロとなれば望外の喜びである．

$\mathbf N$ を（$1$以上の）自然数全体のなす集合, $\mathbf Z$を整数全体のなす集合とする．$\mathbf Z$は加法$+$
に関して群\footnote{集合$G$が\term[]{群}であるとは任意の$a$, $b\in G$に対して積$ab\in G$が定まり，次の三条件を満たすことである：
\[
   (ab)c = a(bc)\quad \forall a, b, c\in G
\]
\term[]{単位元}$e\in G$が存在し，
\[
   ea=ae=a\quad \forall a\in G
\]
任意の$a\in G$に対して\term[]{逆元}$a^{-1}$が存在して
\[
   aa^{-1}= a^{-1}a = e
\]
$ab=ba$をみたすとき\term[]{可換群}といい，$+$で表すことが多い．}
になっている．
$\mathbf Z$には掛け算$\cdot$ ( $\times$と書くこともある)もあるが$0$を除いた$\mathbf Z\setminus\{0\}$では割り算ができない(任意の$n\neq\pm1$は$\mathbf Z$内で逆元が無い)ので群ではない．
$\mathbf Z$は$+$と$\cdot$ に関して環\footnote{集合$R$が\term[]{環}であるとは演算$+$と$\cdot$ をもち$+$に関しては可換群で，$\cdot$ については結合律をみたし可換であり，$\cdot$と$+$については分配則をみたすこと．$\mathbf Z$や多項式環$\mathbf Z[X]$が典型例}であるが，体\footnote{環$K$が体とは$K$の$0$以外の元が$\cdot$についての逆元を持つこと．$\mathbf R$や$\mathbf C$が例．}ではない．
$m+n=n+m$なので可換である，%加法$+$に関して加群ということがある．
$m\in \mathbf N$を固定し，$x, x'\in \mathbf Z $
\[
   x\equiv x' \mod m \Leftrightarrow 
   \exists q \in  \mathbf Z \text{ such that }x-x'=qm
\]
と置くとこれは同値関係である．
$x\in \mathbf Z $を代表元とする同値類を$\bar x$と置く．
%$\mathbf Z $を加法$+$に関する可換群と見なすと
\[
\begin{split}
   &x\equiv x' \mod m  \text{ and } y\equiv y' \mod m\\
   \Rightarrow \quad
   & x + y \equiv x' + y' \mod m,\quad
   xy\equiv x'y' \mod m
\end{split}
\]
なので
\[
   \bar x + \bar y = \overline{x+y},\quad
   \bar x\cdot \bar y = \overline{xy}
\]
と定め暫く$G$と置く．
$(G, +)$ は$\bar 0$を単位元とする可換群である．
\[
   f:\mathbf Z \ni x \mapsto \bar x \in G
\]
とすると$f(x+y)=f(x)+f(y)$\footnote{このようなものを$+$についての群の準同形という．}であり，
\[
   \{x\in \mathbf Z\mid  f(x)=\bar 0\}= m\mathbf Z
\]
見なせるので
%$ m\mathbf Z$は$\mathbf Z$の正規部分群（$\forall a\in \mathbf Z$について$ a+m\mathbf Z-a=m\mathbf Z$）であるので，群の準同形定理より
\[
   G \sim \mathbf Z/m\mathbf Z
\]
となる．
\term[]{巡回群}とは一つの元から生成される群のことであり，
$\mathbf Z/m\mathbf Z$は$+$については$\bar 1$で生成されているので
有限巡回群である．これを$C_m$と表すことがある．
%(また，無限集合になる巡回群を無限巡回群といい，$\mathbf Z$と同形になる．)
\emptyline

$\mathbf Z/m\mathbf Z$は$+$ と$\cdot$について$\mathbf Z$の計算と同じルールで計算していい．例えば
\[
  \bar x (\bar y+\bar z)
  = \overline{xy+xz}
  = \bar x \bar y+\bar x \bar z
\]
になる\footnote{$f(x+y)=f(x)+f(y)$, $f(xy)=f(x)f(y)$をみたすものを環の準同形という．}．以下，
$\mathbf Z/m\mathbf Z$の積$\cdot$について詳しい構造を調べる．
最初にいくつか計算してみる．
\begin{ex}
   $\mathbf Z/4\mathbf Z$の足し算と掛け算の演算表を作ってみる．
   足し算については巡回群$C_4$なのでそのままであるが，
   積に関しては$\bar 2$が零因子\footnote{$a\neq 0$, $b\neq 0$ が$ab=0$となるとき零因子という．}となっているので群ではない．
%   従って$\mathbf Z/4\mathbf Z$は体でない．
   \begin{table}[htb]
      \centering
      \begin{minipage}{0.4\columnwidth}
         \centering
         \begin{tabular}{l|llll}\noalign{\hrule height 1pt} 
             $+$&$\bar 0$&$\bar1$&$\bar2$&$\bar3$ \\ \hline
            $\bar 0$&$\bar 0$&$\bar1$&$\bar2$&$\bar3$ \\
            $\bar1$& $\bar1$&$\bar2$ &$\bar3$&$\bar0$ \\
            $\bar2$ &$\bar2$&$\bar3$&$\bar 0$&$\bar1$ \\ 
            $\bar3$& $\bar3$ &$\bar0$& $\bar1$&$\bar2$\\ \noalign{\hrule height 1pt} 
         \end{tabular}
         \smallemptyline
         \caption{$+$ $\mod 4$}\label{mod4plus}
      \end{minipage} 
      \begin{minipage}{0.4\columnwidth}
         \centering
         \begin{tabular}{l|llll}\noalign{\hrule height 1pt} 
             $\cdot$&$\bar 0$&$\bar1$&$\bar2$&$\bar3$ \\ \hline
            $\bar 0$&$\bar 0$&$\bar0$&$\bar0$&$\bar0$ \\
            $\bar1$& $\bar0$&$\bar1$ &$\bar2$&$\bar3$ \\
            $\bar2$ &$\bar0$&$\bar2$&$\bar 0$&$\bar2$ \\ 
            $\bar3$& $\bar0$ &$\bar3$& $\bar2$&$\bar1$\\ \noalign{\hrule height 1pt} 
         \end{tabular}
         \smallemptyline
         \caption{$\cdot$ $\mod 4$}\label{mod4times}
      \end{minipage} 
   \end{table}
\end{ex}

\begin{ex}
   $\mathbf Z/5\mathbf Z$の足し算と掛け算の演算表を作ってみる．
   掛け算では零因子がないので$\bar0$を除く$(\mathbf Z/5\mathbf Z)^{\times}=(\mathbf Z/5\mathbf Z)\setminus \{\bar 0\}$は群となる．
   これは巡回群$C_4$と同型である(表~\ref{mod4plus}
   と異なって見えるかもしれないが，$\bar 1_4$, $\bar 2_5$のように表し区別し，$f:(\mathbf Z/5\mathbf Z)^{\times}\to \mathbf Z/4\mathbf Z$を
   \[
      f(\bar 1_5)=\bar 0_4,\quad
      f(\bar 2_5)=\bar 1_4,\quad
      f(\bar 3_5)=\bar 3_4,\quad
      f(\bar 4_5)=\bar 2_4
   \]
   とすると$(\mathbf Z/5\mathbf Z)^{\times}$は掛け算，$\mathbf Z/4\mathbf Z$は足し算として同型\footnote{$f$が(この場合は群の)準同形で，更に全単射になること}になることが分かる)．
%   このことから$\mathbf Z/5\mathbf Z$は体でもある．
   
   \begin{table}[htb]
      \centering
      \begin{minipage}{0.4\columnwidth}
         \centering
         \begin{tabular}{l|lllll} \noalign{\hrule height 1pt} 
             $+$&$\bar 0$&$\bar1$&$\bar2$&$\bar3$ &$\bar4$\\ \hline
            $\bar 0$&$\bar 0$&$\bar1$&$\bar2$&$\bar3$&$\bar4$ \\
            $\bar1$& $\bar1$&$\bar2$ &$\bar3$&$\bar4$& $\bar0$\\
            $\bar2$ &$\bar2$&$\bar3$&$\bar 4$&$\bar0$&$\bar1$ \\ 
            $\bar3$& $\bar3$ &$\bar4$& $\bar0$&$\bar1$&$\bar2$\\ 
            $\bar4$& $\bar4$& $\bar0$&$\bar1$&$\bar2$&$\bar2$ \\\noalign{\hrule height 1pt} 
         \end{tabular}
   \smallemptyline
         \caption{$+$ $\mod 5$}
      \end{minipage} 
      \begin{minipage}{0.4\columnwidth}
         \centering
         \begin{tabular}{l|lllll} \noalign{\hrule height 1pt} 
             $\cdot$&$\bar 0$&$\bar1$&$\bar2$&$\bar3$ &$\bar4$\\ 
            \hline
            $\bar 0$&$\bar 0$&$\bar0$&$\bar0$&$\bar0$&$\bar0$ \\
            $\bar1$& $\bar0$&$\bar1$ &$\bar2$&$\bar3$ &$\bar4$\\
            $\bar2$ &$\bar0$&$\bar2$&$\bar 4$&$\bar1$ &$\bar3$\\ 
            $\bar3$& $\bar0$ &$\bar3$& $\bar1$&$\bar4$&$\bar2$\\ 
            $\bar4$& $\bar0$& $\bar4$&$\bar3$&$\bar2$&$\bar1$ \\\noalign{\hrule height 1pt} 
         \end{tabular}
   \smallemptyline
         \caption{$\cdot$ $\mod 5$}
      \end{minipage} 
   \end{table}
\end{ex}

%上で指摘したように$\mathbf Z/m\mathbf Z$には$\cdot$という積があり環となっている．しかし，$\bar 0$を除いても$\cdot$に関して一般には群ではない．


$m, n\in \mathbf Z$に対して$(m,n)$で$m$ と $n$の最大公約数を表す．
\[
   (m,n)=1
   \quad \Leftrightarrow\quad
   \exists a, b\in \mathbf Z \text{ such that } am + bn =1
\]
がユークリッドの互除法から分かる．このことから
\[
   (m,n)=1\quad 
   \Leftrightarrow\quad
   \exists  \bar b\in \mathbf Z/m\mathbf Z \text{ such that } 
         \bar b\bar n= \bar 1
\]
$(m,n)=1$のとき$m$と$n$は互いに素という．
$m\in \mathbf N$を固定し$(m,n)=1$なる
$\bar n\in\mathbf Z/m\mathbf Z$を\term[]{既約剰余類}という．

既約剰余類$\bar a\in \mathbf Z/m\mathbf Z$ を固定すると$\bar b\bar a = \bar 1$なる$\bar b\in \mathbf Z/m\mathbf Z$が見つかるので
\[
   f:\mathbf Z/m\mathbf Z\ni \bar n 
    \mapsto \bar{a}\cdot\bar{n}\in  \mathbf Z/m\mathbf Z
\]
を考えると$f$は$+$に関して準同形写像で，$\bar b$を掛けることで同形写像であることが分かる．
\[
   \varphi(m)= |\{n\mid n=1,\dots, m-1 \text{ such that } (m,n)=1\}|
\]
と置くと$\mathbf Z/m\mathbf Z$には$\varphi(m)$個の既約剰余類が存在することが分かる．
$\varphi$は\term[]{オイラー関数}と呼ばれる．
$\mathbf Z/m\mathbf Z$には$\varphi(m)$個の生成元が存在し，
$\mathbf Z/m\mathbf Z$はどの生成元をとっても加群としては
同型な巡回群となる．

\begin{ex}
   $\varphi(4)=2$, $\varphi(5)=4$.
\end{ex}

\begin{Q}
   $p$を素数とするとき
   \[
      \varphi(p^n)=p^n-p^{n-1}
   \]
   を示せ．
\end{Q}



$(m,n)=1$なる$n$を固定する．
$\bar x_1, \cdots, \bar x_{\varphi(m)}$を$\mathbf Z/m\mathbf Z$の既約剰余類全体とする．$\bar n \bar x_1, \cdots , \bar n\bar x_{\varphi(m)}$も既約剰余類全体なので，集合としては一致している．したがって
\[
   \bar x_1\cdots \bar x_{\varphi(m)}=\bar n\bar  x_1 \cdots \bar n \bar x_{\varphi(m)}
   =\bar n^{\varphi(m)}\bar x_1\cdots \bar x_{\varphi(m)}
\]
$\bar x_1\cdots \bar x_{\varphi(m)}$は逆元をもつので
$
   \bar n^{\varphi(m)}=\bar 1\in\mathbf Z/m\mathbf Z
$,  つまり$(m,n)=1$なる任意の$n$に対して
\begin{equation}\label{varphi}
   n^{\varphi(m)}\equiv 1 \mod m
\end{equation}
$p$を素数とすると，$\varphi(p)=p-1$なので任意の$x\not\equiv 0 \mod p$に対して
\[
   x^{p-1}\equiv 1 \mod p
\]
となり，\term[]{フェルマーの小定理}が得られる．これから任意の$x\in \mathbf Z$に対して
\begin{equation}\label{funPoly}
   x^{p}\equiv x \mod p
\end{equation}
も分かる．

\begin{Prop}[中国の剰余定理]
   互いに素な$m$, $m'\in \mathbf N$を取る．任意の$a$, $a'\in\mathbf Z$に対して
   \[
      x\equiv a \mod m,\quad x\equiv a' \mod m'
   \]
   を満たす$x$が存在して$\mod mm'$で一意．
\end{Prop}
\begin{proof}
   $m$, $m'\in \mathbf N$は互いに素なので
   \[
      nm+n'm'=1
   \]
   なる$n$, $n'\in\mathbf Z$が存在する．つまり
   \[
      n'm'\equiv 1 \mod m,\quad nm\equiv 1 \mod m'
   \]
   なので，$x= an'm'+a'nm\in\mathbf Z$の$\mod mm'$の代表が求めるものである．
   
   $x$, $x'$が方程式を満たすと
   $x-x'$は$\mod m$でも$\mod m'$でも$0$
   よって$x-x'$は$\mod mm'$でも０になり，一意性が分かる．
\end{proof}
\begin{Prop}
   互いに素な$m$, $m'\in \mathbf N$に対して
   \begin{equation}\label{eq3}
      \varphi(mm')=\varphi(m)\varphi(m')
   \end{equation}
\end{Prop}
\begin{proof}
$k$ を$1, \cdots, mm'-1$の間にある自然数で$mm'$と互いに素なものとする．
すると
\[
   k\equiv  j\mod m, \quad k\equiv  j'\mod m'
\]
なる$j\in \{1, \dots m-1\}$, $j'\in \{1, \dots m'-1\}$が見つかる．
逆に中国の剰余定理は$j$, $j'$から$k$が$\mod mm'$で
一意的に定まることを意味する．
$k$は$mm'$と互いに素なので，
ある$a$, $b\in\mathbf Z$が存在して
\[
   ak+b m m'=1
\]
よって$k = j + c m$($c\in\mathbf Z$)と置くと
\[
   1 = a(j+c m)+b mm'=aj +(ac+bm')m
\]
となり $j$は$m$と互いに素になる．
同様に$j'$は$m'$と互いに素で，
これから示したい等式が分かる．
\end{proof}

\begin{ex}
$m=3$, $m'=5$のとき$\varphi(15)=\varphi(3)\varphi(5)=2\cdot 4=8$

 \begin{table}[htb]
      \centering
         \centering
         \begin{tabular}{l|lllllllllllllll}\noalign{\hrule height 1pt} 
            $\mod 15$&\sout{$0$}&$1$&$2$&\sout{$3$}&$4$&\sout{$5$}&\sout{$6$}
            &$7$&$8$&\sout{$9$}&\sout{$10$}&$11$ &\sout{$12$}&$13$&$14$\\ \hline
            $\mod 3$&$0$&$1$&$2$&$0$ &$1$&$2$&$0$&$1$&$2$&$0$&$1$&$2$ &$0$&$1$&$2$ \\
            $\mod 5$& $0$&$1$ &$2$&$3$&$4$&$0$&$1$&$2$&$3$&$4$ &$0$&$1$&$2$&$3$&$4$ \\
\noalign{\hrule height 1pt} 
         \end{tabular}
         \smallemptyline
         \caption{$\times$ $\mod 15$,  $\mod 3$,  $\mod 5$}\label{}
\end{table}
\end{ex}


ここではこれ以上必要ないが，将来的には楕円曲線暗号などで必要になるであろう有限体について簡単にまとめておく．
フェルマーの小定理は$(\mathbf Z/p\mathbf Z)^{\times}=(\mathbf Z/p\mathbf Z)\setminus \{\bar 0\}$が積$\cdot$について巡回群であることを意味している．
特に$\mathbf Z/p\mathbf Z$が体であることも意味していている．
このとき$\mathbf F_p$と表すことが多い．これは標数$p$の有限体である．
一般に$F$を体とする．環の準同形
$f:\mathbf Z\ni n\mapsto n\times 1\in F$が単射で無いとき，$\mathbf Z$ は単項イデアル整域なので$\Ker f=(p)$なる素数$p$がある．
これを標数という．
$f$が単射のときは\term[]{標数}は$0$と言う．
有限個の元からなる体が有限体で，定義からその標数はある素数$p$となる．
このとき$f$より$\bar f:\mathbf F_p\to F$は体の埋め込みである．
つまり，$F$は$\mathbf F_p$の拡大体となる．
このとき$\mathbf F_p$を\term[]{素体}という．
体の定義より$F$ は$\mathbf F_p$上のベクトル空間と見なせるが，
仮定から有限次元ベクトル空間であり，
この次元を$n$と置くと$F$は$q=p^n$の元からなることが分かる．
$\mathbf F_p$上のベクトル空間$F$ の$\mathbf F_p$上の同型群が
この代数拡大（は$\mathbf F_p$の正規分離拡大なので$F/\mathbf F_p$）のガロワ群となる．

\begin{Q}
   $n$を$2$以上の自然数とすると
   $\mathbf Z/p^n\mathbf Z$は$p^n$個の元からなる集合だが，
   体でないこと（したがって，これは
   上の有限体$F$とはなり得ないこと）を示せ．
\end{Q}


\begin{exQ}
   標数$p$の有限体は$\mathbf F_p$のガロワ拡大であることを示し，
   そのガロワ群を決定せよ（ヒント$:$
    \eqref{funPoly}の$p$を$q=p^n$に代えた多項式でこの拡大を特徴付けよ．またそのガロワ群はフロベニウス写像で生成されることを示せ）．
\end{exQ}


\subsection{公開鍵暗号---RSA}
アリス(Alice, 後でAと略す)とボブ(Bob, 後でBと略す)が共通の秘密の鍵を共有していて
その鍵を使ってアリスが文章(平文)を暗号化した暗号文を，
インターネット等を通じてボブに送り
ボブがその鍵を使って復号するような暗号を
\term[]{秘密鍵暗号}（共有鍵方式と呼ばれることの方が多い）という．
シーザー式暗号はその例である．送りたい文のA〜Z \footnote{この話を初めて聞いたとき小文字のa〜zや句読点や空白がないのに違和感を感じたが，この時代は実際そうであったようで小文字や句読法が発明されたのは中世後期であったそうである．もし今でも英語の文が全て大文字で語の間に空白や句読点がないものであったら英語を読むのはいかに困難であったであろうか．}
を順に0から25に変換し，$\mod 26$で$+3$した結果の0〜25をA〜Zに変換した文を送り，相手は逆算して暗号を特く．$3$が共有鍵（秘密鍵）になる．詳しい説明は控えるが秘密鍵暗号には以下のような方式がある：

  \begin{table}[h]
      \centering
      \begin{tabular}{llll} 
           \noalign{\hrule height 1pt}
               暗号& 鍵の長さ & 発表年& コメント \\ \hline \hline
               DES &54bit& 1977&危殆化 \\ \hline
               トリプルDES&$3\times 54$bit &&危殆化\\ \hline
               AES & 128/192/256bit&   2000&\cite{Y}に詳しい説明\\ \hline
           \noalign{\hrule height 1pt}
      \end{tabular}
    \smallemptyline     
      \caption{秘密鍵暗号}
   \end{table}


二人が知り合いで二人で合って秘密の鍵を共有することは簡単なように思えるが，
インターネットでこれまで会ったことのない二人，
例えばamazon とその新しい顧客が秘密の鍵を共有するには
どのようにすれば良いだろうか？
このような問題を解決するために生まれたのが\term[]{公開鍵暗号}である．

1976年デフィーとホフマンは公開鍵暗号のアイデアを発表したが，
1978年にリベスト・シャミール・アドルマンが発表した\term[]{RSA}アルゴリズムを解説したい：
Bob にAliceが秘密情報を送りたいとする．Bob は
同じぐらいの大きさのランダムな素数$p$, $q$を選び$m=pq$を計算する．\eqref{eq3}より
\[
   \varphi(m)=\varphi(p)\varphi(q)=(p-1)(q-1)=pq -p -q +1
\]
である．$pq -p -q +1$と互いに素な$0<e<pq -p -q +1$をランダムに選び，$\mod \varphi(m)$で$e$の逆元$0<f<\varphi(m)$ を互除法で計算する．つまり，
\begin{equation}\label{eq4}
   ef \equiv 1 \mod \varphi(m)
\end{equation}
Bob は$[m, e]$を公開鍵として公開し，
$[m, f]$を秘匿鍵として自分の所に秘密に保存しておく．
Alice はBob の公開鍵$[m, e]$を使って，
自然数の秘密情報$0\le C <m$から
\[
   C^e \equiv  D \mod m \quad 0\le D<m
\]
を計算しBobに$D$を送る．Bobは秘匿鍵$f$を使って
\[
   D^f \equiv C\mod m \quad 0\le C <m
\]
を得ることができる．実際，ある$k\in \mathbf Z$が存在して
\eqref{eq4}より
\[
\begin{split}
   D^f &= C^{fe} \\
   &= C^{k \varphi(m)+1}  \\
   &\equiv C  \mod m \quad 0\le C <m
\end{split}
\]
となるからである．攻撃者 Eve が公開鍵$[m, e]$を知りAliceの$D$を傍受したとしても，
$m$の因数分解を知らなければ$\varphi(m)$が分からず，
$\mod\varphi(m)$での$e$の逆$f$を求めることが出来ないので
$C$を復号出来ないわけである．
この方法でAlice はBobに例えば共通の秘密鍵を送り，
共有することができる．
公開暗号方式は計算が複雑で時間も掛かるので，
その後は効率よく秘密鍵を使った暗号を使うことが可能になる．

\begin{Q}
   $p=3$, $q=5$のとき$ef\equiv 1\mod \varphi(15)$となる
   $e$, $f$を探し，適当な$D\in \{1,2,\dots,15\}$を
   公開鍵で暗号化しそれを解読して下さい．
\end{Q}

\begin{Q}
github の自分のアカウントにssh で公開鍵を登録：
\begin{verbatim}
   % cd  # ホームディレクトリに移動
   % ssh-keygen # RSA で秘密鍵と公開鍵の作成
\end{verbatim}
%\begin{verbatim}
%   % cd myGithub # 適当な場所にディレクトリを作成し移動
%   % ssh-keygen # RSA で秘密鍵と公開鍵の作成
%\end{verbatim}
その後公開鍵をGitHubの自分のホームに登録して下さい．
\end{Q}


\begin{exQ}
   $p=11$, $q=13$で$ef\equiv 1\mod \varphi(pq)$となる
   $e$, $f$を探し，適当な$D$を公開鍵で暗号化し
   それを解読するプログラムを作成して下さい．
   出来るならもっと大きな素数$p$, $q$にプログラムを拡張して下さい．
\end{exQ}

\subsection{デジタル署名}
物理的書類の真実性を担保するのにサインが有効であった．
これを電子的な文書に実現するのが \term[]{デジタル署名}であり，RSAを使い以下のようにデジタル署名を作ることができる．
Alice が公開鍵$[m_A, e_A]$ と秘匿鍵$[m_A,f_A]$を作り
公開鍵を公開しているとする．


Alice がデジタル署名したい平文を$M$とし，
Bobがそれを送りたい人間で
Bob以外の人間Victorがそれを検証するとする．
Alice はまず一方向ハッシュ関数で$M$のハッシュ値を求め，
それを秘匿鍵$[m_A,f_A]$で暗号化したものを$D$とする．
$(D,M$)をBobに送る．
Bob及びVictorは公開鍵$[m_A, e_A]$で$D$を復号化し，
$M$のハッシュ値を求め同じかどうかを調べ，
同じならばデジタル署名は正しいと判断できる．

\begin{figure}[htb]
   \centering
   \includegraphics{~/Documents/GitHub/introInfo/ch5/images/figure8.1}
   \caption{デジタル署名}\label{}
\end{figure}

さらに，Alice が公開鍵$[m_A, e_A]$ と秘匿鍵$[m_A,f_A]$を作り
公開鍵を公開していて，Bobも公開鍵$[m_B, e_B]$ と秘匿鍵$[m_B,f_B]$を作り
公開鍵を公開しているとする．AliceがBobに短いメッセージ$M<\min(m_A,m_B)$を送りたいとする．
\[
   M^{f_A} \equiv C\mod m_A \quad 0\le C <m_A
\]
\[
   C^{e_B} \equiv D\mod m_B \quad 0\le D <m_B
\]
として$D$をBob に送る．Bob は
\[
   D^{f_B} \equiv C\mod m_B \quad 0\le C <m_B
\]
となるので
\[
   C^{e_A} \equiv M\mod m_A \quad 0\le M <m_A
\]
となり，$M$を復元できる．
この方式もデジタル署名と呼ばれることがある．

%https, ssh ではこのようなアルゴリズムを実装している．

\subsection{認証局}
\begin{figure}[htb]
   \centering
   \includegraphics{~/Documents/GitHub/introInfo/ch5/images/figure8.2}
   \caption{中間者攻撃}\label{}
\end{figure}
一つだけ問題なのは，公開暗号でもデジタル署名でも
Aliceが本当にAliceかAliceを騙った別人ではないか，
Bobは本当にBobなのかBobを騙った別人ではないかを
別の方法で確認しておく必要があることである．
例えば，Eveが最初Bobを騙ってAliceと公開キーを共有し，
EveがAliceを騙ってBobと公開キーを共有すると
AliceとBobの間の通信がEveを通して行われるようになり，
二人の通信が筒抜けになるという 
\term[]{中間者攻撃}(man-in-the-midle Attack)が知られている．


そのために使われるのが\term[]{認証局}(Certificate Authority)である．
まず，Alice とBobは既にある認証局に信頼を確立しているとする．
BobからAlice への公開キーを送るとき，
Bobから公開キーとそのBobによるデジタルサインを認証局へ送り，
認証局でBobから送られてきたことを確認する．
その公開キーと認証局によるデジタルサインをAlice に送り，
認証局から送られてきたことを認証し，
Bobが本当にBobであることをAlice に認証する(図~\ref{CA})．
秘密鍵はどこでもインターネットを通っていないことに注意する．
Alice がBobへの公開キーを送るときも同様である．

\begin{figure}[htb]
   \centering
   \includegraphics{~/Documents/GitHub/introInfo/ch5/images/figure8.3}
   \caption{Bob がAlice へ認証局を通して公開キーを送付}\label{CA}
\end{figure}

ただし，Alice とBobは既にある認証局に信頼を確立しているとは限らない．
そこで，まずブラウザーが信頼できる認証局を幾つか
初期設定で知っているようになっている．
幾つかの認証局同士はお互いに公開キーと基本情報を認証した
デジタル署名を共有していて，
それをたどって認証できる認証局のみを信頼し，
その認証する相手を認証するという
信用の輪を広げる仕組みになっている．

\subsection{暗号を実装するプロトコル}
暗号や電子署名は色々な方法で実装化されている：
\begin{itemize}  
   \item 
       IPsec：IPパケット単位でのデータの改竄防止を行うプロトコル
  \item
     SSL =secure socket layer, TLS =transport layer security：
     暗号通信を行うためのプロトコル
  \item
     HTTPS, SMTPS, FTPS：
     元々のプロトコルのセキュリティーを強化したプロトコル
     (図~\ref{digitalSign}).
     
\begin{figure}[htb]
   \centering
   \includegraphics[width=9cm]{~/Documents/GitHub/introInfo/ch5/images/certificate.png}
   \caption{ブラウザーでの認証証明}\label{digitalSign}
\end{figure}
   \item  
      SSH =secure shell：公開鍵暗号を使って，ネット上でログインしたり，ファイルを転送するプロトコル
\end{itemize}

また，それらの実装は日々問題点が表れ，それに対応する拡張が行われているので，CRYPTRECなどで確認する必要がある．


%\begin{thebibliography}{9}
%   \bibitem{H} 
%      堀田 良之, 加群十話: 代数学入門 (すうがくぶっくす 3), 朝倉書店, 1988.
%   
%   \bibitem{Ko} Koblitz, A course in Number theory and Cryptography, GTM,  
%   \href{https://link.springer.com/search?package=11649&facet-content-type=%22Book%22}{ダウンロード}
% 
%    \bibitem{N}
%      NISC, インターネット安全・安心ハンドブック
%   \href{https://security-portal.cyber.go.jp/guidance/handbook.html}{ダウンロード}
%
%   \bibitem{ST} Silverman, Tate,
%   Rational Points on Elliptic Curves, Springer, 2015 
%   \href{https://link.springer.com/search?package=11649&facet-content-type=%22Book%22}{ダウンロード}
%   \bibitem{T} 
%      高木貞治, 初等整数論講義, 共立出版, 1971.
%   \bibitem{Y}
%      結城浩, 暗号技術入門 第３版, 2015.
%\end{thebibliography}

\section{セキュリティーマネージメント}
\subsection{セキュリティーの基本概念}
コンピュータやネットを使った活動が安全かつ正確に出来るようにすることを\term[]{情報セキュリティ}(Information Security)といい，それへの対応を
\term[]{情報セキュリティマネージメント}(Information Security Management)
という．

情報セキュリティについての基本的な概念として以下のものがある：
\begin{itemize}
   \item 
      情報セキュリティーのCIA：基本の3つの柱
      \begin{itemize}
         \item \term[]{機密性}(Confidentiality)：権限を持った人やデバイスのみが情報にアクセスできること
         \item \term[]{統一性}(Integrity)：情報の改竄や破壊や消去が出来ない，起こらないこと
         \item \term[]{有用性}(Availability)：情報にアクセスしたい時にその情報にアクセスできること
      \end{itemize}

   \item 
      次の4つも合わせて情報セキュリティの7要素という：
      \begin{itemize}
         \item 真正性(Authenticity)：個人，組織，情報システム，情報
            が正当でなりすましがない状態．
         \item 責任追従性(Accountability)：情報を閲覧変更した個人，組織を
            特定，追跡できる状態．
         \item 否認防止(Nonrepudiation)：情報を閲覧変更したことを
            後から否認できないこと．
         \item 信頼性(Reliability)：情報システムが意図した通りに動くこと．
      \end{itemize}

   \item 
      組織や個人に必要な価値のあるものを資産というが，
      特に情報システムや情報自体を\term[]{情報資産}という．

   \item   
    
      攻撃側と守備側の関係：
      \begin{itemize}
         \item 
            \term[]{脅威}(Threat)：情報システムや情報に危害を与えうる状況や事実
         \item 
            \term[]{脆弱性}(Vulnerability)：情報システムや情報に脅威が影響を与える
            部分があること．
         \item 
            \term[]{リスク}(Risk)：脅威が脆弱性のある情報システムや情報に
            影響を与える確率の度合い．
      \end{itemize}
      リスクは大雑把に以下のように評価できる：
      \[
      \begin{split}
      	&\text{リスク}\\
      	&=\text{情報資産の価値}\times\text{脅威の大きさ}
         	\times\text{脆弱性の度合い}
      \end{split}
      \]
   \item 
   \[
      \text{セキュリティーシステムの強度}
      \le
      \text{構成要因のもっとも弱い部分の強度}
   \]
   \item 
      脅威には三つの側面がある：
      \begin{itemize}
         \item 物理的脅威 
         \item 人的脅威
         \item 技術的脅威  
      \end{itemize}
      
      情報セキュリティーのことを(ネットワークの)技術的脅威のことに
      だけに限定するのは危険である．
      
   \item 
      攻撃の目的：
      \begin{itemize}
         \item 愉快犯
         \item スクリプトキディ
         \item お金 RaaS 
         \item ハクティビズム：ハッカーの主義主張を広げるため
         \item 貴重な個人や組織の情報の収集やダークウェブでの売買
         \item 内部関係者の組織や個人への恨みへの復讐
         \item サイバーテロリズム・サイバー戦争：
         
            ロシアのウクライナ侵攻や米軍・イスラエル軍のイラン戦争の始まる少し前，
            全土で例を見ないサイバー攻撃が起こり，
            社会の通信網を混乱させたり重要インフラを使用できなくした
            後に現実の戦争が起こった．
         
         \item 敵国による世論操作・選挙介入：
         
            民主国の選挙で陰謀論を広めて社会に混乱を起こしたり，
            選挙妨害などの選挙介入が行われ
            敵国に有利な指導者を生み出す工作，認知戦が行われている
            (\cite{Ba}, \cite{NKO})．
      \end{itemize}
\end{itemize}


\subsection{技術的脅威 }
技術的脅威，特にネットワーク技術に関連した脅威であるサイバー攻撃について纏める．
\begin{itemize}      
      \item 
      \term[]{サイバーキルチェーン} (Cyber Kill Chain)：
      キルチェーンは軍事用語で戦闘の進行段階を示すものだが，
      それをサイバー攻撃の分析に使ったのがサイバーキルチェーンである．
      次の様な攻撃の段階がある：
      \begin{enumerate}
         \item
            偵察(Reconnaissance) ：標的やその脆弱性を探すこと．
            
            OSINT(Open Source Intelligence)によるデータ収集を含む．
         \item
            武器化(Weponization)：攻撃ツールを準備する段階
         \item
            配達(Delivery)：攻撃ツールを標的に届けること
         \item
            エクスプロイト(Exploit)：攻撃ツールを情報システムで実行すること
         \item
            インストール(Install)：永続的なアクセスを確保するために
            マルウエアなどをインストールすること
         \item
            Command and controll(C\verb/&/C)：
            C\verb/&/Cサーバーとの通信の確立
         \item
            目的の実行：データの収集・破壊など
            脅威の目的を達成すること
      \end{enumerate}
      逆に言うと，
      もし途中の段階で阻止できれば被害を食い止めることが出来る．
   \end{itemize}
   具体的な攻撃方法を紹介する(がもちろん完全ではない．
   常に新しい攻撃法が開発されているので，まめに確認するようにしてほしい)：
   \begin{itemize}   
    \item
      一般的攻撃手法
         \begin{itemize}  
             \item       
               不正アクセス
            \item    
               クレデンシャルスタッフィング
            \item    
               脆弱性の攻撃
            \item    
               ゼロデイ攻撃
            \item    
               マルウエア攻撃
            \item    
               ドライブバイダウンロード：
               ウェブなどでユーザーに意識させずにソフトウエアをダウンロードさせる攻撃
            \item
                標的型攻撃：攻撃相手を決めて徹底的に攻撃するやり方．
        \end{itemize}
     \item
         パスワード攻撃：
         \begin{itemize}
            \item
               ブルートフォース攻撃
            \item
               リバースブルートフォース攻撃：
               パスワードを固定して利用者id を変化させる攻撃
            \item
              パスワードリスト攻撃：
              何処かで入手したパスワードのリストを使う攻撃
            \item
              レインボーテーブル攻撃
            \item
             スニッフィング：盗聴でパスワードを盗む攻撃
         \end{itemize}
      \item 
         マルウエア・不正プログラム：
         \begin{itemize}   
            \item   
               ルートキット
            \item   
                バックドア
            \item   
               スパイウエア
            \item   
                ウイルス
            \item 
               ランサムウエア
            \item 
               ボット
            \item
               トロイの木馬
            \item
              ワーム
            \item 
              偽セキュリティ警告
         \end{itemize}

   \item 
      Webを使ったサイバー攻撃の具体例：
      \begin{itemize}   
         \item    
            SQL インジェクション：
            Web ではログインのためのデータベースとしてSQLを使う
            ことが多く，適切な対策が取られていないと，
            そこに直接SQL文を書き込むことでログインしたり，
            利用者の情報を詐取されることがある．sqlmap
         \item    
            OS コマンド・インジェクション：言語によりShellが呼び出せる
            関数を持つものがあり，そこのコマンドを直接入れて実行可能である．
         \item    
            ディレクトリ・トラバーサル：相対パスなどを使って公開されていない
            ディレクトリに外部からアクセスし，情報を詐取すること．
         \item
            バッファオバーフロー
         \item    
            セッション・ハイジャック：ウェブサービスでセッションid を発行する
            ものがある．その管理がずさんだとセッションid を推測され，
            それを使って セッションの乗っ取りを行う
          \item    
           クロスサイトスクリプティング(XSS)：
           不正なスクリプトをクライアント側で行わせる攻撃
         \item    
            クロスサイトリクエストフォージェリ(CSRF)：
            不正なスクリプトをサーバ側で行わせる攻撃
         \item    
            HTTPヘッダインジェクション：  
            HTTPヘッダの中に 不正なスクリプトを埋め込むなどの攻撃
         \item    
            フィッシング, Click Fix：
            信頼できる機関を装い偽のWebサイトに誘い出し
            個人情報などを集める中間者攻撃
         \item    
            CAPTCHAを使った攻撃
      \end{itemize}
      
   \item 
      Networkを使ったサイバー攻撃の具体例：
      \begin{itemize}   
         \item    
            サービス拒否攻撃
            (DoS, DDoS =Distributed Denial of Service, ディードス)：
            多量のネットアクセスを起こさせサービスを麻痺させる攻撃
          \item     
           メールヘッダインジェクション
         \item
            プロンプトインジェクション
         \item    
            FreeWiFi の使用に付け入る攻撃
         \item    
            リモートディスクトップ
         \item    
            VPNを使った攻撃
      \end{itemize}


\end{itemize}


\subsection{技術的脅威への対応}
技術的脅威における攻撃のそれぞれに対応の仕方があり
(もちろん，新しい脅威には対応の方法が分からない場合もあり，その場合は脆弱性のあるサービスを無効化する必要がある)，
それぞれに高度な専門性が必要になるが，
ここでは一般的な最低限の対応について纏める．

\emptyline
\begin{itemize}
   \item
      \term[]{認証}(authentification)：ある人が本人であることを確かめること
   \item
      \term[]{認可}(authorization)：
      どこまでアクセスを許可するかの決めること，その決まり

\end{itemize}

\begin{ex}
   ホテルに予約した後，
   チェックインするとき予約した人間あることを確認することが認証であり，
   自分の泊まる部屋の鍵をうけとり，部屋の使用権を確保するのが認可に当たる．
   このとき鍵は自分の泊まる部屋のみにしか使えないので
   使用範囲を最小限に抑えていることがわかる．
\end{ex}

認証の手段として
%所持認証，知識認証，生体認証の三つに分けられる．
\begin{itemize}
   \item
      所持認証：持っているもので認証すること．例は自分の家の鍵   
   \item
      知識認証：知っている知識で認証すること．例はパスワードや暗唱番号
   \item
      生体認証：体の情報を元に認証すること．例は指紋認証や顔認証
\end{itemize}

% 一個だけでは破られる可能性があるので，幾つかを組み合わせた認証方式が考えられている．

%\emptyline

%セキュリティに関して認可も大切で，

%\subsection{ }

\emptyline
技術的脅威への対応
\begin{itemize}
   \item
      万が一その人の権限が攻撃者に奪われた場合も最小の被害で済むように
      アクセス権は必要最小限に限る(\term[]{最小アクセス権の原理})のが
      基本の方針
   \item 
      パスワード管理：ただし後で指摘するように，以下を実現するのは
      普通の人間には極めて難しことである．
      \begin{itemize}
         \item 
            予想が難しいパスワード
         \item 
            定期的に変更
         \item 
            使い回しはダメ
      \end{itemize}
   \item
      特に二つの認証方式を組み合わせたものを\term[]{二要素認証}という．
      キャッシュカードはその所持と
      暗証番号PIN(Personal Identification Number)の知識
      を組み合わせた二要素認証の例である．
      最近ではホームページのサービスの登録の本人確認で，
      自分のアドレスを登録するとそのアドレスに数字をメールしてきて，
      それをhttpsのホームページに貼り付けると本人確認が済む
      という二要素認証がよく使われている．
      三つ以上の認証方法を組み合わせた
      \term[]{多要素認証}が使われる場合もある．

   \item
      \term[]{パスキー}{Passkey}：
      各サーバーごとに最初にローカルで生成した
      公開鍵をサーバーに登録し，
      秘密鍵はログイン名とサーバーのアドレスと共に大切に
      ローカルの端末に保存しておく．
      以降保存されたサーバーのアドレスを訪れると，
      サーバー側で毎回ランダムに生成される数字を
      公開鍵で暗号化した \term[]{チャレンジ}というものを
      送り返してくるので，
      ローカルの秘密鍵で解読し，
      サーバーと照合し合致すればログイン出来る．
      秘密鍵を使うための本人確認は
      ローカルの認証器(所有認証)の生体認証や
      パスフレーズ(パスワードに近いが秘密鍵の解除のみに使われる，
      知識認証)などで行う．
      秘密鍵・生体認証・パスフレーズはログインするとき
      インターネット上に一切送られていないこと，
      チャレンジは毎回廃棄されるので第三者に見られて構わないこと
      に注意してもらいたい．
      また，フィッシングサイトを自分から訪れることもないし，
      フィッシングサイトにはアルゴリズムからもログイン出来ない
      ことに注意してほしい．
\begin{figure}[htb]
   \centering
   \includegraphics{~/Documents/GitHub/introInfo/ch5/images/figure9.1}
   \caption{パスキー}\label{}
\end{figure}

   \item
      ファイヤーウォール(FireWall)：
      内部ネットワーク とWANの間おいて
      通過できるパケットのIPアドレスとサービスの種類を制限する
      装置やアプリケーション．
      外部に公開したいメールサーバーやWebサーバが置かれる
      DMZ=DeMilitarized Zone が設置されることが多くなった(図~\ref{FireWall})．

\begin{figure}[htb]
   \centering
   \includegraphics{~/Documents/GitHub/introInfo/ch5/images/figure9.2}
   \caption{ファイヤーウォール}\label{FireWall}
\end{figure}

   \item
      OSやアプリを常に最新のものに保つ．
      
   \item
      ウイルス対策ソフトの導入．

   \item
      VPN(Virtual Private Network)：通信を秘匿化して送る方法．
      LANの外からLANの中へアクセス出来るようになるが，
      最近ではそれが脆弱性の元になることが分かり，
      使うことを規制する動きがある．
      
\end{itemize}



\subsection{人間的・物理的脅威と対応}
セキュリティーでは技術的側面だけでなく，人間的・物理的側面も重要になる．
例えばシステム管理者に「難しいパスワードを定期的に変更し，パスワードの使い回しはするな」と言われていても
面倒くさくてその指導に従っていない人の方が多いと推測される（皆さんの中にはこんな
\href{https://nordpass.com/most-common-passwords-list/}{酷いパスワード}を使っている人はいないと思うが$\cdots$)．
しかし，一般にセキュリティーを強くすると使い勝手が不便になることも多い．
%また，倫理・組織の防御方針・法律・一部企業の利益主義・他国による攻撃など政治・経済に関連した側面も大事である．


\emptyline



本気で攻撃を仕掛けてくる相手(国)は
物理的・人間的なセキュリティーの脅威も突いてくる．
特に色々な組織の中枢に近いものは
そこに近づく前からセキュリティー全般に注意を払う訓練が必要である．
また，フィッシングはネットワークセキュリティーと
人間的なセキュリティーの二つの側面を持っているように，
この三つは独立している訳では無いが，便宜上三つに分けて解説する．
\emptyline

物理的なセキュリティー ：
\begin{itemize}
   \item 
      自分の部屋を出るときには必ず施錠する．

   \item      
      自分の部屋の外では鍵やカードや認証器は肌身離さず持っている．
         
   \item 
      他人がいる可能性のある場所でデジタル器機を離れるときは
      画面をロックしてパスワード，生体認証がなければ使えないようにする．

   \item 
      機密情報の入ったUSBメモリーは極力持ち歩かない．
      保管場所には施錠する．
      
      その組織のロゴなどの付いたUSBメモリーを駐車場などにわざと落として，
      それを拾わせて内部のパソコンにウイルスを感染させる攻撃もある．

   \item 
      次のような古典的な攻撃にも気を付ける：
      \begin{itemize}
         \item ショルダーハッキング：肩越しに暗証番号を盗み見
         \item トラッシング：ゴミ箱あさり
         \item スキミング：クレジットカード等の情報の電子的な盗み取り
         \item
            ピギーバック(piggyback), 共連れ：
               鍵のない人が鍵を持った人と鍵が無いと入れないところに一緒に入ること
            
            ただし，目の前に鍵がなくて困っている人がいたら
            邪険にすることも出来ないので，
            声を掛けて顔を確認して入ってもらうことにしても良いだろう．
	\end{itemize}
  
   \item      
      自分や知人の写真・PID・マイナンバー・暗証番号・生体認証情報をみだりに公表しない．
      
   \item 
      写真の日付・GPS情報などは不用意に公表しない．
      また映り込みから場所を特定されることがあるので，
      それにも十分に気を付ける．
      ネットで「特定屋」と検索すると数多く見つかる．
   \item
      SNSの個人情報はダダ漏れなので個人情報はあげない．
      
      \noindent
      例えば
      旅行中の写真をその時点でインターネットにアップすると
      その時で家にいないことを泥棒に知らせてしまう危険がある．
   \item 
      紛失防止タグを使った家の特定：ストーカーなども使っている．

   \item
      スマホをオンにしての盗聴が可能になっている．
      秘密の会議の時はスマホを預けるなどの
      セキュリティー対策が必要になっている．
      会議室の最大限のセキュリティーを保つには
      電磁シールドの敷設も必要だろう．

   \item 
      ランサムウエア攻撃でバックアップも同時に暗号化されてしまい
      データを復元できないことが頻繁に起こっている．
      ネットに繋がっていないデバイスへの定期的なバックアップを
      お奨めする．
      
   \item
      政治家や企業経営者は
      パワハラ・セクハラ・不倫・週刊誌の取材・税務調査・スパイに気を付ける．
      
   \item
       \begin{itemize}
         \item 
            某国は闇バイトによる使い捨て工作を行っていると聞く．
         \item 
            Google アースやストリートビューの情報から
            具体的な攻撃を考える攻撃者もいる．
            例えば，駐車場に止まった業者のトラックの情報を使って
            その施設への侵入を試みるという攻撃が知られている．
            
       \end{itemize}
      
      セキュリティ責任者はそのような観点にも注意が必要．

\end{itemize}

\begin{Q}　
   \begin{itemize}
      \item 
         スマホの写真の位置情報・時間情報を確認・削除する方法を調べてください．
         他人と写真を共有するときはこのようなメタ情報を消すことをお奨めします．
   
      \item 
         また，スマホやパソコンは貴方の個人的な情報を常に収集しています．
         (私のiPhone だと)
         Setting/Privacy and Security の\href{https://www.youtube.com/watch?v=XtOD7Cqd_H0}{アプリとアップルのアクセス権}を必要最小限に設定してください．
   
      \item    
         Google などはCookieを使って皆さんの検索履歴などから情報を分析し，
         インターネットでオークションに掛け，巨万の富を得ています．
         自分が使っているBrowser 毎にCookie を消す方法を調べてください．
         定期的に消すことをお奨めします．
      \item    
         Google ストリートビューで自分の家・実家がどう写っているか
         確かめてください．危険性が確認できた場合対策を取って下さい．
   \end{itemize}
\end{Q}

\term[]{ソーシャルエンジニアリング}：人間の心理的弱点を突く攻撃

\begin{itemize}   
   \item 
      不審な直接のコンタクト・電話・メールに気を付ける．
   
   \item 
      不審でないと思われる直接のコンタクト・電話・メールにも
      ソーシャルエンジニアリングでないかと気を付ける．
   
   \item 
      Web ページのハイパーリンクの下線やQR コードは
      見た目だけからは何処へのリンクか分からないので，
      軽率にそのリンクに飛ばず，リンク先を確認してから訪問する．
      
%    \item 
%    
%      \begin{itemize}
%         \item 
%             インターネットのアルゴリズムの引き起こすフィルターバブル\footnote{「インターネットの検索サイトが提供するアルゴリズムが，
%各ユーザーが見たくないような情報を遮断する機能」（フィルター）のせいで，
%まるで「泡」（バブル）の中に包まれたように，自分が見たい情報しか見えなく
%なること．(Wikipedia 2024/04/07, 18:34)} (\cite{P})・
%エコーチェンバー\footnote{自分と似た意見や思想を持った人々の集まる空間
%（電子掲示板やSNSなど）内でコミュニケーションが繰り返され，
%自分の意見や思想が肯定されることによって，それらが世の中一般においても
%正しく，間違いないものであると信じ込んでしまう現象．
%(Wikipedia 2024/04/07, 18:35)}，
%アテンションエコノミー\footnote{情報の優劣よりも「人々の関心・注目」という希少性こそが
%経済的価値を持つようになり、それ自体が重要視・目的化・資源化・交換財化
%されるようになるという実態を指摘した概念．(Wikipedia 2025/03/21, 8:12)}
%         \item 
%          AIやbot によるフェイク動画・音声の拡散
%             
%       \item 
%         敵国の情報戦・認知戦(\cite{NKO})
%      \end{itemize}
%により
%       現在の情報空間・インターネットは
%       誤情報・フェークニュース・プロバガンダであふれている．
%       
%    \item 
%      そのために世界認知が歪み，陰謀論に堕ちる人も多い．
%      
%      ネットや知り合い，各国の指導者からの情報も鵜呑みにしてはいけない．
      \item 
         AIによるソーシャルエンジニアリングの高速化

      \item 次のような情報収集が行われている：
      \begin{itemize}
         \item 
            フィッシングなどネットを使った違法な情報収集
            
         \item 
            個人情報を調べるための景品付きアンケートなどの
            合法な情報収集(\cite{Ba})
                        
         \item
            格安のスマートテレビ(Android TV)に搭載されたACR (=Automated Content Recognition 自動コンテンツ認識)による情報漏洩
         \item
            イランへの攻撃では，モサド・アマン(8200部隊)が
            監視カメラをハッキングしたり，
            攻撃先に秘密カメラを仕込んだりしていたことが明らかになった．
            
         \item
            一部の国は国家情報法を持ち，
            海外の自国民を監視するため外国に秘密警察を持っている．
            他国民を監視していない保証もない．
           
      \end{itemize}
 
      
      
   \item 
      エプスタイン事件からも分かるように
      ハニートラップも現実に行われている．
      裏に他国の諜報機関がいたとも言われている．
      某大統領はその国に弱みを握られているという噂もある．
      
      そのような罠に気を付ける．

   \item 
      最近の事例：
      \begin{itemize}
         \item 
            個人の生成AIへの機密情報の提示，漏洩
         \item
            BeReal などからの機密情報の提示，漏洩
         \item
            認知戦
      \end{itemize}
\end{itemize}

\begin{Q}
   Web ページのハイパーリンクの下線やQR コードで
   文字でリンクを表示する方法を調べてください．
   これからは毎回クリックする前にチェックすることをお奨めします．
\end{Q}

\cite{GM}にはKali Linuxや
ネットを使ったソーシャルエンジニアリングの具体的手法の解説があるので
是非参考にしてほしい．
ただし，これは攻撃のマニュアルともなるのでそこには十分注してください．

\begin{Q}
   \href{https://haveibeenpwned.com/}{Have I Been Pwned}
   を使って
   メールアドレスに付随した
   password が Dark Web に漏洩していないかチェックして下さい．
   もし漏洩が確認できた場合対策をお願いします．
   まめにチェックすることをお奨めします．
\end{Q}

%\begin{thebibliography}{99}
%   \bibitem{Ba} 
%         ジェイミー・バートレット, 
%         操られる民主主義: デジタル・テクノロジーはいかにして社会を破壊するか, 
%         草思社, 2018
%
%      \bibitem{GM}
%       Joe Gray, Jin Maeda,
%      人を動かすハッカーの技術：ソーシャルエンジニアリングの実践と防御 
%
%    \bibitem{IPA}
%       情報処理推進機構, 
%       \href{https://www.ipa.go.jp/security/vuln/websecurity/about.html}{安全なウェブサイトの作り方}
%
%    \bibitem{IPA2}
%       情報処理推進機構, 
%      \href{https://www.meti.go.jp/policy/netsecurity/downloadfiles/guide_v3.0.pdf}{サイバーセキュリティ経営ガイドライン}
%
%    \bibitem{IPA3}
%       情報処理推進機構, 
%     \href{https://www.ipa.go.jp/security/guide/sme/ug65p90000019cbk-att/sme_guideline_v4.0.pdf}{中小企業の情報セキュリティガイドライン}
%    \bibitem{JPCERT} 
%       一 般社団法人 JPCERT コーディネーションセンター, 
%       \href{https://www.jpcert.or.jp/csirt_material/files/manual_ver1.0_20211130.pdf}{インシデントハンドリングマニュアル}
%    \bibitem{HB}
%      NISC,
%      \href{https://security-portal.cyber.go.jp/guidance/handbook.html}
%      {インターネット安全・安心ハンドブック} 
%    \bibitem{N} 
%    一般社団法人 日本シーサート協議会シーサート人材ワーキンググループ, 
%    改訂新版 セキュリティエンジニアの教科書 一般社団法人 日本シーサート協議会シーサート人材ワーキンググループ  
%    
%
%   \bibitem{NKO} 
%      長迫, 小谷, 大澤, SNS時代の戦略兵器 陰謀論 民主主義をむしばむ認知戦の脅威,
%       ウェッジ, 2025.
% 
%     \bibitem{O} 
%       小川和久, 総理、国防も安全も穴だらけ！国民を守れない国・ニッポン
%
%  
%   \bibitem{P} 
%      イーライ・パリサー, 閉じこもるインターネット——
%      グーグル・パーソナライズ・民主主義 , 早川書房, 2012.
%      改題：フィルターバブル──インターネットが隠していること
%      (ハヤカワ文庫NF), 早川書房, 2016.
%
%   \bibitem{SS}
%    瀬戸, 斎藤, 徹底攻略情報セキュリティマネージメント, インプレス   
%      
%    
%\end{thebibliography}
%
%\end{document}


%\documentclass[a4paper,12pt]{amsart}
%\usepackage[dvipdfmx]{graphicx,xcolor}%graphics
%\usepackage{float}
%\usepackage{amsmath}
%\usepackage{verbatim}
%\usepackage{latexsym}
%\usepackage{amssymb}
%%\usepackage{emath}
%\usepackage{listings}
%\usepackage{ccicons}
%\usepackage{mflogo}
%%\usepackage{fancybox}%enclose by box
%\usepackage[dvipdfmx]{hyperref}
%\usepackage{verbatim}
%\hypersetup{% hyperrefオプションリスト
%   setpagesize=false,
%   bookmarksnumbered=true,%
%   bookmarksopen=true,%
%   colorlinks=true,%
%   linkcolor=red,
%   citecolor=blue
%} 
%
%\newtheorem*{Thm}{定理}
%\newtheorem*{Prop}{命題}
%
%\newtheorem*{exQ}{レポート}
%
%\theoremstyle{remark}
%\newtheorem*{ex}{例}
%\newtheorem*{Q}{演習}
%\newtheorem*{Quiz}{問題}
%
%\renewcommand{\figurename}{図}
%\renewcommand{\proofname}{証明}
%\renewcommand{\tablename}{表}
%\setlength{\belowcaptionskip}{0mm}
%
%\DeclareMathOperator{\Ker}{Ker}
%\DeclareMathOperator{\Image}{Im}
%
%\def\emptyline{\vspace{11pt}}
%\def\smallemptyline{\vspace{6pt}}
%\begin{document}
%
%\begin{flushright}
%   %2026/04/08   
%   2026/07/14
%   
%   九大・数理 大津幸男
%\end{flushright}
%
%
%\begin{center}
%   {\Large 情報科指導法II---11}
%   \href{https://creativecommons.org/licenses/by-nc-nd/4.0/}{\ccbyncnd}
%   \par{情報・社会・AI・健康}
%\end{center}

\chapter{情報・社会・AI・健康}
\section{情報とは}
 Data(データ)は数字や文字や図形などで表されたもの，
Information (情報) はデータに意味を持たせたもの，
Knowledge(知識)は情報を解釈して体系化したもの，
Wisdom(知恵)は知識を体得して人生に役立つものとしたものである．DIKW ピラミッドとはそれらを以下のように図式化したものである：
\begin{figure}[htb]
   \centering
   \includegraphics{~/Documents/GitHub/introInfo/ch6/images/figure.1}
   \caption{DIKW ピラミッド}\label{DIKW}
\end{figure}
 
 例えば生成AIにレポートを丸投げして内容を確認せず提出するのは，
 情報のみ得ることに対応する．出てきたレポートを自分なりに解釈し
 自分の言葉で喋れるようにするのは知識の習得にあたる．
 生成AIと壁打ちしながら自分に欠けていた視点からの考察を深め，
 自分だけでは書けなかったレベルのレポートをAIと共同で作成するのは
 知恵を得る使い方である．
 
%\section{情報社会とビジネス}
%近年ではビジネスや経営，マーケティングで情報技術を使うことが多くなった．
%
%CRM（Customer Relationship Management：顧客関係管理）

\section{情報社会と法律}
情報社会と情報セキュリティーに関係する日本の法律：
\begin{itemize}
   \item
      知的財産権

   \item
      著作権

   \item
      個人情報保護法
   \item      
      不正アクセス禁止法
   \item      
      サイバーセキュリティ基本法
      
   \item
      電気通信事業法
     \item    
      電子署名及び認証業務に関する法律
    \item     
      電波法
   \item      
      特定電子メールの送信の適正化等に関する法律
   \item      
      有線電気通信法
\end{itemize}


\section{AIの発達の歴史}
\begin{itemize}
   \item
      ロジックセオリスト：
      1955年から1956年にかけてアレン・ニューウェル，ハーバート・サイモン，
      J・C・ショーが開発したプログラム．
      人間の問題解決能力を真似するよう意図的に設計された世界初のプログラム
  
   \item 
      イライザ ELIZA：1964年ワイゼンバウムが開発した対話システム．
 
   \item  
      エキスパートシステム
   
%   チューリングテスト：
%   チューリングの考案した思考実験。機械と人間を別々の密室に入れ、第三者がそれぞれと交信して、機械と人間の判別がつかなければ、機械は人間と同じ思考能力をもつにいたったとする。
%   chatGPTにより不十分であることが示された．
   \item  
     neural network ニューラルネットワーク，
     machine learning 機械学習，
     deep learning  深層学習などのモデルが研究された．
 
    \item    
      ディープブルー： 1997年チェスの世界チャンピオンカスパロフに勝った
    \item    
      ワトソン： 2011年クイズ番組ジョパディで人間に勝利
   
%   ケンブリッジ・アナリティカ：
%   2016の大統領選でフェースブックのデータを使いトランプ大統領の勝利を導いた．アメリカで大問題となった．
    \item       
      AlphaGo 2016年碁で名人に勝利，  
      GAN =Generative Adversarial Network：敵対的生成ネットワークを使った
  
%   CNN たたみ込み型ニューラルネットワーク 
%   generative 生成的
%   adversarial 敵対的
%   generator 生成器
%   discriminator 識別器
%    RAG =retrieval Augmented genration
%   Claude 4 アンソロピック
%   Claude Mythos ミトス：

     \item
        LLM =large language model大規模言語モデルが発達し
        GPT-=Generative Pre-Trained Transformerが開発され，
         2023年それを対話に使うChatGPT が発表され
         世界に生成AIブームをもたらし，株もAIバブルの様相を示し，
         現代に至る．

%    OpenAI  Sam Altman CEO
%   ChatGPT 
%   GPT-3の次世代バージョン　
%   GPT-4 2023/3
   
\end{itemize}


\section{情報社会と心と体の健康}
最近，体やメンタルの不調を訴える人・社会的な問題を起こす人が増加している．

%
%{\Large 質問：}
%\begin{enumerate}
%	\item 
%	スマホを一日何時間使っていますか？
%	（iPhone ならScreen TimeでAndroidならDigital Wellbeing で確認できます）
%
%	\item 
%	いつも「ながらスマホ」していませんか？
%
%
%	\item 
%	猫背・ストレートネック・視力の低下・その他の目の不調・耳の不調・
%	頭痛・腰痛などありませんか？
%
%	\item 
%	運動はどれ位していますか?
%
%	\item 
%	一日何時間椅子に座っていますか？
%
%	\item 
%	一日何時間寝ていますか？
%
%	\item 
%	朝何時に起きますか？朝日を浴びていますか？
%
%	\item 
%	朝食は食べていますか？
%
%	\item 
%	コンビニで売っているような食べ物やファストフードをどれ位
%	食べていますか？
%\end{enumerate}
%
%\emptyline
%
\emptyline

\subsection{データ}
\begin{enumerate}
   \item 
      2023年のネット利用時間は高校生が6時間14分，中学生が4時間42分，
      小学生（10歳以上）が3時間46分だっただった
      \footnote{こども家庭庁 2023年度のインターネット利用環境実態調査}.
   
   \item 
      2024年の裸眼の視力が1.0未満の割合は高校生で71％，中学生で60％，
      小学生で36％だった
      \footnote{文部科学省 2024年度学校保健統計}．
   
   \item 
      2021年度中等度以上のうつ症状がある子どもの割合が高校生30％，
      中学生24％であると発表された
      \footnote{国立成育医療研究センター「コロナ×こどもアンケート」
      第4回調査報告}．
   
   \item 
      小中高校生の不登校の人数が2023年には過去最多の34万人を越えた
      \footnote{文部科学省 令和5年度児童生徒の問題行動・不登校等生徒指導上の
      諸課題に関する調査}．
   
   \item 
      2024年オーストラリアでは16歳未満のSNSを禁止する法律を定めた．

   \item 
      ネットカジノ依存
   \item 
      闇バイトの1年間の逮捕者の推移：
      年分トクリュウ全体の年間摘発者数うち「闇バイト」応募・関与者の人数2023年（令和5年）10,105人約3,925人（全体の約4割）2024年（令和6年）12,178人4,759人（全体の約4割）
\end{enumerate}
\emptyline


\subsection{対策の前提となる考え方}
\begin{enumerate}
  \item 
   科学的な（実験や統計や論理や数式で白黒付けられるような）
   反証可能(\cite{PK})な議論・
   対話可能な議論を行う．
   
   具体的対処の有効性は，体や心の状態・症状の改善度から評価する．
   特に，排便の状態や疲労感・睡眠不足感をモニターする．
%   \emptyline
   
   自分のことは自分では分かりづらい．
   例えば目の10cm 前でスマホを見ている人がいたら
   自分のスマホを目から離す行動をとるなど，
   他人を詳細に観察しそれが自分に当てはまらないか思い返してみるのは
   良いことである．
   
   \item 
   進化生物学：ヒトが農耕を始めてから未だ400世代程しかたっていないので，
   それ以来の生活の変化に我々の体は適応していない．
   ヒトが適応している1・２万年前の生理学的環境・状態と近い（
   逆に人工的な環境の）生活を送る必要がある(\cite{KL}, \cite{L})．
   
   \begin{enumerate}
      \item 
      ヒトは狩猟採集生活に適応しているが，運動しない生活，
      座り続ける生活，夜寝ない生活，
      長時間近くを見続ける生活，イヤホンを使い続ける生活
      には適応していない．
      
      \item 
      ヒトは狩猟採集生活の雑食の生活に適応しているが，
      多様な食物を摂らない生活，超加工食品には適応していない．
      
      \item 
      ヒトはスマホを使い続ける生活には適応していない
      (\cite{Han}, \cite{I}, \cite{Ko}, \cite{KS}, \cite{M2})．
   \end{enumerate}
	
   \item 
   動物は環境に合わせて行動を変えるように進化してきた．
   これを制御しているのが脳の報酬系
   \footnote{1954年オールズらにより最初に発見されたときは快楽中枢と
   呼ばれていたが，後に快楽ではなく，
   新しい行動や情報への報酬を与えていることが明かになり，
   報酬系と呼ばれるようになった．}
   から放出されるドーパミン\footnote{現在では脳内物質で
   説明されているが，古くからスキナーによる小動物を使った
   実験から現象としては知られていた．
   そのため，スマホをスキナー箱という研究者もいる．
   この心理学の分野を行動分析学\cite{M}という．}
   であり，
   ドキドキわくわくする感情を引き起こし，
   我々が新しい行動を行ったり新しい情報を得るのを促進するように働いている．
   ドーパミンはアドレナリンに代謝され攻撃性を増す効果もある．

   
   スマホ（のアルゴリズム）はドーパミンを継続的に大量に放出させ，
   使用者を依存症にするように作られており，
   それを作っている会社はそこから巨万の富を得ている．
   ただし，ドーパミン自体は善でも悪でもない．
   ドーパミンを上手くコントロールして，
   なりたい自分に近づいてほしい(\cite{LL})．
   
   \item
   我々のデジタルの行動データはブラウザーのCookie などに蓄積されていて，
   それが売買されたり，アルゴリズムがフィルターバブル
   \footnote{「インターネットの検索サイトが提供するアルゴリズムが，
   各ユーザーが見たくないような情報を遮断する機能」（フィルター）のせいで，
   まるで「泡」（バブル）の中に包まれたように，自分が見たい情報しか見えなく
   なること．(Wikipedia 2024/04/07, 18:34)} (\cite{P})
   ，SNSがエコーチェンバー現象
   \footnote{自分と似た意見や思想を持った人々の集まる空間
   （電子掲示板やSNSなど）内でコミュニケーションが繰り返され，
   自分の意見や思想が肯定されることによって，それらが世の中一般においても
   正しく，間違いないものであると信じ込んでしまう現象．
   (Wikipedia 2024/04/07, 18:35)}を引き起こしている．
   更に，最近ではこれらの現象に加えてアテンションエコノミーの拡大
   \footnote{情報の優劣よりも「人々の関心・注目」という希少性こそが
   経済的価値を持つようになり、それ自体が重要視・目的化・資源化・交換財化
   されるようになるという実態を指摘した概念．(Wikipedia 2025/03/21, 8:12)}
   や諸外国からの情報戦・認知戦により誤情報・偽情報が生み出され，
   それらが生成AI により著しく増幅されている．
   それらにさらし続けられることで，それらの中に所謂「真実」を見つけ
   \textcolor{red}{陰謀論に絡め取られる}人々が増えている(\cite{M1}, \cite{NKO})．
   陰謀論に惹かれることもドーパミンの放出と関連している．

   \item 
      それらによりトランプ大統領のようなポピュリズム政治家が世界中に台頭している．
    
\end{enumerate}
\emptyline

\subsection{デジタルデトックス}
\begin{enumerate}
%   \item 
%   スマホやデバイスの使用制限\footnote{\cite{PC}などを参考に．}
%   スマホやデバイスの使用禁止\footnote{\cite{PC}などを参考に．}---
%   これは一時的なもので，将来はスマホを自制して使えるようになることを
%   目標にする．
   \emptyline
   
   \item 
   生活改善：現在の生活はデジタルを除いても昔はなかった危険に囲まれている．
   もちろんスマホ＞ストレートネック＞肩こりのように直接の関連がある場合もある
   ので，どちらにしてもここから改善する必要がある．
   \begin{enumerate}
      \item 
      バランスのとれた（多様性の高い）質の高い食事を適量，よく噛んで食べる．
      \begin{enumerate}
         \item 
         野菜・果物・豆類・ナッツ・魚・肉・タマゴ・乳製品・
         発酵食品・ビタミン・ミネラルをバランスよく摂る．
         
         \item 
         コンビニで売っているような食品・ファストフードなどの
         超加工食品を食べない．これらにもドーパミンを放出させる働きがあり
         強い依存性を持っている．
         例：野菜ジュース・フルーツジュース・清涼飲料・カップ麺・
         ポテトチップス・揚げもの・加工肉・コンビニ弁当・サンドイッチ・
         ハンバーガー
         
         \emptyline				
         ただし，時たま友達や家族と一緒にこのようなものを食べるのは特に
         問題ない．逆に，他人に自分と同じ基準での食生活を求めるのは
         現実的でなく，場合によっては人間関係をこわすのでお奨めしない．
         
         \item 
         毎日朝食をとる．
         
         例：プチトマト・スプラウト・ゆで卵・納豆・豆腐・バナナ・ミカン・
         ヨーグルト・ゴマ・お酢・アマニ油・スープ・紅茶・コーヒー
         の中から何点かを適宜組み変えて摂る
         \footnote{基本パッケージ化されていて冷蔵庫を出して直ぐに食べられる
         もので(i)の要素を多数含むように選んでいる．}．
%         ご飯・味噌汁・焼きシャケ・卵焼き・キャベツの千切り・漬物
%         のような旅館の朝ご飯のようなもので続けられれば
%         それで構わないが，
%         長続きして(i)の要素を多数含むようなものを選んでほしい．
%         その意味では小型の食洗機の導入も検討に値する．}
      \end{enumerate}
   
      \item 
      バランスのとれた質の高い運動を適量，正しい姿勢で行う．
      \emptyline
   
      注意点：
      口呼吸の矯正，猫背・巻き肩・ストレートネックの矯正，
      肩甲骨の可動域の確保，
      反り腰の矯正，体の左右差の矯正，
      アキレス腱の柔軟性・アーチの保全・足指可動域の確保．
      \emptyline
   
      \begin{enumerate}
         \item 
         毎日1, 2時間ほど日光を浴び遠方を見ながらの
         ウォーキング・自転車・スポーツを行う．
         
         \item 
         長時間椅子に座り続けない．
         
         \item 
         階段の使用・片足立ち・発声練習・ストレッチ
         （専門家の指導が必要な場合は理学療法士にご相談ください）・
         足指のストレッチなどを隙間時間に行う．
      \end{enumerate}
      
      \item  
      質の高い睡眠と休養を適量，規則正しくとる．
      \begin{enumerate}
         \item 毎日8時間以上の規則正しい睡眠．
         \item 毎日朝日を浴びる．
      \end{enumerate}
      
      \item 
      バランスのとれた質の高い，社会活動を含む高度な知的活動を適量行う．
      
      \item 
      昔（殆ど）存在しなかった化学物質への被爆を避ける．
      
      %昔（殆ど）存在しなかった化学物質：
      砂糖・人工甘味料・合成洗剤(\cite{Uts})・大量のサラダ油(\cite{Ha})・
      トランス脂肪酸・PFAS・マイクロプラスチック・サプリ・抗生物質(\cite{Br})・
      薬など
      			
      \item 
%      フィルターバブル・エコーチェンバー・アテンションエコノミー・認知戦
%      について理解を深め，
      我々の世界認知がアルゴリズムや金儲けや敵国や
      生成AIから影響を受けていることを認識し，
      自分の認知が歪まされていないか常に検討することが必要である．
   \end{enumerate}
   	これらはできる限り続けることをお奨めする．
   \item
      デバイスとの付き合い方：
      \begin{enumerate}
            \item
            様々な通知機能は絶対外せないものを除いて停止．
            
            \item
            メール，SNS，ニュースの確認は一日数回に限定．週末は見ない
            と決めて知人に知らせるのもいい．
            
            \item
            集中しないといけないときはデジタルデバイスの電源を落とす．
            
%            \item
%            毎日1・2時間はスマホを見ずに，遠くの景色を見ながらの早歩きなどの
%            屋外活動を行う．日光がビタミンDを増やし近視の進行を止めることも
%            知られている．
%            ストレッチ・姿勢の矯正・片足立ち・階段の積極的使用などの運動なども
%            継続して行うことをお奨めする．
            
            \item
            スマホを含めデバイスを見るときはなるべく大きな外付けモニターを使い，
            表示フォントも大きなものを使う．
            モニターはなるべく遠くに置き，目に過度な負担をかけないようにする．
            iPhone では Setting /Screen Time/Screen Distance で画面との距離の警告を
            設定できる．
            また，
            イヤホンやヘッドホンではなく外付けのスピーカーを使い，
            他人にうるさくない音量に保つ．
            悪い姿勢を続けないように気を付ける．
            
            目も耳も体も（頭脳も）一生もの，大切に扱ってください．
      \item
         まめに\textcolor{red}{Cookie のリセット}などを行って
         フィルターバブルを防ぐことをお奨めする．
         詳しいやり方はブラウザー毎に検索するように．
      \item
         \textcolor{red}{プログラミングでデジタルデトックス}：
         プログラミングやデバッキング，論文やプレゼンの準備は
         アナログな作業．コードを書く前に，
         アナログでどれだけ綿密に計画を立てられるかで
         プロジェクトとその保守拡張の成功が決まる．
         データ構造とアルゴリズムをよく検討し，
         必要に応じて手で図を作成し，
         場合によっては数学的な性質を証明する必要がある．
         また，コーディングではパソコンでの作業が増えて
         スマートホンの小さな画面から離れられる．
         
         私の経験からのアドバイス：
         プログラミングに集中するとデジタルに距離が置けるようになる．
   \end{enumerate}

%   \emptyline
%   
%   \item 医療との向き合い方：
%   医療は必ずしも患者に良いことをやっているわけではない．
%   医療も商売の一つである．もちろん，
%   急激に体調が悪化した場合や大けがを負った場合直ぐに医療を受けることを勧める
%   が，医療とはほどほどの距離を保つことをお奨めする
%   (\cite{A},  \cite{Hi}, \cite{Me}, \cite{US})．
%   
%   \begin{itemize}
%      \item 商売で行われている治療も沢山ある．
%      		
%      \item 製薬会社のプロパガンダで作り上げられた病気もある．
%      	
%      \item 薬や治療が原因の不調も多い．
%      	
%      \item 体調が悪くないのに健康診断で病気を探すようなことは有害無益．
%   \end{itemize}
\end{enumerate}

%これをひと月続けて症状が改善するようなら，それを続けてください．
%改善しないようなら，やり方に問題なかったか反省して，やり方を変えて続けてください．それでもダメなら，既に生活改善では改善出来無いようなトラウマを抱えているか，全く別の原因がある可能性があります．
%その時こそ医師に相談することお奨めします．
%\emptyline
%
%ここに書いてあること全てを無批判に信じてはいけません．
%それではネットの情報を鵜呑みするのと変わりありません．
%以下に挙げた資料等を検討し必要ならそれ以外の文献も参考にし，
%心の底から納得したものがあればそれについてここに書いてあることを実践して下さい．



%\section{熱力学・生命・人間・社会・AI}
%\begin{itemize}
%   \item
%      エントロピーと熱力学の第二法則
%
%   \item
%      マックスウェル：マックスウェルの悪魔
%      
%   \item
%      シュレーディンガー：生命は負のエントロピーを食べて生きている．
%
%   \item
%      ショーペンハウエル：生命とは生きようとする盲目的な意志である．
%
%   \item
%      フォン・ノイマン：自己増殖オートマトン
%   \item
%      マルグリス：真核細胞の進化の内部共生説
%   \item
%      言葉・貨幣・神・法・文明の創成
%   \item
%      AI：人間と異質の知能      
%   \item
%      AIと人間との共存？
%\end{itemize}
%

%\section{授業で出来なかったこと}
%今回の授業で紹介できなかったこと：
%\emptyline
%
%基礎理論：
%\begin{itemize}
%   \item
%      サンプリング定理の証明
%   \item
%      計算可能性・再帰関数：チャーチの提言
%   \item
%      アルゴリズム
%   \item
%      計算量
%   \item
%      情報量
%\end{itemize}
%
%計算機：
%\begin{itemize}
%   \item
%      ノイマン型計算機のメカニズム
%   \item
%      オペレーティングシステム
%   \item
%      ユーザーインターフェース
%\end{itemize}
%
%プログラミング：
%\begin{itemize}
%   \item
%      Python 入門
%   \item
%      Shell script  入門
%   \item
%      Mail のプロトコル
%   \item
%      データベース, SQL入門
%   \item
%      プログラミング・プログラマーの思想と技法
%\end{itemize}
%
%
%統計と生成AI：
%\begin{itemize}
%   \item
%      データ分析：実践的技術
%   \item
%      統計理論
%   \item
%      機械学習
%   \item
%      LLM
%   \item
%      Fine turning
%   \item
%      生成AI への質問技術
%   \item
%      生成AI と著作権
%   \item
%      生成AI と兵器・戦争
%\end{itemize}
%
%これらについては，多くの本が書かれているので
%興味のある部分を自分で調べてください．

%\begin{thebibliography}{99}
%
%   \bibitem{Ba} 
%      ジェイミー バートレット, 
%      操られる民主主義: デジタル・テクノロジーはいかにして社会を破壊するか , 草思社
%
%   \bibitem{BF} 
%      ボズウェル, フィッチャー, リーダブルコード
%
%   \bibitem{B} 
%      ブレーザー, 失われてゆく我々の内なる細菌, みすず書店, 2015.
%   \bibitem{C}    
%      チャルディーニ：影響力の武器
%
%   \bibitem{H} 
%      アンデシュ・ハンセン, スマホ脳, 新潮新書, 2020.
%   \bibitem{H} 
%      ユバル・ノラ・ハラリ, サピエンス全史
%   
%   \bibitem{Ha} 
%      林 裕之, 「隠れ油」という大問題, 三五館, 2017.
%   
%   \bibitem{I} 
%      石井光太, ルポ スマホ育児が子どもを壊す, 新潮社, 2024.
%   
%   \bibitem{K} 
%      川本 晃司, スマホ失明, かんき出版, 2022.
%   
%   \bibitem{KL} 
%      ヴァイバー・クリガン=リード, 
%      サピエンス異変──新たな時代「人新世」の衝撃, 
%      飛鳥新社, 2018.
%   
%   \bibitem{KS}
%       川島 隆太, 榊 浩平, スマホはどこまで脳を壊すか, 朝日新書,  2023.
%   
%   \bibitem{L} 
%      ダニエル・リーバーマン, 人体六〇〇万年史 
%      上・下──科学が明かす進化・健康・疾病, 早川書房, 2015.
%   
%   \bibitem{LL} 
%      ダニエル・リーバーマン, マイケル・ロング , 
%      もっと! : 愛と創造，支配と進歩をもたらすドーパミンの最新脳科学, 
%   インターシフト (合同出版), 2020.
%
%   \bibitem{M} 
%   マロット, 行動分析学入門, 産業図書, 1998.
%
%   \bibitem{TH} 
%      達人プログラマー
%   
%   \bibitem{M1} 
%      物江 潤, デマ・陰謀論・カルト (新潮新書) , 新潮社, 2022.
%   
%   \bibitem{M2} 
%      物江 潤, デジタル教育という幻想 GIGAスクール構想の過ち, 
%      平凡社, 2023.
%   
%   \bibitem{NKO} 
%      長迫, 小谷, 大澤, SNS時代の戦略兵器 陰謀論 民主主義をむしばむ認知戦の脅威,
%    ウェッジ, 2025.
%   
%   \bibitem{KP} 
%      カール・ポッパー, 科学的発見の論理 上下, 恒星社厚生閣, 1971.
%   \bibitem{P} 
%      イーライ・パリサー, 閉じこもるインターネット——
%      グーグル・パーソナライズ・民主主義 , 早川書房, 2012.
%      改題：フィルターバブル──インターネットが隠していること
%      (ハヤカワ文庫NF), 早川書房, 2016.
%   \bibitem{PC}
%      キャサリン・プライス,  
%      スマホ断ち 30日でスマホ依存から抜け出す方法, 
%      角川新書, 2024.
%   
%   \bibitem{Ut} 
%   宇津木龍一, シャンプーをやめると髪が増える, KADOKAWA, 2013.
%\end{thebibliography}


%\end{document}



\appendix

%\documentclass[a4paper,11pt]{amsart}
%\usepackage[dvipdfmx]{graphicx,xcolor}%graphics
%\usepackage{float}
%\usepackage{amsmath}
%\usepackage{verbatim}
%\usepackage{latexsym}
%\usepackage{amssymb}
%%\usepackage{emath}
%\usepackage{listings}
%\usepackage{ccicons}
%%\usepackage{fancybox}%enclose by box
%\usepackage[dvipdfmx]{hyperref}
%\usepackage{verbatim}
%\hypersetup{% hyperrefオプションリスト
%   setpagesize=false,
%    bookmarksnumbered=true,%
%    bookmarksopen=true,%
%    colorlinks=true,%
%   linkcolor=red,
%   citecolor=blue,
%} 
%
%\renewcommand{\labelenumi}{\textbf{\theenumi.}}
%\renewcommand{\theenumii}{\arabic{enumii}}
%
%\renewcommand{\figurename}{図}
%\renewcommand{\proofname}{証明}
%
%\newtheorem*{PP}{今日の目標}
%\newtheorem*{Thm}{定理}
%\newtheorem*{Prop}{命題}
%\newtheorem*{exQ}{レポート}
%
%\theoremstyle{remark}
%\newtheorem*{ex}{例}
%\newtheorem*{Q}{問題}
%\newtheorem*{Quiz}{小テスト}
%
%\def\emptyline{\vspace{11pt}}
%
%
%\begin{document}
%
%\begin{flushright}
%   2026/08/20
%   %2026/08/01
%   %2026/04/21
%   %2025/12/02
%   %2025/11/22
%   %2025/03/05
%   %2024/04/06
%   
%   %大津幸男(九大・数理) 
%   snow00two
%\end{flushright}
%\begin{center}
%	\Large{Excelの関数とVBAによるプログラミング}
%	\small{
%      	\href{https://creativecommons.org/licenses/by-nc-nd/4.0/}{\ccbyncnd}
%	}
%\end{center}
%\emptyline

\chapter{Excelの関数とVBA}
\section{Excelの関数とVBA}
\subsection{Excel の関数でのプログラミング}
Excel の関数によるプログラミングではそれなりに色々なことが出来るし，
仕事ではデファクトスタンダードなので使わざるをえない．
幾つかの方法で階乗の計算をしてみる
%(\href{ftp://github.com/snow00two/docs/excelVba/factorial.xlsm.zip}{ここからダウンロード}．ただし最初はmacroを無効にして開き確認の後macroを有効にして開いてください)
．

\begin{figure}[htb]
       \centering
      \includegraphics[width=8cm]{~/Documents/GitHub/introInfo/appa/images/ScreenCapture5.pdf}
      \caption{階乗計算}\label{fun}
\end{figure}
\begin{figure}[htb]
      \centering
      \includegraphics[width=7cm]{~/Documents/GitHub/introInfo/appa/images/factTrace.pdf}
      \caption{Trace Precident の(何度かの)実行}\label{fun2}
\end{figure}


%\begin{figure}[htb]
%   \centering
%   \begin{minipage}{0.45\columnwidth}
%      \centering
%      \includegraphics[width=6cm]{~/Documents/GitHub/introInfo/appa/images/ScreenCapture5.pdf}
%      \caption{}\label{fun}
%   \end{minipage}
%   \begin{minipage}{0.45\columnwidth}
%      \centering
%      \includegraphics[width=6cm]{Screen_Capture5.pdf}
%      \caption{}\label{fun2}
%   \end{minipage}
%\end{figure}

図 \ref{fun} の2行目はセルの間の掛け算を使った関数計算である．セルあるいは数式バーに数字を普通に入力し，リターンキーを押すと通常の設定だとセルの右寄りに表示される．
数式バーでは文字列は 'の後に続け，リターンキーを押すと通常の設定だと左寄りに表示されるが，ここでは右寄りになるように表示している．
数式バーでの数式を入力するにはキーボードから$=$を入力する．
セルの値はB5のように参照され，セルをクリックすると$=$の後に自動で入力される．
上の場合セルC6 には =C5*B6 が入力されていて，そこでリターンキーを押すと計算した値がセルに現れる．
ただし，$*$は積を表す記号でキーボードから正しい位置に入力する必要がある．


一個づつ入力するのは大変で間違いが起こりやすいので以下のようにする：C5に正しい式を書いて，ここを選択し，選択の部分を少し動かすと十字の記号があらわれるので，そのまま，下までドラッグすると正しい式がコピーされる．これは相対参照という機能であるが参照先を動かしたくないときは=\$C\$5*B6 のような絶対参照を使う必要がある．
ここはよく間違えて微妙に狙ったものと異なる結果になり，慣れない人は途方に暮れるところでなので注意してほしい．Formula のタブに Trace Precident のボタンがあるのであるセルがどこを参照しているかを矢印で確認できる(図 \ref{fun2})．我々の例は小さすぎてあまり役立たないが，複雑な参照では視覚的に確認が出来て便利である．

 3行目では繰り込み関数Fact(n)を使って計算している．
 計算的にはどの方法でやっても同じであるが，
 やりたいことが既に繰り込み関数で実現されていることも多いので
 必要に応じて Help やインターネットの検索等で探してほしい．

\subsection{VBAでのプログラミング}
次に4行目で行っている
VBA を使った計算を説明する．
VBAとは Visual Basic for Applicationsの略でExcel・Word のスクリプト言語で，Basicから派生したものである．
Officeに標準搭載されているので，自分で開発環境を導入する必要が無いのは
初学者にとっては魅力的である．
\begin{figure}[htb]
	\centering
	\includegraphics[width=7cm]{~/Documents/GitHub/introInfo/appa/images/ScreenCapture1.pdf}
	\caption{Excel の設定画面}\label{}
\end{figure}

ExcelのメニューバーのExcel メニューからPreferencesを選び，
Preferencesのパネルで
Ribbon＆Toolbar＞Ribbon＞Main Tabs で Developer をチェック．
閉じてファイルを見るとDeveloper タブが見えるようになる：

\begin{figure}[htb]
	\centering
	\includegraphics[width=12cm]{~/Documents/GitHub/introInfo/appa/images/ScreenCapture2.pdf}
	\caption{Developer のタブを押した後}
\end{figure}
\begin{figure}[htb]
	\centering
	\includegraphics[width=12cm]{~/Documents/GitHub/introInfo/appa/images/ScreenCapture8.pdf}
	\caption{Visual Basic のボタンを押した後}\label{withoutCodeWindow}
\end{figure}

\begin{figure}[htb]
	\centering
	\includegraphics[width=12cm]{~/Documents/GitHub/introInfo/appa/images/ScreenCapture3.pdf}
	\caption{Moduleの追加後}\label{withCodeWindow}
\end{figure}

\begin{figure}[htb]
	\centering
	\includegraphics[width=12cm]{~/Documents/GitHub/introInfo/appa/images/ScreenCapture6.pdf}
	\caption{factorial()のコード}\label{factorial}
\end{figure}



Developer タブを選び左側のVisual Basic のボタンをクリックする
とVBA の開発環境が開く．
左側の上の窓がプロジエクトエクスプローラー，
その下がプロパティウインドウ，
左側の上の窓がコードウインドウだが最初は表示されていない(図~\ref{withoutCodeWindow})．
その下がミデエットウインドウであり，
ここに出力することでプログラムの確認が出来る．

メニューバーのinsertのModule を選ぶと，
プロジエクトエクスプローラーにModule1が追加され，
コードウインドウが使えるようになる(図~\ref{withCodeWindow})．

次に開くとき VBA を有効にするためには fuge.xlsm という拡張子で保存する必要ある．
また，
そのファイルを開くときはVBAを使用して良いかの確認アラートがでるので，自分で作ったファイルの場合 Enable Macro のボタンを押してほしい．
ただし，他人の作った.xlsmファイルのときは十分注意が必要である．

先程の例で VBA のコードをあげる (図 \ref{factorial}).


まず，変数の宣言を強制する設定にすることが推奨されているので，メニューバーのExcel のメニュー Preferences  を開き，editer のoption の二つ目の部分にチェックを入れておく．この設定が有効になっていると Option Explicit と示されている．

\begin{figure}[htb]
	\centering
	\includegraphics[width=8cm]{~/Documents/GitHub/introInfo/appa/images/ScreenCapture7.pdf}
	\caption{VBAの設定画面}\label{}
\end{figure}


プログラムに戻る．' がコメントの記号で，この部分はプログラムに何の影響も与えないが，
将来の自分・他人が理解し易いようにコメントを書くように努めてほしい．

文の後に ; は書かない．

Sub のあとに関数の名前を書く，この場合は factorial() で
必ず() が必要で，(var1, var2) のようにして変数を入れることで関数への変数の受け渡しが出来る．
次の行からEnd Subの間に定義を書く．
これらは自動補完されるので正しい入力をすれば改行で自動に現れる．
Dim が変数の宣言で
As の後に型の宣言を行っている．ここでは long は長い整数を表す．
maxNは定数のつもりで a, nが計算に使う補助的な変数である．
\begin{verbatim}
   For n = 1 To maxN
        a = a * n
        Cells(n + S, E).Value = a 'ex C(1+S,E) = E4 cell
   Next n
\end{verbatim}
がVBAの for 文でn = 1 から maxN まで a に n を掛けた
  a * n を a に書き換えそれをCells(n + S, E)に書いている．
ただし，\verb/S = 3, E = 5 /である．
この表現Cells(n + S, E)は行列と同じで(n + S, E)成分，つまり，
n+3行5列の成分を表す．
.Value はこのセルのプロパティValue , つまり，値を表している．
様々なmethod （Obj.hoge()のようなもの）も用意されているが
オブジェクト指向のプログラミングをしているというより，
Excel のbookがオブジェクトの塊となっているので，
VBA からそれを操作しているという方がもっともらしい．

\emptyline
プログラムを実行するにはメニューバーの Run のメニュー Run Sub/UseForm を選択すると実行され，シートのE行に結果を書き出しているはずである．
あと最低必要な文は if 文で，例えば
\begin{verbatim}
Sub myfun()
    If Cells(1, 1).Value < 60 Then
            Cells(1, 2).Value = "F"
    ElseIf Cells(1, 1).Value < 70 Then
        Cells(1, 2).Value = "C"
    ElseIf Cells(1, 1).Value < 80 Then
        Cells(1, 2).Value = "B"
    ElseIf Cells(1, 1).Value < 90 Then
        Cells(1, 2).Value = "A"
    Else
        Cells(1, 2).Value = "S"
    End If
End Sub
\end{verbatim}
のようにするとCells(1, 1)の値に応じてCells(1, 2)に成績を記入することが出来る．

他にも色々な機能があるが，取りあえず今説明したことでVBAを動かすことが出来るであろう．
外部ファイルとの入出力も可能なので，Excel の関数では実現できないようなプログラムも作ることが出来る．
上でも注意したがウイルスのような破壊的なプログラムを書くことも可能であるので注意が必要である．色々調べてほしい．



\begin{Q}
Excelの参照機能を使って
Fibonacci 数列を書き出せ(図~\ref{Fibonacci})．
また，同じことをVBAを使って書き出せ．
余裕のある人はExcelのシートから幾つ出力するか制御出来るようにシステムを変更せよ．
\end{Q}

\begin{figure}[htb]
	\centering
	\includegraphics[width=5cm]{~/Documents/GitHub/introInfo/appa/images/Fibonacci.pdf}
	\caption{Fibonacci 数列}\label{Fibonacci}
\end{figure}

\begin{Q}
   Excelの参照機能を使って
   Pascalの三角形を書き出せ(図~\ref{Pascal})．
   また，同じことをVBAを使って書き出せ．
   余裕のある人はExcelのシートから幾つ出力するか制御出来るようにシステムを変更せよ．
\end{Q}

\begin{exQ}
   ここまで作ってきたプログラムを生成AIに
   Python やJapaScript (+node.js )のプログラムに書き換えてもらい，
   必要に応じて訂正しそれを実行せよ．
\end{exQ}

\begin{figure}[htb]
	\centering
	\includegraphics[width=12cm]{~/Documents/GitHub/introInfo/appa/images/Pascal2.pdf}
	\caption{参照機能を使ったPascal's triangle}\label{Pascal}
\end{figure}
%\clearpage



\subsection{関数によるデータ分析}
Excel関数のもう少し詳しい使い方とデータ分析，グラフなどによる可視化について解説する．先程も説明したが，
まずセルを選ぶと数式バーに入力出来るようになる．
\verb/=/を入力すると数式を入力出来るようになる．
ここで他のセルをクリックするとそのセルの内容を参照できることになるし，
関数を書くとその関数が使えるようになる．

図~\ref{histogram}の成績集計用のExcelシート%(\href{ftp://github.com/snow00two/docs/excelVba/template.xlsm.zip}{ここからダウンロード})．
を用いて
簡単なデータ分析について解説する：
連続したセルの指定はC2:C10や C列を表すC:Cで行う．
例えばC4からC15までの平均を求めるにはC16に
\begin{verbatim} 
   =AVERAGE(C4:C15)
\end{verbatim}
と書けばいい．あるいは関数を書いてC4からC15までのセルをシフトキーを押しながら選ぶと自動で挿入される．

ヒストグラムを書くためにH4に
\begin{verbatim} 
   =FREQUENCY(C4:C15,F4:F15)
\end{verbatim}
と入力してリターンキーを押すとH4からH15に頻度が計算される．
この内容からグラフを作るとヒストグラムが出来る．
ただし，これにはそれなりの設定が必要である
(図~\ref{data}, \cite{DE}12章参照)．
H4からH15までの総和をもとめるにはH16に
\begin{verbatim}
   =SUM(H4:H15)
\end{verbatim}
と書く．これは集計が確実に行われているかのチェックにも成る．
 
条件分岐はD4の内容のように
\begin{verbatim} 
   =IF(C4>=$K$4,$J$4,IF(C4>=$K$5,$J$5,
   IF(C4>=$K$6,$J$6,IF(C4>=$K$7,$J$7,$J$8))))
\end{verbatim}
とすれば良い．その後十字の記号を下までドラッグすると
正しい式がコピーされる．
絶対参照と相対参照の使い方に注意してほしい．

L4のように
\begin{verbatim} 
   =COUNTIF($D$4:$D$15,J4)
\end{verbatim}
とすると集計がされる．これも十字の記号を下までドラッグ．
集計はヒストグラムにまとめてある(図~\ref{histogram}参)．

あと，スタイルの問題だが手入力部分の数字のフォントは青で，
いつもの固定の部分と自動計算のところは黒になるように
字の色が設定してある．
ここでは個別の成績部分と最終成績の閾値の部分がそれに当たる．


\begin{figure}[htb]
	\centering
	\includegraphics[width=11cm]{~/Documents/GitHub/introInfo/appa/images/excelAnalysis.pdf}
	\caption{データとヒストグラム}\label{histogram}
\end{figure}

\begin{figure}[htb]
	\centering
	\includegraphics[width=7cm]{~/Documents/GitHub/introInfo/appa/images/graphD.pdf}
	\caption{Graph のデータ}\label{data}
\end{figure}

\subsection{歴史と参考文献}
IBM PC市場ではLotus 1-2-3が覇権を握っていたが，
1984年のMacintoshの発売に合わせて
MicrosoftはMacintosh 版Excel を開発し1985年に発売を開始した．
その後1985年MicrosoftはWindows 1.0を開発し，
Macintosh 版Excel をWindows 版に移植し1987に発売を開始した
(\cite{Jobs} の16章にはこれらのこととS.~Jobs とB.~Gateの関係についての記述がある)．
VBAの起源はもともとは1964年にJ.~KemenyとT.~Kurtzによって発表されたBASIC =Beginner's  All-purpose Symbolic Instruction Codeであるが，
1991年にMicrosoftがオブジェクト指向の考えを取り入れたVisual Basic 1.0 を発表し，
1994年にVBA =Visual Basic for Application を開発し 
Excel 5.0に搭載されたのが始まりである(\cite{W})．

参考文献について：
\cite{Y}はExcelの関数でどのようなことが出来るかを分かり易く解説したもの．
\cite{Kum}はこの資料であげた図~\ref{fun}, \ref{Fibonacci}, \ref{Pascal} のような見栄えのExcelを作るにはどうすればいいかを解説したもの．
\cite{Kud}はVBAについて非常に短いが分かり易く解説したもの．
\cite{T}はVBAについて非常に詳しく解説したものである.
Excelの機能は膨大なので，
辞書として\cite{EV}, \cite{EL}, \cite{EF} も必要となるだろう.
\clearpage

%\end{document}



%\documentclass[a4paper,12pt]{amsart}
%\usepackage[dvipdfmx]{graphicx,xcolor}%graphics
%\usepackage{float}
%\usepackage{amsmath}
%\usepackage{verbatim}
%\usepackage{latexsym}
%\usepackage{amssymb}
%\usepackage{ccicons}
%\usepackage{mflogo}
%%\usepackage{emath}
%\usepackage{MePoTeX}
%\usepackage[dvipdfmx]{hyperref}
%\hypersetup{% hyperrefオプションリスト
%   setpagesize=false,
%   bookmarksnumbered=true,%
%   bookmarksopen=true,%
%   colorlinks=true,%
%   linkcolor=red,
%   citecolor=blue,
%} 
%\input xy
%\xyoption{all}
%
%\newtheorem*{exQ}{レポート}
%\theoremstyle{remark}
%\newtheorem*{Q}{演習}
%
%\renewcommand{\figurename}{図}
%\renewcommand{\proofname}{証明}
%\DeclareMathOperator{\const}{const}
%\def\emptyline{\vspace{12pt}}
%
%\begin{document}
%
%\begin{flushright}
%   %2024/04/24
%   %2024/09/30
%   %2025/01/05
%   %2025/02/13
%   %2025/02/19
%   %2025/03/29
%   %2025/05/19
%   %2025/05/22
%   %2025/08/06
%   %2025/10/23
%   %2025/11/19
%   %2025/11/29
%   %2026/04/29
%   %2026/07/25
%   2026/08/26
%   
%   snow00two
%\end{flushright}
%\begin{center}
%   \large{\TeX\footnote{\TeX  は American Mathematical Society の登録商標だが，
%   このように明記する必要は無いらしい．} の導入}
%   \small{
%   \href{https://creativecommons.org/licenses/by-nc-nd/4.0/}{\ccbyncnd}}
%\end{center}

\chapter{TeXとMarkdown}

\section{\TeX}
\subsection{\TeX のインストール}
\TeX とは:  
\begin{itemize}
   \item 
   \TeX:  フリーで使える数式を含む文書作成言語(compiler). 
   
   D.~Knuth により開発された．
   
   \item
   \LaTeX:   L.~Lamport が\TeX を使いやすくしたマクロパッケージ．   
      数式・相互参照・表・作図・プレゼン・hyperlink などが使える
   
   \item
       高校の先生・高校生でも使ってる人がいるらしい．
\end{itemize}


\emptyline
英語表示のMacの場合のインストールパッケージ:  

\begin{itemize}
   \item 
      MacTeX:  mac 用 installing package （以下を全て含んでいる）
   
   \item 
      TeXLive:   \TeX の本体,  mac だと /usr/local/texlive にインストールされる．
   
   \item
      TeXShop:   \TeX の mac 用の editor （開発環境で, VSCodeなどでも使用可能）
   
   \item
      beamer:   プレゼンのための \LaTeX の documentclass
   
   \item
      \MP:   \TeX の描画ソフト  (font 作成ソフト \MF(\cite{K}) から派生した描画言語)
\end{itemize}

\begin{Q}
   \LaTeX のインストール（Mac の英語表示のみ）:  
   \begin{enumerate}
   \item
      google で MacTeX を検索
   
   \item
      MacTeXをダウンロード
   
   \item
      MacTeX.pkg をダブルクリックし，インストーラーの指示に従いインストール．
   
   \item
      texshop で 設定が必要かも知れない．
   \begin{figure}[htb]
      \centering
      \includegraphics[width=7cm]{~/Documents/GitHub/introInfo/appb/images/preferences}
      \caption{texshopの設定}\label{preferences}
   \end{figure}
   詳細は
   \href{https://texwiki.texjp.org}{texwiki}を参照
   
   
    \end{enumerate}
\end{Q}
      texlive のファイル群は/usr/local/texlive という不可視フォルダーに
      インストールされている．コマンドラインユーザーインターフェイスでは
      簡単にアクセスできるが，
    （Macの）グラフィカルユーザーインターフェースでは 
       Finder のメニューバーのgo のメニューの Go to Folder を選ぶと
       フォルダーの選択窓が開くので
      /usr/local/texlive  と打ち込むと移動できる．

      Windows などの場合には\cite{OK}などを参考にして
      インストールして下さい．

\subsection{\TeX のタイプセット}
次のソースファイル   
\begin{verbatim}
   \documentclass[11pt]{article}
   \usepackage{graphicx, xcolor}
   \usepackage{latexsym}
   \usepackage{amsmath}
   \usepackage{amssymb}
   \usepackage{amsthm}
   \renewcommand{\figurename}{図}
   \renewcommand{\proofname}{証明}
   \renewcommand{\tablename}{表}
   \DeclareMathOperator{\const}{const}
   \title{Hello world \TeX !}
   \begin{document}
   \maketitle
   
   因数分解 $x^3-1=(x-1)(x^2+x+1)$. 
   不定積分
   \[
      \int \sin x\, dx =-\cos x,\quad 
      \int \frac1xdx=\log |x|
   \]
   行列
   \begin{equation}\label{eqMat}
      B
      =
      \begin{pmatrix}
         1&0&0\\
         1&1&0\\
         1&1&1
      \end{pmatrix}
   \end{equation}
   \eqref{eqMat}のように引用できる．標準でない記号例えば$\const$
   も使うことが出来る(\verb/$const$/ とすると$const$になる)．
   \textcolor{red}{色も付けられる}．表も作れる:  
   \begin{table}[thb]
      \centering
      \begin{tabular}{l|llll}\noalign{\hrule height 1pt}
         &$\varepsilon$&$\alpha$&$\beta$&$\gamma$ \\ 
         \hline
         $\varepsilon$&$\varepsilon$&$\alpha$&$\beta$&$\gamma$ \\
         $\alpha$& $\alpha$&$\varepsilon$ &$\gamma$&$\beta$ \\
         $\beta$& $\beta$&$\gamma$&$\varepsilon$&$\alpha$ \\ 
         $\gamma$& $\gamma$ &$\beta$&  $\alpha$&$\varepsilon$\\
        \noalign{\hrule height 1pt}
      \end{tabular}\caption{Kline group}
   \end{table}
   
   \noindent
   以下のように文献表も付けられる(\cite{K})．
   \begin{thebibliography}{9}
   \bibitem{K} Donald E. Knuth, \textit{ The \TeX\, book}, 
      Addison-Wesley, 1984.
   \end{thebibliography}
   \end{document}
 \end{verbatim}

\begin{figure}[htb]
   \centering
   \includegraphics[trim=0 30  0 40, clip, width=12cm]{~/Documents/GitHub/introInfo/appb/images/ScreenCapture1.pdf}
   \caption{Hello world \TeX !}\label{compile}
\end{figure}

\noindent
を helloWorld.tex\footnote{どんなコンピュータ言語でも最初に書くプログラムの名前を helloWorld のようなものにするのがプログラマーの伝統である．また，どうでも良い名前の場合 hoge とか fuga のような名前を使うのも伝統である．} として保存し（読み込みたい画像とかのファイルがあればそのディレクトリに置き），
helloWorld.tex を TeXShop.app で開き， typeset をクリックするか， cmd T とタイプするとコンパイルできる(図 \ref{compile})．

\begin{Q}
   上のソースファイルをコンパイルしてpdf を表示してください．
\end{Q}
\clearpage

\begin{verbatim}
　\begin{document}
\end{verbatim}
の前を preamble (序文)という．
取りあえず preamble はコピーして使い，簡単な数式を \cite{TeXShop} などを調べながら書いてみることをお奨めする．ひとことコメントすると
\begin{verbatim}
　\DeclareMathOperator{\const}{const}
\end{verbatim}
が標準にない記号$\const$を定義しているところである．
あと，\TeX の機能と云うよりTeXShop ( Editor ) の機能だが，
\begin{itemize}
\item
cmd キーを押しながらクリックすると対応するpdf の場所とソースファイルの場所を行き来することができる．

\item 
コンパイルしてエラーが出て止まったとき，console の上部の Goto Error をクリックすると，ソースファイルの該当部分に飛ぶことができる．

\item
ソースファイルでEsc を押すと補間機能を使うことが出来る．
例えば\verb/\begin{eq/まで打ってEsc を押すと
\begin{verbatim}
　\begin{equation}
　・
　\end{equation}・
\end{verbatim}
と補完される．

\end{itemize}

また，
\emptyline
 emath という大熊一弘氏の開発した日本の小中高数学のプリント作成用のlataxの style file があり，教科書の数式に近い出力を得ることが出来る．
\href{http://emath.s40.xrea.com/ }{ここから} 
ファイルを取り寄せ，

\noindent
/usr/local/texlive/texmf-local/tex/latex/local/にコピーし, 
terminal から 

\verb/  % /sudo texhash

\noindent
としてpathを通す． ただし\verb/%/ はプロンプトで，タイプしていけない．
preamble に
\begin{verbatim}
   \usepackage{emath}
\end{verbatim}
を追加すると使えるようになる．詳しくはemath のマニュアルを参照．

\subsection{\TeX の図式作成ソフト}  
\[
   \xymatrix{
   S^2(1) \ar[d]_{\mu'} \ar[r]
    & \mathbf R^3\setminus \{\mathbf 0\} \ar[d]^{\mu} \\
   S^2(1)/ \sim  \ar[r]_{\pi} & \mathbf P^2 
   }
\]
のような図式を描く, texlive に含まれるソフトの一つが XY-pic である．まず，
\begin{verbatim}
   \input xy
   \xyoption{all}
\end{verbatim}
をpreamble に付け加え，本文中に
\begin{verbatim}
   \[
      \xymatrix{
         S^2(1) \ar[d]_{\mu'} \ar[r]
          & \mathbf R^3\setminus \{\mathbf 0\} \ar[d]^{\mu} \\
        S^2(1)/ \sim  \ar[r]_{\pi} & \mathbf P^2 
      }
   \]
\end{verbatim}
を入れると，上の図式を描くことができる．
\verb/\xymatrix{...}/ の中でマトリックスを描く要領で成分の数式の記号を並べている．\verb/&/は列の分割を表し，\verb/\\/は改行である．
\verb/ \ar[d]/は下への矢印，
\verb/ \ar[r]/ は右への矢印を描く指令で，それに続く\verb/ _{\mu'}/
は下付き，\verb/^{\mu}/は上付きの矢印の修飾を描く命令である．
   
   
%もう少し複雑な可換図式
%\[
%   \xymatrix{
%    0\ar[r]
%    &H^0(X,\mathfrak S') \ar[r] \ar[d]
%    &H^0(X,\mathfrak S) \ar[r] \ar[d]
%    &H^0(X,\mathfrak S'') \ar[r]^{\delta^{*}} \ar[d]
%    &H^1(X,\mathfrak S')\ar[r] \ar[d]
%    &\\
%    0\ar[r]
%    &H^0(X,\mathfrak T') \ar[r]
%    &H^0(X,\mathfrak T) \ar[r]
%    &H^0(X,\mathfrak T'') \ar[r]^{\delta^{*}} 
%    &H^1(X,\mathfrak T')\ar[r] 
%    &\\
%    }
%\]
%も
%\begin{verbatim}
%   \[
%   \xymatrix{
%       0\ar[r]
%       &H^0(X,\mathfrak S') \ar[r] \ar[d]
%       &H^0(X,\mathfrak S) \ar[r] \ar[d]
%       &H^0(X,\mathfrak S'') \ar[r]^{\delta^{*}} \ar[d]
%       &H^1(X,\mathfrak S')\ar[r] \ar[d]
%       &\\
%       0\ar[r]
%       &H^0(X,\mathfrak T') \ar[r]
%       &H^0(X,\mathfrak T) \ar[r]
%       &H^0(X,\mathfrak T'') \ar[r]^{\delta^{*}} 
%       &H^1(X,\mathfrak T')\ar[r] 
%       &\\
%    }
%   \]
%\end{verbatim}
%で描くことができる(詳しいことは\cite{texlive}参照)．


\subsection{\TeX の描画ソフト}   
\MP のコンパイルは，ソースファイル figure.mp をtexshopで開き，ファイルのウインドウの上部にある
ファイルタイプのセレクターをLaTeX からMetaPost に変更し，タイプセットする(図 \ref{texhopMetapost})と，figure.1, figure.2のようなpostscript ファイルが生成される．あるいは，ファイルのあるディレクトリで terminal から

   \verb/% /mpost figure.mp

\begin{figure}[htb]
   \centering
   \includegraphics[width=12cm]{~/Documents/GitHub/introInfo/appb/images/texhopMetapost.pdf}
   \caption{\MP のコンパイル}\label{texhopMetapost}
\end{figure}

\noindent
でコンパイル出来る．
\newpage

%\begin{figure}[htb]
%   \centering
%   \includegraphics{~/Documents/GitHub/introInfo/appb/images/figure.1}
%   \caption{\MP の描画例}\label{~/Documents/GitHub/introInfo/appb/images/figure.1}
%\end{figure}
\begin{figure}[htb]
　\centering
　\begin{minipage}{0.4\columnwidth}
　　\centering
　　\includegraphics{~/Documents/GitHub/introInfo/appb/images/figure.1}
   \caption{\MP の描画例}\label{figure.1}
　\end{minipage} 
　\begin{minipage}{0.4\columnwidth}
　　\centering
　　\includegraphics[trim=90 10  130 30, clip]{~/Documents/GitHub/introInfo/appb/images/figure.2}
   \caption{\MP の描画例}\label{figure.2}
　\end{minipage} 
\end{figure}
\noindent
のような絵を描く\MP のソースファイルの例を次にあげる:  

\begin{verbatim}
   prologues := 3;
   verbatimtex
   %&latex
   \documentclass{article}
   \usepackage{verbatim}
   \usepackage{latexsym}
   \usepackage{amssymb}
   \usepackage{amsmath}
   \begin{document}
   etex;
   
   numeric a, b;
   a=10pt;
   b=6pt;
   
   beginfig(1);
      z0=(5a,3a);
      z1=z0-(3a,0);
      z2=z0+(7a,0);
      draw z1--z2; %line \ell
      label.top(btex $\ell$ etex,z2); 
      z5=z0+(7a*cosd(60), 7a*sind(60));
      z6=z0-(3a*cosd(60), 3a*sind(60));
      draw z5--z6;%line m
      label.top(btex $m$ etex,z5); 
   
      for j=-2 upto 10:             
         z[100+j]=z0+6a*(cosd(10j),sind(10j));
      endfor;
   
      draw 
         for i=-2 upto  9:
            z[100+i]..
         endfor
         z110 dashed evenly;
   
      z150=whatever[z0,z2]=whatever[z98,z102];
      z160=whatever[z0,z5]=whatever[z102,z110];
      z151=z150-(b,0);
      z152=z151-(0,b);
      z153=z152+(b,0);
      z161=z160-(b*cosd(60),b*sind(60));
      z162=z161+(b*cosd(150),b*sind(150));
      z163=z162+(b*cosd(60),b*sind(60));
      
      draw z98--z102--z110;
      draw z151--z152--z153;
      draw z161--z162--z163;
      
      label.bot(btex $\mathbf 0$ etex,z0);
      label.bot(btex $\mathbf x$ etex,z98);
   
      pickup pencircle scaled 1pt;
      drawarrow z0--z98;
      drawarrow z0--z102;
      drawarrow z0--z110;
   endfig;
   
   beginfig(2);
      def voro(expr p, q, r) = 
         begingroup
             save a, b, c, u, v, w, e, f, g, h; 
             pair a, b, c, u, v, w;
             numeric e, f, g, h;
             a = 0.5[q,r];
             b = 0.5[r,p];
             c = 0.5[p,q];
             (e,f)= r-q ;
             (g,h) = p- r;
             u = (f, -e);
             v = (h, -g);
             w = whatever[a + u, a - u] = whatever[b + v, b - v] ;
             draw a--w--b;
             draw w--c;
         endgroup
      enddef;
   
      z0=(-5a,-5a);
      for i=1 upto 6:
         for j=1 upto 6:
         	z[10i +j]= z0 + (5i*a, 0) +5j*(a/2, sqrt(3)*a/2) 
            	+(a*cosd(30(i+j)+i*j/2), a*sind(30(i+j) + i*j/2) );
         endfor;
      endfor;
      
      z100 = z32+ (a,1.4a);
      dotlabel(btex etex, z100);
      draw z100--z32;
      draw z100--z33 dashed evenly;
      draw z100--z23 dashed evenly;
      draw z100--z42 dashed evenly;
   
      for i=1 upto 6:
         for j=1 upto 4:
         	dotlabel(btex etex, z[10i +j]);
         endfor;
      endfor;
      for i=1 upto 5:
         for j=1 upto 3:
         	voro(z[10i +j], z[10(i+1)+j], z[10i +j+1]);
         	voro(z[10i +j+1], z[10(i+1) +j+1], z[10(i+1)+j]);
         endfor;
      endfor; 
   endfig;
   
   verbatimtex
   \end{document}
   etex;
   end
\end{verbatim}
をfigure.mp と言う名前で保存し，コンパイルすると
figure.1, figure.2 というファイルが作られる．
\MP のプログラミング言語としての仕様
(詳しくは\cite{texlive}):  

\begin{itemize}
   \item
   beginfig(1); からendfig; までが figure.1の描画内容である． 
   figure.2も同様．
   これ以外の最初の部分と最初の部分はコピペすることをお奨めする．
   
   \item
   文末には ; が必要．\verb/%/ の後がコメント文
   
   \item
   点を表すのに, \verb/z10/ のような \verb/z/に数字を並べた表記を使う．
   インデックスを変数として扱うときは\verb/z[i]/のような記法を使う．
   
   \item
   点\verb/z1/から\verb/z2/への直線を書くのは
   \begin{verbatim}
   draw z1--z2;
   \end{verbatim}
   点\verb/z1/から\verb/z2/への矢印を書くのは
   \begin{verbatim}
   drawarrow z1--z2;
   \end{verbatim}
   \verb/z1/, \verb/z2/, \verb/z3/ を補間するベジェ曲線は
   \begin{verbatim}
   draw z1..z2..z3;
   \end{verbatim}
   
   \item
   \MP 特有の機能として
   \begin{verbatim}
   z150=whatever[z0,z2]=whatever[z98,z102];
   \end{verbatim}
   のように直線と直線の交点を求めることが出来る．
   様々な初等関数も用意されている．
   関数は \verb/sind(10)/ のように呼び出される．
   
   \item
   for 文で点を定義したり，曲線を描くことも出来る．
   
   \item
   図への数式や文字のラベルでは \LaTeX と同じものを使える．
   例えば
   \begin{verbatim}
   label.bot(btex $\mathbf 0$ etex,z0);
   \end{verbatim}
   は点\verb/z0/の真下の部分に$\mathbf 0$ ( \verb/$\mathbf 0$/ ) を書くことを
   命令するコマンドである．%(この文や関連した文のソースファイルをMoodleのホームページに置いておくので参考にせよ．)
   \item
   beginfig(2)に関数の使い方の例をあげたのでぜひ読み解いてほしい．
   この例ではdef からenddef の中身voloが三つの入力変数をもつ関数である．
   $=$の後がその定義である．定義の中で使う変数をまず保存して，
   その後変数の宣言を書く．
   次に二点の中点をもとめ，二点の作る辺に直行するベクトルをもとめ
   その中点に直交する線分を求めるている．
   もう一つ二点についても同様に中点に直交する線分を求め
   それらの交点として三点に外接する円の中心をもとめている．
   そこから各中点まで線分を3っつ書いている．
   
   あとは正6角形格子の点を少し振動させ，
   それらの点についてのボロノイ分割をvolo()を使って書いている．
   \item
      ここではこれ以上述べないがもちろん if 文も使える．
\end{itemize}

\begin{Q}
   上の\MP のソースファイルをfigure.mp と言う名前で保存し，コンパイルし，
   上で作ったtex のファイルに取込み pdf に表示してください．
\end{Q}


Ti\textit{k}Z と言う style file もあり，それを使うと直接 \LaTeX の中で\MP に似たコマンドで描画できる．また metafun という\MP を現代風に拡張したstyle file もある．

MePoTeXという（
\href{http://minamo.my.coocan.jp/metapost/mepotex.html}{ここから入手可能} ）
\MP を拡張する style file もあり，(emath と同じようにインストールし path を通せば） \LaTeX の中で使える． preamble に
\begin{verbatim}
   \usepackage{MePoTeX}
\end{verbatim}
を追加し，本文に

\begin{verbatim}
   \begin{MPpic}<5mm>(7|2,4|2)
   \mptXaxis[l]<1mm,0mm>{$x$}
   \mptYaxis[b]<0mm,1mm>{$y$}
   \sendMP{
      xdraw() kansu(sin(t))
      (-1,6.5,20);
      xdraw(red) kansu(t)
      (-1,3.5,20);
   }
   \mptLabel{origin}[tr]
   <0mm,4mm>{O}
   \end{MPpic}
\end{verbatim}
を挿入すると

\begin{center}
   \begin{MPpic}<5mm>(7|2,4|2)
   \mptXaxis[l]<1mm,0mm>{$x$}
   \mptYaxis[b]<0mm,1mm>{$y$}
   \sendMP{
      xdraw() kansu(sin(t))
      (-1,6.5,20);
      xdraw(red) kansu(t)
      (-1,3.5,20);
   }
   \mptLabel{origin}[tr]
      <0mm,4mm>{O}
   \end{MPpic}
\end{center}
のような図が描ける． 
\emptyline
\begin{exQ}
   自分のセミナーの内容を説明する\LaTeX のファイルを作成し，コンパイルせよ
   （目標は下の作業を含めてpdfで5〜10ページほど）．
   それに表・図式・参考文献を付け加えよ．
   \MP で図を描き，それをその説明の \LaTeX のpdfの中に表示せよ．
   
\end{exQ}

%Moodle で講義ノートの\TeX のソースファイルを公開するので参考にして下さい．また，
%以下の参考文献も参照%してください．
%
%\begin{thebibliography}{9}
%   \bibitem{TeXShop}
%   TeXShop のメニューバーの Helpメニューにある説明文書と動画
%   
%   \bibitem{texlive}
%   /usr/local/texlive/2024/texmf-dist/doc/latex の中のドキュメントファイル:  
%    User's guide for amsmath package, Using the amsthm package, 
%    XY-pic User's Guide, XY-pic Reference Manual,
%    The \MP\, system, A users manual for \MP など
%   
%   \bibitem{OK}
%   奥村晴彦, 黒木裕介, 
%   ［改訂第9版］\LaTeX  美文書作成入門, 技術評論社, 2023.
%   
%   \bibitem{LCs}
%   {\em The  \LaTeX \, Companions Third Revised Boxed Set}, Addison-Wesley, 2007.
%   
%   \bibitem{K} 
%   Donald E.~Knuth, {\em The \MF book},  Addison-Wesley,1986.
%\end{thebibliography}

\section{Markdown}
\TeX とは直接関係はないが，最近よく使われるMarkdown言語についても紹介したい．
これはWeb page を編集するための言語 html (=hyper text markup language) 
を簡易的に表すために，つまり，markdown を目指して，
2004年ジョン・グルーバー(J.~Gruber)によって開発された言語である．
VSCodeで IntroToMarkdown.md というファイルを作り，
以下のソースをコピーし，cmd+shift+Vとタイプすると（
またはウインドウの右上にあるノートに虫眼鏡のアイコンをクリックすると）
図~\ref{preview}, \ref{preview1}のような html のプレビュー画面が現れる．
%\newpage

%remove tab 
\begin{verbatim}
   <!-- Introduction to markdown created by snow00two -->
   
   # Introduction to Markdown
   
   ## Markdown
   
   - list1
   - list2
     - sublist1
     - sublist2
   - list3
   
    1. list1
    2. list2
    3. list3
   
   Separator:
   
   ---
   
   *italic* ,  **bold** , ***italic bold***
   
   > quotation
    >> second quotation
   
   Programming Language `JavaScript` list:
   
   ```javascript
   function hoge () {
     console.log("problem");
   }
   ```
   
   Table:
   
   | Data | A | B |
   | ---: |: ---: |: ---: |
   | 0 | 1 | 2 |
   
   Mathematics:
   $\sqrt2$, $\pi$, $\frac1{x + 1}$
   
   Hyperlink to: [amazon](https://www.amazon.co.jp/amazonprime)
   
   Image file: ![image file](dodecahedron.jpg)
     
\end{verbatim}

文法についてはこのファイルを見れば直ぐに分かるだろう．
最近の技術系ホームページ (例えば GitHub) は
Markdownで作られていることが多い．
また，生成AIのファインチューニングという訓練でMarkdownが使われていることはよく知られていて，ChatGPTへのプロンプトの入力もMarkdownを使うことが推薦されている(\cite{LLM}, \cite{Hi}). 
最後にmarkdown から\LaTeX を作ることも出来る（とは言っても数式等は\LaTeX で作る必要がある）．

\begin{figure}[htb]
   \centering
   \includegraphics[width =12cm, trim=0 80  0 0, clip]{~/Documents/GitHub/introInfo/appb/images/IntroToMarkdown.pdf}
   \caption{VSCode での IntroToMarkdown.md のプレビュー1}
   \label{preview}
\end{figure}

\begin{figure}[htb]
   \centering
   \includegraphics[width =12cm, trim=0  0 0 125, clip]{~/Documents/GitHub/introInfo/appb/images/IntroToMarkdown1.pdf}
   \caption{VSCode での IntroToMarkdown.md のプレビュー2}
   \label{preview1}
\end{figure}

\begin{Q}
\begin{enumerate}
   \item
   VSCode で上のソースファイルを参考にして，1〜2ページ程度で
   markdown を説明する markdown の文書を作って
   プレビューして下さい．
   \item
   そのソースファイルを，生成AIを使ってlatex のソースファイルに変換してください(その中でコマンドの実行を頼まれるかも知れません)．
   
   \item
   その過程で生成AIが何をしたのか分析してください．
   
   \item
   その分析に基づき，自分一人でもう一度そのlatexの書類をmarkdownに変換して下さい．
\end{enumerate}
\end{Q}

%\begin{thebibliography}{9}
%   \bibitem{LLM}
%   Berryman, Zigler, LLMのプロンプトエンジニアリング, オライリージャパン, 2025.
%
%   \bibitem{H}
%      平田 賀一, 
%      ITエンジニアのためのMarkdown実践入門 ── 生成AI時代の"伝わる"書き方, 
%      技術評論社, 2026.
%
%\end{thebibliography}
%
%\end{document}


%\documentclass[a4paper,12pt]{amsart}
%\usepackage[dvipdfmx]{graphicx,xcolor,hyperref}%graphics
%\usepackage{float}
%\usepackage{amsmath}
%\usepackage{verbatim}
%\usepackage{latexsym}
%\usepackage{amssymb}
%\usepackage{ccicons}
%%\usepackage{emath}
%%\usepackage{maplestd2e}
%%\usepackage{fancybox}%enclose by box
%\hypersetup{% hyperrefオプションリスト
%setpagesize=false,
% bookmarksnumbered=true,
% bookmarksopen=true,
% colorlinks=true,
%linkcolor=red,
%citecolor=blue,
%} 
%
%\renewcommand{\figurename}{図}
%\renewcommand{\proofname}{証明}
%\renewcommand{\tablename}{表}
%\newtheorem*{PP}{授業の目標}
%\newtheorem*{Thm}{定理}
%\newtheorem*{Prop}{命題}
%\newtheorem*{Q}{考えたい課題}
%\newtheorem*{Quiz}{問題}
%\newtheorem*{exQ}{レポート}
%
%\theoremstyle{definition}
%\newtheorem*{Def}{定義}
%
%\theoremstyle{remark}
%\newtheorem*{ex}{例}
%
%\def\emptyline{\vspace{12pt}}
%\DeclareMathOperator{\const}{const}
%\DeclareMathOperator{\Ker}{Ker}
%\DeclareMathOperator{\Image}{Im}
%\DeclareMathOperator{\id}{id}
%\DeclareMathOperator{\Is}{Is}
%\DeclareMathOperator{\stab}{stab}
%\DeclareMathOperator{\tr}{tr}
%\DeclareMathOperator{\vol}{vol}
%
%\begin{document}
%\begin{flushright}
%   2026/06/23
%   
%   九大・数理 大津幸男
%\end{flushright}
%
%\begin{center}
%   {\Large 情報科教育法II---08}
%   \href{https://creativecommons.org/licenses/by-nc-nd/4.0/}{\ccbyncnd}
%   \par{JavaScript と p5.js}
%\end{center}

\chapter{JavaScript と p5.js}

\href{https://snow00two.github.io/mySite/}{私のポートフォリオ}
に私がJavaScript と p5.jsで作ったアプリが置いてある．
その幾つかについて説明する．

%\section{JavaScript と p5.js}
\section{ p5.js の2Dアプリ}
\subsection{p5.js}
p5.js は Processing というJava用の描画プログラミング言語を Javascript に移植したものである．以下のウェブ上のエディターで実行可能である．

\href{https://editor.p5js.org}{https://editor.p5js.org} （以下この色はハイパーリンク）

\begin{figure}[htb]
   \centering
   \includegraphics[width=12cm]{~/Documents/GitHub/introInfo/appc/images/ScreenCapture3.pdf}
   \caption{ウェブ上の p5.js エディター}\label{}
\end{figure}

\noindent
公開した sketch.js を sketch.js の部分にコピーし，
（ファイル octahedron.stl をアップロードし）
左上の三角のボタンを押すと右画面に描画が始まる．
もちろん隣りの四角のボタンで止まる．
%プログラムのソースファイルに是非適宜手を入れながら動かして見て欲しい．
ただし .js はJavaScriptのファイルであることを示す拡張子 ( extension ) である．


また以下のようにすると自分のパソコン上で動かせる．
 p5.js というライブラリーを\href{https://p5js.org/download/}{ダウンロード}するとp5というディレクトリが作られる．
 公開されたsketch.jsをそのディレクトリーの sketch.js と置き換え，
フリーのエディターVSCode \href{https://code.visualstudio.com}{Visual Studio Code}でそのディレクトリp5を開き，
 Live Server の拡張をVSCodeにインストールして
Go Liveのボタンを押すと，chrome で再生を始める．


\subsection{一般のJavaScriptの動かし方}
JavaScriptの処理の仕方は幾つかある．
JavaScriptはブラウザーに組み込まれているので，まず
\begin{verbatim}

   <!DOCTYPE html>
   <html>
      <head>
          <meta charset="UTF-8" />
          <title>Hello JavaScript</title>
      </head>
      <body>
         <script type="text/javascript" src="test.js"></script>
      </body>
   </html>
   
\end{verbatim}
のようなhtml のファイルを hoge.html として保存し，次に
\begin{verbatim}
   window.alert('hello world!');
\end{verbatim}
という行のみからなるファイルを test.js として hoge.htmlと同じディレクトリに保存する．
ブラウザー Chrome の
メニューバーのFileのメニューOpen Files... からhoge.htmlを開くと

\begin{figure}[htb]
   \centering
   \includegraphics[width=11cm, trim=0 1000  0 0, clip]{~/Documents/GitHub/introInfo/appc/images/ScreenCapture1.pdf}
   \caption{}\label{alart}
\end{figure}

\begin{figure}[htb]
   \centering
   \includegraphics[width=11cm, trim=0 800  0 0, clip]{~/Documents/GitHub/introInfo/appc/images/ScreenCapture2.pdf}
   \caption{}\label{console}
\end{figure}

\noindent
図 \ref{alart}のような結果が得られる．あるいは
\begin{verbatim}
   console.log('hello world!');
\end{verbatim}
をtest.jsとして保存し，またhoge.htmlを開くと
Fileのメニュー View から Developer の Developer Tools を確認すると
図 \ref{console}のような結果が得られる．
%\newpage

他にも Terminal でtest.jsのあるディレクトリに移動し
\begin{verbatim}
   % osascript test.js
\end{verbatim}
でTerminalに結果が表示される．
ただし，これはMacのみで，逆に言うとMacのスクリプト言語としてJavaScriptを使うことが出来る．
一般的には Node.js を\href{https://nodejs.org/en}{ダウンロード}
しインストールすると
\begin{verbatim}
   % node test.js
\end{verbatim}
でTerminalに結果が表示される．したがってMac を直接コントロールする必要が無いときは osascriptより node.js を使うことをお奨めしたい．
%Node.js + Electron + Chromium で macOS, Windows, スマートホンに対応したデスクトップアプリケーション(One page application)を一度に作ることが出来る (例 VScode 図 \ref{vscode})．
%
%\begin{figure}[htb]
%\centering
%\includegraphics[width=7cm, trim=0 40  0 30, clip]{~/Documents/GitHub/introInfo/appc/images/ScreenCapture4.pdf}
%\caption{}\label{vscode}
%\end{figure}
%
%CGに興味を持った人には\cite{CG}, 遠近法の歴史については\cite{D},
%具体的な遠近法の絵の描き方については\cite{3DT}, \cite{RS}, \cite{RSII} などをお奨めする．


\subsection{JavaScriptの言語仕様}

\begin{itemize}
\item
JavaScript には変数の型はなく，文の終りで ; が必要

\item
let は変数の宣言，以前は var も使われていた．

\item
const は定数の宣言

\item
\verb////はコメント文でこれ以降の文は無視される．

\noindent
/* */はブロックコメント文で途中の文は無視される．
\end{itemize}


よく使う制御構文：
for 文の例：
\begin{verbatim}
   let I = 5;
   for (let i=0; i < I; i++) {
      console.log(i);
   }
   
\end{verbatim}

if 文の例：
\begin{verbatim}

   let m = 80;
   if ( m<60 ){
      console.log('F');
   } else if ( m<70  ) {
      console.log('C');
   } else if ( m<80  ) {
      console.log('B');
   } else if ( m<90  ) {
      console.log('A');
   } else {
      console.log('S');
   }

\end{verbatim}


function hoge(var1,\dots,var3) は関数hogeの宣言で，続く$\{ \}$の中に関数の定義を書く．var1,\dots,var3が関数の変数で定義の中で使える．
関数の中で定義された変数はローカル変数で，関数外で使われる変数はユニバーサルな変数になる．

\subsection{2Dアプリのソース}
以下がSierpinski gasketを描くアプリのソースファイル sketch.js の内容である．

\begin{figure}[htb]
	\centering
	\includegraphics[width=8cm]{~/Documents/GitHub/introInfo/appc/images/ScreenCapture7.pdf}
	\caption{Sierpinski gasket}\label{Sierpinski gasket}
\end{figure}

\begin{verbatim}
/* Sierpinski gasket etc., abbreviated version
 * 2026/06/22, by snow00two
 * \href{https://creativecommons.org/licenses/by-nc-nd/4.0/}{\ccbyncsa}
 */
/* canvas setting */
const WIDTH_CANVAS = 720 * 3/2  ; //=1080
const HEIGHT_CANVAS = 405 * 3/2 ; //=607.5
const CENTER_X = WIDTH_CANVAS/2 ;
const CENTER_Y = HEIGHT_CANVAS/2 ;

/* color setting */
const DARK_COLOR = [150,30,50];
const WHITE_COLOR = 255;
let colorChoice ;
let choiceC ;/* variable for color scheme */
let backColor ;
let frontColor ;

/* Sierpinski gasket setting, abbr. 3 */
const BASE_SEGMENT_3 = 500;
let height3;
let radius3;
let shift3Y;
let fixed3X;
let fixed3Y;

function setup() {
  /* Sierpinski gasket additional setting */
  height3 = sqrt(3)* BASE_SEGMENT_3 / 2;/* sqrt(3) is defined here */
  radius3 = 2 * height3 /3 ;/* the radius of circumscribed circle  */
  fixed3X = ( WIDTH_CANVAS - BASE_SEGMENT_3 ) / 2 ;
  fixed3Y = ( HEIGHT_CANVAS - height3  ) / 2 ;
  shift3Y= ( HEIGHT_CANVAS - height3 ) / 3 ;

  colorChoice = createSelect();/* object setting color scheme */
  colorChoice.option('normal color','normal');
  colorChoice.option('reverse color','reverse');
  colorChoice.selected('normal');
  colorChoice.position(30, 8);
  colorChoice.changed( resetBackground ) ;

  frameRate(0.6);
 
  /* create canvas */
  createCanvas(WIDTH_CANVAS, HEIGHT_CANVAS);
  background(DARK_COLOR);
}

let n = 0; //clock parameter
/* main function controlling this application */
function draw() {
  let choiceC = colorChoice.value();
  if (choiceC =='normal' ){
    backColor = DARK_COLOR;
    frontColor = WHITE_COLOR;
  } else {
    backColor = WHITE_COLOR;
    frontColor = DARK_COLOR;
  }
  background(backColor);

  translate(CENTER_X, CENTER_Y - shift3Y);
  fill(frontColor);
  iterateTriangleS(n);

  let I = 9 ; //Number of iterations 
  if( n < I ){
    n++ ;
  } else {
    background(backColor);
    n = 0 ;
  } 
}

/* recursive function creating Sierpinski gasket */
function iterateTriangleS(num){
  push();
    if (num == 0){
      fill(frontColor);
      noStroke();
      triangle(-BASE_SEGMENT_3 / 2, -radius3 /2, 
      BASE_SEGMENT_3 /2,　-radius3 /2, 0, radius3 );
  //Here we erase the center reverse triangle 
    } else {
      scale( 1 / 2 );
      fill(frontColor);
      noStroke();
      for (let l=0;l < 3;l++){
        translate(0, radius3 );
        iterateTriangleS(num - 1);
        translate(0, -radius3 );
        rotate(2*PI/3);
      }
    }   
  pop();
} 

function resetBackground () { 
  background(backColor);
  n = 0;
}\end{verbatim}


\subsection{p5.js の2Dモードの説明とソースファイルの解説}
p5.js の2次元モードでは$y$軸の方向が通常とは逆で下向きになっている．
これはテレビなどの表示では，一次元で送られてくるデータを図\ref{~/Documents/GitHub/introInfo/appc/images/figure.2}のように上から下に折り返しながら表示していたからである．
したがって，円周上で角度が増加するプログラムでは\textcolor{red}{時計回り}に回転する．
%以下とp5.js のプログラムでは動く円は時計回りに動くようにパラメータをとることにする．
もちろん，反時計回りにするためにはパラメータの符号を代えれば良い．
\begin{figure}[htb]
   \centering
   \begin{minipage}{0.4\columnwidth}
      \centering
      \includegraphics{~/Documents/GitHub/introInfo/appc/images/figure1.2}
      \caption{}\label{~/Documents/GitHub/introInfo/appc/images/figure.2}
   \end{minipage} 
   \begin{minipage}{0.4\columnwidth}
      \centering
      \includegraphics{~/Documents/GitHub/introInfo/appc/images/figure1.1}
      \caption{}\label{~/Documents/GitHub/introInfo/appc/images/figure.1}
\end{minipage} 
\end{figure}



\emptyline
関数setup()の中でウィンドウ等を初期化している．
createCanvas( widthCanvas, heightCanvas) で描画領域を指定している．
関数draw()が呼び出されると変数 n を1フレーム毎に1づつ増やしながら，
その都度フレームを描画し続ける．いじってみるとだんだん理解が深まると思う．
本質的な部分は帰納関数 iterateTriangleS() である．

\emptyline
p5.js 特有の関数を説明したい：
\begin{itemize}
   \item
      translate($a, b$); $(a, b)$だけ平行移動し，その点を$(0,0)$として描画
   
   \item
      rotate(theta); 角度 theta だけ直交座標系を回転し，
      それを新たな直交座標系として描画
   
   \item
      triangle($a, b, a', b', a'' b''$); $(a, b)$, $(a', b')$, $(a'', b'')$を頂点とする
      三角形を描画
   
   \item
      background($[a,b,c]$); 画面をRGB $[a,b,c]$で埋める．
      ただし$a, b, c$は$0$から$255=2^8-1$ の間の整数．
      
   \item 
      fill($[a,b,c]$); 次の図形をRGB$[a,b,c]$ で埋める．
   
   \item
      push(), pop() はその間に変えられた設定を元に戻すためのスタックに
      書き込むための関数  
\end{itemize}

関数
iterateTriangleS(num)がこのプログラムの本体である．
push() でそれまで呼び込まれた設定を保存し，pop()で元の設定に戻っている．
$y$軸が逆だと分かりずらいので，$y$軸を通常のようにとり，
回転方向も反時計回りだと思って正しくプログラムすると（図 \ref{trainagle}），
出力は$y$軸が逆になって描画される．
三角形の重心が原点にあるとする．
num $=0$ のとき fill(frontColor)で次の図をfrontColor で満たすことを， noStroke() で輪郭を描かないことを指示し，
\begin{verbatim}
   triangle(-baseSegment3 / 2, -radius3 /2, 
      baseSegment3 /2,　-radius3 /2, 0, radius3 );
\end{verbatim}
で，辺の長さが baseSegment3 でその底辺が$x$軸に平行な正三角形を書いている．
このとき, radius3は正三角形の外接円の半径である．
\begin{figure}[htb]
   \centering
   \includegraphics{~/Documents/GitHub/introInfo/appc/images/figure1.3}
   \caption{}\label{trainagle}
\end{figure}

num $= 1$のときは
scale  $1/2$ にして$(0,0)$から $(0, $radius3$)$ 進むので
実際は$(0,0)$から$ (0, $radius3$/2)$ 進んで三角形を描く．
その後 $(0, -$radius3$)$ 進み $(0,0)$ に戻り$2\pi/3$回転する．
これを三回繰り返し，3つの三角形を描く．

num$ >1$のときはこれを繰り返すことでSierpinski gasketを書くことが出来る．



\begin{exQ}
   Sierpinski Carpet(図~\ref{Sierpinski Carpet}) を描くようにプログラムを変更せよ．
\end{exQ}
\begin{exQ}
   Koch flake(図~\ref{Koch flake}) を描くようにプログラムを変更せよ．
\end{exQ}

 \begin{figure}[htb]
   \centering
　\begin{minipage}{0.4\columnwidth}
　　\centering   　　
      \includegraphics[width=4cm]{~/Documents/GitHub/introInfo/appc/images/ScreenCapture5.pdf}
      \caption{S.~Carpet}\label{Sierpinski Carpet}
　\end{minipage} 
　\begin{minipage}{0.4\columnwidth}
    　　\centering
   　　\includegraphics[width=4cm, trim =0 0 0 0]{~/Documents/GitHub/introInfo/appc/images/ScreenCapture6.pdf}
   　　\caption{Koch flake}\label{Koch flake}
  　\end{minipage} 
 \end{figure}

\section{ p5.js の 3D-アプリ}
\subsection{3Dアプリのソース}
$3$次元の描画を行うことも出来る．createCanvasを

　createCanvas(widthCanvas, heightCanvas, WEBGL);

\noindent
に替えれば良い．ただし，WEBGLとはブラウザ上でCGを描画するためのAPIである．
以下がソースファイル sketch1.js の本質的な内容である．

\begin{verbatim}
   /* Cubes landing before the Great Pyramids abbreviated version
    *  2026/06/17, by snow00two,
    * \href{https://creativecommons.org/licenses/by-nc-nd/4.0/}{\ccbyncsa}
    */
   const WIDTH_CANVAS = 1080  ; // 720x3/2 
   const HEIGHT_CANVAS = 607.5 ; // 405x3/2
   const BACK_COLOR = [150, 200, 250] ;//[200, 220, 200]
   
   let octahedron ; 
   function preload(){
     octahedron = loadModel('octahedron.stl'); 
     //https://www.printables.com/model/411768-octahedron-freecad/files
   }
   
   function setup(){//canvas : window of the projection coordinates
     createCanvas(WIDTH_CANVAS, HEIGHT_CANVAS, WEBGL);
     background(BACK_COLOR);
   }
   
   let i = 0;//basic parameter for frames
   
   function draw(){
     let I = 9425; //change the priod condition to 3x2Pi/0.002=9424.7... 
   
     if (i<I) {
     } else {
       i = 0;// reset basic parameter
     }
   
     // define the camera coordinates from a camera on the helicopter
     // of the following coordinates in the world coordinates : 
     let cameraX = cos(0.5 + i * 0.002)*250 ;
     let cameraY = sin(0.5 + i * 0.002)*250;
     let cameraZ = 200+sin(0.5 + i * 0.002)*30;
     camera(cameraX,cameraY, cameraZ, 100, 0, 100, 0, 0, -1);
     // the output images are in the projection coordinates.
   
     //The world coordinates : 
     background(BACK_COLOR);
     directionalLight(250, 200, 200, 0, -1, -1); //light from the sun 
     fill(250, 200, 200);
     plane(10000,10000);
   
     stroke(0);
     strokeWeight(2);
     fill(150,150, 150);
     pyramid(3);
   
     translate(-250,0,0);
     pyramid(3);
   
     translate(-250,0,0);
     pyramid(1.5);
   
     ambientMaterial(125);
     ambientLight(100,100,100);
     shininess(100.0);
   
     translate(700, 200, max(500 - i * 0.3, 67) );
     sCube(5);
     
     let M = 25;
     for(let m = 0 ; m < M ; m++){
       translate(140,140, max(650 + 150*m - i * 0.3, 0));
       sCube(5);
     }  
     i++; 
   }
   
   // The modeling coordinates :
   function pyramid(scaleN){
     push();
       rotateZ(QUARTER_PI);
       scale(scaleN);
       model(octahedron);
     pop();
   }
   
   function sCube(scaleN){
     push();
       rotateZ(QUARTER_PI);
       scale(scaleN);
       box(25);
     pop();
   }
\end{verbatim}


\subsection{p5.js の3Dモードの説明とソースファイルの解説}
関数preload()でディレクトリにある図形 octahedron.stl を読み込み，
関数setup()の中でウィンドウ等を初期化する．
関数draw()が呼び出されると変数 i を1フレーム毎に1づつ増やしながら，
その都度フレームを描画し続ける．いじってみるとだんだん理解が深まると思う．

\emptyline

p5.js 特有の組込関数の説明：
\begin{itemize}
   \item
      translate(a, b, c); $(a, b,c)$だけ平行移動し，その点を$(0,0,0)$として描画
   
   \item
      rotateZ(theta); $z$軸を中心に角度 theta だけ直交座標系を回転し，それを新たな直交座標系として描画．rotateX, rotateY も同様．
   \item
      scale(a) 座標系をa 倍する
   
   \item  
      strokeWeight(a) 太さaで描く
     
   \item
      background(a,b,c);  背後の画面をRGB [a,b,c] で埋める．
      ただし$a, b, c$は$0$から$255=2^8-1$ の間の整数．
   
   \item directionalLight(a, b, c, d,e,f) RGB[a, b, c]の光で向き(d,e,f) $-1\le d,e,f\le 1$ で光を照らす
   
   \item
       ambientLight(a, b, c) RGB[a, b, c]の光で背景となる光を照らす
       
    \item 
       ambientMaterial(a, b, c) 次の図形の壁の色をRGB [a,b,c] で設定
   
   \item 
      fill(a,b,c) 次の図形をRGB [a,b,c] で埋める．
   
   \item 
      plane(a,b), box(c) などは3次元の組込立体図形
   
   \item
      camera(a, b, c, a', b', c') カメラの位置$(a, b, c)$ 対象の中心$(a', b', c')$ に向けてウィンドウを描画
   
   \item
      push(), pop() はその間に変えられた設定を元に戻すためのスタックに書き込むための関数
   
   \item
      model(modelName) modelNameを描く

\end{itemize}



\begin{exQ}
GitHub のアカウントがあればGitHub のPages という機能を使うと
自分のホームページが作れます．
私のポートフォリオもそのようにして作ったものです．
やり方を調べて自分のホームページを作って，自分の作った作品を公開してください．
\end{exQ}


%\begin{thebibliography}{9}
%   \bibitem{P5}
%      松田晃一, p5.jsプログラミングガイド  改訂版, 2021, カットシステム
%   
%   \bibitem{YY}
%      山田祥寛, 
%      改訂3版JavaScript本格入門, 2023, 技術評論社
%\end{thebibliography}
%\newpage
%
%\section{補足：透視投影と3DCGと線形代数}
%3DCG(3次元コンピューターグラフィックス)の数学的側面とプログラミング的側面が線型代数や行列の計算に関連していることを説明したい．
%%JavaScriptについても説明する．
%
%\subsection{3DCGと遠近法}
%3DCGは現実の見え方をシュミレートしたものである．
%ここではピラミッドの写真(図\ref{pyramid})を使って3DCGの原理を図示する(図\ref{model})．
%\begin{figure}[htb]
%   \centering
%   \includegraphics[width=7.5cm]{~/Documents/GitHub/introInfo/appc/images/2428333990_d3c913b2b0_c.jpg}
%   \caption{\href{https://www.flickr.com/photos/girolame/2428333990/in/photolist-4GvJhV-4GzR6q-68Y3JE-4GzTym-5DUvzA-4Gvzcg-5UXT1K-tAvViv-dmRS35-6j7Tdf-65175S-7hkV8W-7ZP8Wc-f7stxJ-6kADQr-92FxE-9MpvoF-7DWyqL-7DWzkj-7trRb8-7kV2Kh-74NeWm-6fnJRK-gBqRa-A4k2s-5DLk1m-A4jWP-6JPcdB-erLBS-7KLGJq-fWSgy-7bxjK9-5DRPd1-dduMZ-5DUGgw-6iWpt7-oueimu-oeLdfu-6GS2bv-6JCrjr-5Phx2c-7KGHVB-29jwn5-6JQBeo-oeDc4c-5DUBKb-dvw1V-5UXWtP-ouBgy3-gBn4K}{Giza Pyramids} by Jerome BonFollow 
%   \href{https://creativecommons.org/licenses/by/2.0/?ref=openverse}{\ccby}}\label{pyramid}
%\end{figure}
%カメラでは像はレンズの後の壁にあり，しかも180°回転しているためイラストを作図しづらいので（他にもレンズがあると被写界深度が問題になることや，人間の目の構造上網膜が平坦でないことによるカメラや遠近法とのずれなど微妙な違いが有る．詳しくは図~\ref{pinCamera}, \ref{camera}, \ref{eye}を参考にしてほしい )，以下では遠近法を用いた作図をつかい，画家の視点から(透明な)キャンバス面への被写体を投影したものを考える．
%\begin{figure}[htb]
%   \centering
%   \begin{minipage}{0.4\columnwidth}
%      \centering
%      \includegraphics{~/Documents/GitHub/introInfo/appc/images/figure.14}
%      \caption{カメラ・オブスキュラ%フェスメール？
%      }\label{pinCamera}
%   \end{minipage} 
%   \begin{minipage}{0.4\columnwidth}
%      \centering
%      \includegraphics{~/Documents/GitHub/introInfo/appc/images/figure.15}
%      \caption{カメラ}\label{camera}
%   \end{minipage} 
%\end{figure}
%\begin{figure}[htb]
%   \centering
%   \includegraphics{~/Documents/GitHub/introInfo/appc/images/figure.18}
%   \caption{目---カメラとの微妙な違い}\label{eye}
%\end{figure}
%
%光は直進するので，視点と被写体上の点を結ぶ直線(視線)とキャンバスの平面との交わりを求めることが遠近法になる．
%%これは連立一次方程式を解くことにあたる．
%%特に，空間内の直線と視点を結ぶ直線全体は空間内の平面であるので，キャンバスとの交わりはキャンバス上の直線に投影される．
%
%\begin{figure}[htb]
%   \centering
%   \includegraphics{~/Documents/GitHub/introInfo/appc/images/figure.2}
%\caption{遠近法(透視投影)}\label{model}
%\end{figure}
%
%\subsection{(透視)投影}
%\textcolor{red}{遠近法(perspective)}でどのように直線が描かれるかを説明したい：
%原点$\mathbf 0$に視点があるとする．
%$(x,y,z) \in \mathbf R^3$  ($x\neq 0$ )とし，
%ここでは，$x$軸の正の方向に視線を向け，
%$x=1$ にキャンバスがある場合を考える．
%$U=\{(x,y,z)\in \mathbf R^3\mid x\neq 0\}$として各点$(x,y,z)$と原点$\mathbf 0$を結ぶ直線と$x=1$の面の交点を求める，つまり，連立一次方程式を解くと$(1,\frac yx,\frac zx)$がわかる(図\ref{perspective})．その結果を与える写像
%\[
%   P(x,y,z)=\biggl(\frac yx,\frac zx\biggr)\qquad (x,y,z)\in U
%   \tag{*}\label{projection}
%\]
%$P:U\to \mathbf R^2$が\textcolor{red}{(透視)投影}(あるいは，\textcolor{red}{射影})である．
%この$ \mathbf R^2$, あるいは，$ \mathbf R^2$を $\mathbf R^3$の$x=1$の面とみたものを\textcolor{red}{投影面}と言う(図\ref{model})．
%
%\begin{figure}[htb]
%   \centering
%   \begin{minipage}{0.4\columnwidth}
%      \centering
%      \includegraphics{~/Documents/GitHub/introInfo/appc/images/figure.5}
%      \caption{透視投影}\label{perspective}
%   \end{minipage} 
%   \begin{minipage}{0.4\columnwidth}
%      \centering
%      \includegraphics{~/Documents/GitHub/introInfo/appc/images/figure.17}
%      \caption{直線の射影}\label{perspective2}
%   \end{minipage} 
%\end{figure}
%
%定義域が$U\neq \mathbf R^3$なので線形写像にならないことは明らかだが
%\[
%   P(ax,ay,az)=P(x,y,z) \quad a\neq0
%\]
%という性質がある．
%\eqref{projection}は$y$, $z$については線形写像で，線形写像を少しひねっただけの簡単なものである．
%
%$ \mathbf R^3$で点$(x,y,z)$から向き$(u,v,w)$に向かう直線は
%\[
%   (x,y,z)+t(u,v,w)
%\]
%と表せるので
%\[
%   P(x+ut,y+vt ,z+wt)=\biggl(\frac{y+vt}{x+ut},\frac{z+wt}{x+ut}\biggr)\tag{!}
%\]
%これは$t=0$で$(\frac yx,\frac zx)$を通り，$u\neq 0$のとき$t\to \infty$で$(\frac{v}{u},\frac{w}{u})$に収束する．この$(\frac{v}{u},\frac{w}{u})$を\textcolor{red}{消失点(vanishing point)}という．
%差をみると
%\[
%   \biggl(\frac{y+vt}{x+ut},\frac{z+wt}{x+ut}\biggr)
%   =
%   \biggl(\frac{v}{u},\frac{w}{u}\biggr)
%   +
%   \frac{1}{x+ut}\biggl(y-x\frac{v}{u},z-x\frac{w}{u}\biggr)
%\]
%となり，最後の項は定数ベクトルを$\tfrac{1}{x+ut}$倍したものなので，投影面内で$(\frac yx,\frac zx)$を通る直線で消失点へ至るものであることが分かる(図\ref{perspective2}に初等幾何学的な証明をあげた)．
%線形写像は直線を直線に写すが，投影は直線を直線の一部(半直線)に写すことがわかった．
%また，平行な直線は同じ消失点に収束することも分かる．
%さらに，$u=0$のとき
%$(\frac{y+vt}{x},\frac{z+wt}{x})$となり，等速の直線運動になる．
%
%
%遠近法の数学と直接関係ないが，地面や海面を表す平面が絵を描く（写真をとる）場合には大切になる．
%地面で画家やカメラマンがたっている点を立点，その真上に視点をとり立点と視点を通る軸を$z$軸とする．地面は通常$z=\const$に取る．（これは必ずしも必要でないが）地面に平行などちらかへの視線向きを$x$軸，それに直交する（右手系になる）軸を$y$軸ととると上で考えた座標系と一致する．
%地面$z=\const$の消失点全体が地平線をなす．
%遠近法の基本的な手段は地平線上の
%いくつかの消失点とそこに収束する平行直線群を見いだし，それを元に描画することである(図~\ref{1VanishingPoint}, \ref{2VanishingPoint})．もちろん，地面に垂直な線は$z$軸に平行に描かれる．
%
%\begin{figure}[htb]
%   \centering
%   \includegraphics{~/Documents/GitHub/introInfo/appc/images/figure.7}
%   \caption{1消失点の平行直線群}\label{1VanishingPoint}
%\end{figure}
%
%\begin{figure}[htb]
%   \centering
%   \includegraphics{~/Documents/GitHub/introInfo/appc/images/figure.16}
%   \caption{平行直線群の作成}\label{parallelLines}
%\end{figure}
%
%\begin{figure}[htb]
%   \centering
%   \includegraphics[trim=13  0 305 0 , clip]{~/Documents/GitHub/introInfo/appc/images/figure.8}%
%   \caption{2消失点の平行直線群}\label{2VanishingPoint}
%\end{figure}
%
%よく見られる一消失点と平行直線群($u\neq0$, $v=w=0$ に対応)と直交する交わらない平行直線群($v\neq0$, $u=w=0$ に対応)をどう作図するかを説明する：
%平面を$z=-1$に固定する．この平面とカンバスの最底辺$\ell_1$を固定する．その平面上の$\mathbf R^3$の中での直交する平行直線は
%\[
%    (0,y_0,-1)+s(1,0,0),\quad (x_1,y_1,-1)+t(0,1,0),
%\]
%なので射影面では
%\[
%   \biggl(\frac{y_0}{s},\frac{-1}{s}\biggr),\quad 
%   \biggl(\frac{y_1+t}{x_1},\frac{-1}{x_1}\biggr)
%\]
%と表せる．
%まず，直線群は幅$a$, $b$に並んでいるとし，
%\[
%   y_0 = am,\quad
%   x_1 = bn\quad  m\in \mathbf Z, n\in \mathbf N
%\]
%と取り
%\[
%   {\ell'}_m:\biggl(\frac{am}{s},\frac{-1}{s}\biggr)\quad
%    (s\ge b),\quad
%   \ell_n: \biggl(\frac{y_1+tv}{bn},\frac{-1}{bn}\biggr)
%\]
%と置く．${\ell'}_n$の作図はカンバスの最底辺を表す直線$\ell_1$の$0$を通る等分点と消失点を結べば良い．カンバスの外については後にコメントする．
%
%${\ell'}_m$と$\ell_n$の交点は
%\[
%   s=bn
%\]
%であり，
%$n$, $m$はインデックスに過ぎないので，${\ell'}_n$では
%\[
%   s=b n\quad 
%   \Rightarrow \quad
%   \frac{an}{s} = \frac{a}{b}
%\]
%なので射影面上で$y= \frac{a}{b} $と${\ell'}_n$, $y=- \frac{a}{b} $と${\ell'}_{-n}$の交点を求め，それを通る直線として${\ell}_n$を作図できる．$\ell_1$を$\ell_2$に替え，
%その上で得られた等分点を拡張することで，画面内でどこまでも続けられる．
%
%1消失点を使った構図(図~\ref{1VanishingPoint}), 
%2消失点を使った構図(図~\ref{2VanishingPoint})，3消失点使った構図(図~\ref{3VanishingPoint})とその写真例を挙げる(図~\ref{ex1}, \ref{ex2}, \ref{ex3})
%
%\begin{figure}[htb]
%   \centering
%   \includegraphics[width=7.5cm]{~/Documents/GitHub/introInfo/appc/images/8367600279_ac56e7454d_b.jpg}
%   \caption{写真例 \href{https://www.flickr.com/photos/31583394@N04/8367600279}{Vanishing Storm} by David Kingham\href{https://creativecommons.org/licenses/by-nd-nc/2.0/jp/?ref=openverse.}{\ccbyncnd}}\label{ex1}
%\end{figure}
%\begin{figure}[h]
%   \centering
%   \includegraphics[width=8cm]{~/Documents/GitHub/introInfo/appc/images/9110846773_dad72ed427_k.jpg}
%   \caption{写真例 \href{https://www.flickr.com/photos/30996501@N07/9110846773}{MEZQUITA CATEDRAL DE CORDOBA 6378} by Jose Javier Martin Espartosa \href{https://creativecommons.org/licenses/by-nd-nc/2.0/jp/?ref=openverse.}{\ccbyncsa}}\label{ex2}
%\end{figure}
%
%\begin{figure}[htb]
%   \begin{minipage}{0.5\columnwidth}
%      \centering
%      \includegraphics{~/Documents/GitHub/introInfo/appc/images/figure.9}
%      \caption{3消失点を使った直方体の作図}\label{3VanishingPoint}
%   \end{minipage} 
%   \begin{minipage}{0.4\columnwidth}
%      \centering
%      \includegraphics[width=4cm]{~/Documents/GitHub/introInfo/appc/images/145347226_76283ffb36_c.jpg}
%      \caption{写真例 \href{https://www.flickr.com/photos/46965656@N00/145347226}{Barcelona - Gaudi building} by Erick \href{https://creativecommons.org/licenses/by-nd-nc/2.0/jp/?ref=openverse.}{\ccbyncnd}}\label{ex3}
%   \end{minipage} 
%\end{figure}
%
%\begin{figure}[htb]
%   \centering
%   \begin{minipage}{0.4\columnwidth}
%      \centering
%      \includegraphics[width=5.8cm]{~/Documents/GitHub/introInfo/appc/images/figure.10}
%      \caption{透視投影と平行投影}\label{paraPerspective1}
%   \end{minipage} 
%   \begin{minipage}{0.4\columnwidth}
%      \centering
%      \includegraphics[width=4.7cm]{~/Documents/GitHub/introInfo/appc/images/figure.11}
%      \caption{平行投影}\label{paraPerspective2}
%   \end{minipage} 
%\end{figure}
%
%
%投影面を$x=0$に止めて，あらためて視点を$(-a,0,0)$にとり，$a\to \infty$として無限遠まで動かしたとき得られる投影を\textcolor{red}{平行投影}という．このとき透視投影の極限として
%\[
%   \biggl(\frac{ a}{x+a}y,\frac{ a}{x+a}z\biggr)
%   \to (y,z)\qquad (a\to \infty)
%\]
%となる(図\ref{paraPerspective1}, \ref{paraPerspective2})ので，これは線形写像である．ただし，平行投影は挿絵(この資料の図もそうである)図学等ではよく使われるが，3DCGでは現実感を出すため透視投影が主に使われる．
%
%\subsection{反射，拡散と影}
%ここから3DCGに関連する話題に移る．
%まず大事なのは光とその反射面の取り扱いである．ここでは屋外なので光源として太陽を考えている．
%スタジオ撮影では多数のライトを用い，光の強度は距離の二乗に反比例することも考慮する必要がある．逆に太陽の場合は無限の彼方にあると考え，距離に依らない強度を持つとする．
%光源から出た光は直進し
%被写体の面で反射して，また直進してキャンバスと視点に到達するので，
%光の平面での\textcolor{red}{鏡面反射}を計算する必要がある．
%これは与えられた平面での光源の鏡映を求め，その点と視点を結ぶ直線と平面の交点を求めることで求まる(図12)．与えられた平面が$z=0$とすると，点$(x,y,z)$のこの平面での鏡映は行列を用いて
%\[
%   \begin{pmatrix}
%      1&0&0\\
%      0&1&0\\
%      0&0&-1\\
%   \end{pmatrix}
%   \begin{pmatrix}
%      x\\
%      y\\
%      z\\
%   \end{pmatrix}
%   =
%   \begin{pmatrix}
%      x\\
%      y\\
%      -z\\
%   \end{pmatrix}
%\]
%
%
%\begin{figure}[htb]
%   \centering
%   \includegraphics{~/Documents/GitHub/introInfo/appc/images/figure.3}
%   \caption{鏡映と鏡面反射, 拡散}\label{ref}
%\end{figure}
%
%\noindent
%と表せる．
%光源の位置を$(x,y,z)$とし，視点の位置を$(x_0,y_0,z_0)$
%とすると途中の点は
%\[
%   \begin{split}
%      &t(x, y, -z)+(1-t)(x_0, y_0, z_0)\\
%      =&
%      (tx+(1-t)x_0, ty+(1-t)y_0,-tz+(1-t)z_0)
%   \end{split}
%\]
%と表される．反射する点では$-tz+(1-t)z_0=0$なので，連立一次方程式を解くと
%\[
%   \biggl(\frac{zx_0+z_0x}{z+z_0},\frac{zy_0+z_0y}{z+z_0},0 \biggr)
%\]
%のように反射点が求まる(図 \ref{ref})．
%さらには，物理特性である面の色・材質による拡散も計算する必要がある．
%
%\begin{figure}[htb]
%   \centering
%   \includegraphics{~/Documents/GitHub/introInfo/appc/images/figure.20}
%   \caption{隠面消去・シェーディング・影付け}\label{shadow}
%\end{figure}
%
%また，一つの物体の見えない部分や，二つの物体の重なりを計算し遠くの物体の隠れた部分を消す\textcolor{red}{隠面消去}や，
%光源によって照らされた部分の明るさや色を求める
%\textcolor{red}{シェーディング}, 
%影をシュミレートする\textcolor{red}{影付け}も必要になる(図 \ref{shadow})．
%
%なお，遠近法には\textcolor{red}{空気遠近法}という，遠くのものを
%薄く描いて遠近感を出す方法も知られている．今回は詳しい説明は与えないが，図 \ref{1VanishingPoint}, \ref{2VanishingPoint} はその原理に
%基づいて作図されている．
%
%\subsection{アファイン変換}
%$3$次元の\textcolor{red}{アファイン変換}とは
%\[
%   \begin{split}
%      \begin{pmatrix}
%    c_1\\
%    c_2\\
%    c_3\\
%      \end{pmatrix}
%      +
%      \begin{pmatrix}
%    a_{11}&a_{12}&a_{13}\\
%    a_{21}&a_{22}&a_{23}\\
%    a_{31}&a_{32}&a_{33}\\
%      \end{pmatrix}
%      \begin{pmatrix}
%    x_1\\
%    x_2\\
%    x_3\\
%      \end{pmatrix}
%      &=
%      \begin{pmatrix}
%    c_1 + a_{11}x_1 + a_{12}x_2+ a_{13}x_3\\
%    c_2 + a_{21}x_1 + a_{22}x_2+ a_{23}x_3\\
%    c_3 + a_{31}x_1 + a_{32}x_2+ a_{33}x_3\\
%      \end{pmatrix}\\
%      &=
%      A\mathbf x + \mathbf c
%   \end{split}
%\]
%となる写像で，平行移動と行列の積の合成である．
%通常$A=(a_{ij})
%%\begin{pmatrix}
%%   a_{11}&a_{12}&a_{13}\\
%%   a_{21}&a_{22}&a_{23}\\
%%   a_{31}&a_{32}&a_{33}\\
%% \end{pmatrix}
%$
%は正則行列と仮定する．特に，合同変換はアファイン変換である．
%以後アファイン変換はCGの作成で何度も現れる．
%
%これを直線
%\begin{equation*}
%   \mathbf x = \mathbf a + t \mathbf u
% \end{equation*}
%に適用すると
%\[
%   A(\mathbf a + t \mathbf u)+\mathbf c
%   =
%   \mathbf c+A\mathbf a+tA\mathbf u\quad A\mathbf u\neq \mathbf 0
%\]
%つまり，アファイン変換は直線全体を直線全体に写し，平行な直線を平行な直線に写す．
%また，直線上の三点の比を同じ比の線分上の点に写すことも分かる．
%
%行列の積で
%\[
%   \begin{pmatrix}
%      1&0&0&0\\
%      c_1&a_{11}&a_{12}&a_{13}\\
%      c_2&a_{21}&a_{22}&a_{23}\\
%      c_3&a_{31}&a_{32}&a_{33}\\
%   \end{pmatrix}
%   \begin{pmatrix}
%      1\\
%      x_1\\
%      x_2\\
%      x_3\\
%   \end{pmatrix}
%   =
%   \begin{pmatrix}
%      1\\
%      c_1 + a_{11}x_1 + a_{12}x_2+ a_{13}x_3\\
%      c_2 + a_{21}x_1 + a_{22}x_2+ a_{23}x_3\\
%      c_3 + a_{31}x_1 + a_{32}x_2+ a_{33}x_3\\
%   \end{pmatrix}
%\]
%となるので，その第2第3第4成分がアファイン変換に対応している．
%これはアファイン変換の行列の積を使った表現であり，
%その意味でアファイン変換は線形変換に近いと言えるが，第1成分が$1$に固定されているので，これは線形写像ではない．このような座標系を\textcolor{red}{非同次座標系}という．これは行列の計算なのでコンピュータでの組織的な計算に適している．
%
%\subsection{モデリング，ビューイングパイプライン}
%\begin{figure}[htb]
%   \centering
%   \includegraphics{~/Documents/GitHub/introInfo/appc/images/figure.4}
%   \caption{ビューイングパイプライン}\label{viewingPipeline}
%\end{figure}
%
%
%コンピュータの中で被写体の曲面の三次元データを作ることを\textcolor{red}{モデリング}という．これは\textcolor{red}{モデリング座標系}での，ピラミッドのようなポリゴンの頂点とそれらの張る辺，三角形の作る面の幾何学的情報と色や質感等の物理情報の組からなる．
%シュミレートしたい被写体とカメラの位置・向きの指定された三次元(データの)世界を\textcolor{red}{ワールド座標系}という．幾つかのモデリングのデータを適当なアファイン変換%(下の補足Aで示すように行列の積で表現される)
%でうつすことで，ワールド座標系にシュミレートしたい被写体のデータを移している(図\ref{viewingPipeline})．
%ここでは，2つのピラミッドを一つのデータから構成している．
%ワールド座標系の中での視点の座標と視線の方向を定めることで，
%その視点から見た座標系が定まるが，
%それはワールド座標系をアファイン変換したものであり\textcolor{red}{カメラ座標系}という．
%視点から投影面に（透視投影した）二次元座標系を\textcolor{red}{投影座標系}という．
%ただし
%実際は投影座標系の中の最終的な画面出力となる
%\textcolor{red}{ウィンドウ}と呼ばれる長方形部分のみを計算すればよい．
%
%ワールド座標系とカメラ座標系のデータ及び，光源，反射，陰影等からウィンドウへの投影を計算し，ウィンドウ内のデジタルデータを作図するのを\textcolor{red}{レンダリング}という．
%これらを少しずつ動かしながら繰り返すことでアニメーションが構成できる．
%視点を動かし，カメラワークをシュミレートすることも可能である．
%このような幾つもの操作の流れを\textcolor{red}{ビューイングパイプライン}という．
%これらのことから，線型代数が3DCGの中で占める役割の大きさが分かると思う．
%
%\emptyline
%CGに興味を持った人には\cite{CG}, 遠近法の歴史については\cite{D},
%具体的な遠近法の絵の描き方については\cite{3DT}, \cite{RS}, \cite{RSII} などをお奨めする．
%
%\begin{thebibliography}{9}
%   \bibitem{CG}
%   画像情報教育振興協会, コンピューターグラフィックス[改訂新版], 2015.
%   
%   \bibitem{D}
%   下村 耕史訳編, アルブレヒト・デューラー「測定法教則」注解, 中央公論美術出版, 2008.
%   
%   \bibitem{3DT}
%   3DTotal.com, 
%   デジタルアーティストが知っておくべきアートの原則 改訂版 -色、光、構図、解剖学、遠近法、奥行き 改訂版, 2021.
%   
%   \bibitem{RS}
%   スコット・ロバートソン, 
%   スコット・ロバートソンのHow to Draw -オブジェクトに構造を与え、実現可能なモデルとして描く方法, 2014.
%   
%   \bibitem{RSII}スコット・ロバートソン, スコット・ロバートソンのHow to Render：アイデアを明確に伝える 光と影、反射の描き方, 2015.
%\end{thebibliography}
%
%\end{document}



%%\documentclass[a4paper,11pt]{jlreq}
%\documentclass[a4paper,12pt]{amsart}
%\usepackage[dvipdfmx]{graphicx,xcolor}%graphics
%\usepackage{float}
%\usepackage{amsmath}
%\usepackage{verbatim}
%\usepackage{latexsym}
%\usepackage{amssymb}
%\usepackage{qrcode}
%%\usepackage{emath}
%\usepackage{listings}
%\usepackage{ccicons}
%\usepackage{mflogo}
%%\usepackage{fancybox}%enclose by box
%\usepackage[dvipdfmx]{hyperref}
%\usepackage{verbatim}
%\hypersetup{% hyperrefオプションリスト
%   setpagesize=false,
%   bookmarksnumbered=true,%
%   bookmarksopen=true,%
%   colorlinks=true,%
%   linkcolor=magenta,
%   citecolor=blue
%} 
%\newtheorem*{exQ}{レポート}
%\theoremstyle{remark}
%\newtheorem*{Q}{演習}
%
%\renewcommand{\figurename}{図}
%\renewcommand{\proofname}{証明}
%\renewcommand{\tablename}{表}
%
%\def\emptyline{\vspace{11pt}}
%\def\smallemptyline{\vspace{6pt}}
%\begin{document}
%\begin{flushright}
%   2026/08/02
%   %2026/06/15
%   %2026/04/14
%   %2026/04/04
%   %2026/03/29
%   %2026/03/11
%   %2026/02/21
%   %2026/02/10
%   %2025/12/22
%   
%   snow00two
%\end{flushright}
%
%\begin{center}
%   {\Large コマンドラインとShell Script 入門}
%   \href{https://creativecommons.org/licenses/by-nc-nd/4.0/}{\ccbyncnd}
%%   \par{Command Line Interface}
%\end{center}

\chapter{ShellとShell Script 入門}

%\section{ファイルシステム}
%現在ではGUI (Graphical User Interface)で色々なことができるようになり，
%実際ほとんどの人がGUI を使ってると思うが，
%Command Line Interface(CLI, CUI =Character User Interface と呼ばれることもある) でないと出来ないことも多い．
%何が出来て何が出来ないかはCLIを使ってみたことがないと分からない
%と思うが，\textcolor{red}{生成AIもCLIを使いこなしている．}
%ここでは CLI のOSの例としてUnixを取り上げ
%(MacOSはUnixの兄弟FreeBSDを基に作られていて，
%Linux もUnixの親戚に当たる), 
%まずUnixのファイルシステムについて紹介する． 
%shell の使い方については後のセクションで説明する．
%\begin{itemize}
%   \item
%      ディレクトリー：% DOM
%      \begin{figure}[htb]
%         \centering
%         \includegraphics{../fig/figure.1}
%         \caption{ファイルシステムのツーリー構造}\label{tree}
%      \end{figure}
%   \begin{figure}[htb]
%         \centering
%         \includegraphics[width=14cm]{../fig/column.pdf}
%         \caption{Column 表示でのツーリー構造の確認}\label{column}
%      \end{figure}
%            
%         ディレクトリーやファイルは図~\ref{tree}のようなツリー構造に配置されている．
%         特にパソコンのハードディスクの一番上の部分をルートといい \verb/// と表す．
%         全てのディレクトリーやファイルはルートから一意的にたどっていく
%         ことが出来る．
%         ファイルやディレクトリが\verb/// から見てどこにあるかを示すのを
%         \textcolor{red}{（絶対） path}という．
%         例えば
%         図~\ref{tree}のprogramming.texは
%         
%         /Users/sonw00two/Documents/classRoom/General/
%         
%         Programming/src/programming.tex
%         
%         \noindent
%         のように表される．
%         後で出てくる .. は一つ上のディレクトリーをあらわす\textcolor{red}{相対 path}の例
%
%      \item
%      不可視ファイル：
%      ディレクトリ /usr  は
%         普通には見えない設定になっているので
%         知らないとアクセスできない．
%         Finder のメニューバーの Go のメニューから Go to Folder を選ぶと
%         四角いウィンドウが開くので /usr と入力し，goをクリックで移動できる．
%         Finderで cmd shift G とタイプしてもよい．
%         
%   \item
%         拡張子：%とpermission 
%         ファイル名には hoge.tex の .tex のような
%         \textcolor{red}{拡張子(extension)}があるが，
%         多くの場合アイコンのみ表示され見えないようになっている．
%         開発のときは様々なファイルを扱うので，
%         拡張子を表示しないと効率的なファイルの管理が難しい．
%         また，メールの添付ファイルのアイコンは偽造できるが
%         拡張子により直ぐに峻別できるので，
%         セキュリティーの観点からも表示するのが望ましい．
%   
%   \item        
%         関係したファイルは一つのディレクトリーに，
%         関連していないファイルは異なるディレクトリーにいれ，
%         関連のある，あるいは類似のディレクトリーは纏めて一つの大きなデレクトリー
%         に入れるなど入れ子構造にして（それがツリー構造の定義である）
%         管理するのが基本である．
%         また，簡単にそのディレクトリーを説明するテキストファイル ReadMe.txt, 
%         README.mdをそのディレクトリーに置くことも良い習慣である．
%
%
%\begin{Q}
%   \begin{figure}[h]
%         \centering
%         \includegraphics[width=8cm]{../fig/ScreenCapture1.pdf}
%         \caption{Macでのエクステンションの表示}\label{extension}
%      \end{figure}
%   拡張子を表示するようにパソコンを設定して下さい．
%\end{Q}
%         
%
%\section{環境構築}
%以下のアプリのインストール・色々なサービスへの登録をしてもらいたい．
%ただし，CLI では破滅的なことも出来るので，
%バックアップは自分でこまめに取っておくようにして下さい．
%この過程で生じた被害への補償は一切行えません
%%\textcolor{red}{
%%   プログラミング実践演習で生じた被害への補償は一切行えません．
%%   この免責事項を認めた人だけがこの授業を受講することが出来ます．
%%}
%
%
%
%  \begin{figure}[htb]
%         \centering
%         \includegraphics[width=14cm]{../fig/vscode.pdf}
%         \caption{VSCode, 下が terminal}\label{vscodeTerm}
%   \end{figure}
%
%\begin{Q}
%\begin{enumerate}
%   \item
%      VSCodeのインストール：
%       
%                   フリーのエディタというより統合開発環境
%            (IDE = Integrated Development Environment)Visual Studio Code
%            （略称VSCode, \href{https://code.visualstudio.com}{ここからダウンロード}
%            可能，図 \ref{vscodeTerm}参)．
%            様々なコンピューター言語のための（自動補完，構文解析のための）
%            インテリセンス，リンター，スニペット，ユニットテスト
%            \footnote{詳しい説明はネット等で調べてください.}等の機能拡張がある．
%            マイクロソフトのVisual Studio  シリーズのエディタのみをフリーで公開した
%            もので，他にも似たものが色々あったが，最近はデファクトスタンダード．
%            JavaScriptで作られたアプリの一つで
%            Chromium+JavaScript +Node.js+Electron 
%            により作られている(図 \ref{vscode})．
%     
%      \begin{figure}[htb]
%               \centering
%               \includegraphics[width=6cm]{../fig/ScreenCapture2.pdf}
%               \caption{VSCodeのクレジット}\label{vscode}
%   \end{figure}
%   
%   \item
%         terminal と shell:
%       
%       Macや Linux 等では既にインストールされている．
%       Windows ではPowershell あるいは
%       WSL(Windows Subsystem for Linux)をインストロールする必要があるだろう．
%       
%      これから必要になるtool等もCLIと上でも述べたtoolなどを使って
%       インストールする必要のあるものが多い．
%       Mac の標準の terminal は
%      
%       /Applications/Utilities/Terminal.app
%   
%       \noindent
%       にある．
%      Mac なら (Home)brew もインストールしておくと都合が良い．
%      \href{https://brew.sh/ja/}{ここからダウンロード}．
%      pathを通す必要もある．
%       
%       不可視なディレクトリやファイルも沢山あるので
%       適宜コマンドで確認する必要がある．
%      VSCode の下の部分にもterminal が付いている（図 \ref{vscodeTerm}）．
%
%%     \item
%%      Web 技術：JavaScript ( TypeScript )
%%      
%%      \begin{enumerate}
%%         \item  http とコマンドラインによるbrowserもどきの作成
%%         
%%         \item 
%%              Chrome のDeveloper tool, Chromeには最初からDeveloper menu がある.
%%             Developer ニューから JavaScript Console を開くと
%%             Console が開き，デバッグが可能．
%%         
%%         \item  
%%            JXA : JavaScriptが Mac のスクリプト言語として採用された：
%%         
%%            Script Editor・ターミナルから osascript・automator で使える．
%%         
%%         \item 
%%            Single page application: JavaScript+Node.js+Electron etc.+Chromium でmacOS, 
%%            Windows, スマートホンのいずれにも対応したデスクトップアプリケーションが
%%            一度に作れる． 例VSCode 
%%         
%         
%%         \item
%%            p5.js ：
%%            JavaScript で記述し，2D, 3Dの映像・動画の作成可能．
%%            Processing というJavaのライブラリーをJavaScript に移植したもの
%%            (\cite{P5})．
%%      \end{enumerate}
%   \item
%      version 管理の基本ソフト Git とネットサーバー GitHub:
%      多人数で開発を行うにはこれを使いこなさないといけないが，
%      個人開発でも有用．
%      GitHubで VSCode の設定ファイルの同期も出来る．
%      あなたが所属する組織がMicrosoft365に加入していればそ
%      のアドレスでGitHubに
%       \href{ https://github.com/}{登録}し
%      Student Developer Pack を使うように設定できる可能性がある．
%
%   \item
%      Microsoft365 にログインすると Copilot
%      （chatGPTを office suite 向きにアレンジしたもの）が
%      使用料無しで使えるようになっている．
%      GitHub のStudent Developer Pack に入っていると，
%      VScode の拡張機能タブで「GitHub Copilot」を検索してインストールすると，
%       使用料無しでVScodeからGitHub Copilot が使える可能性がある．
%%   chatGPT etc
%%      \begin{enumerate}
%%      \item
%%         機械学習
%%      \item
%%         プロンプトエンジニアリング(\cite{OT})．
%%      \end{enumerate}
%
%   \item
%      JavaScript は  Chrome などのBrowser で既に使えるようになっているが，
%      セキュリティのためディスクへのアクセスなどは制限されているので，
%      強力な実行環境となる node.js とその管理ツール npm も必要になる．
%      \href{https://nodejs.org/ja/download}{ここからダウンロード}.
%      これからは
%      \href{https://github.blog/news-insights/octoverse/octoverse-a-new-developer-joins-github-every-second-as-ai-leads-typescript-to-1/}
%      {TypeScript}にも慣れる必要があるだろう．
%            
%   \item
%      Python3 とその管理ツール pip や  (Ana)conda  も必要
%      \href{https://www.python.org}{ここからダウンロード}
%   
%\end{enumerate}
%\end{Q}
%
%
%
%
%%   \item
%%
%%
%%shell 
%%
%%root 
%%
%%ログインディレクトリー
%%作業ディレクトリ
%%
%%ファイル管理
%%原則
%%Login
%
%\end{itemize}
%
%
%\begin{Q}
%   VSCode の設定を，必要なら検索かCopilotに聞いて，
%   GitHubで同期するように設定して下さい．
%\end{Q}
%
%
%\section{CLI}
%      コマンドラインインターフェースを使いこなさないとプログラミングは難しい．
%        以下に簡単なコマンドの使用例をあげる．terminal を立ち上げるとプロンプト
%        が表示される(図~\ref{terminal})．
%        私の端末では$\cdots$\verb/% /となっているのがプロンプトで，
%        そのままキーボードからコマンド（文字）を打ち，リターンキーを押すと
%        そのコマンドが実行され，次の行に結果が表示される．
%        これを標準出力という．標準入力と標準エラー出力もある．
%        \begin{figure}[htb]
%         \centering
%         \includegraphics[width=14cm]{../fig/terminal.pdf}
%         \caption{Terminal}\label{terminal}
%      \end{figure}
%       これがコマンドラインインタープリター shellである．
%       
%       以下では \verb/% /でプロンプトを表し，     
%       \# 以降はコメントを表す．
%        %以下，表示がインデントされすぎるのでインデントを減らす．
%   \begin{lstlisting}
%   % cd #  ホームディレクトリへ移動
%   \end{lstlisting}
%   \begin{lstlisting}
%   % cd .. #  一つ上のディレクトリへ移動 
%   \end{lstlisting}
%   \begin{lstlisting}
%   % pwd  #  現在のディレクトリの名前の表示
%   \end{lstlisting}
%   /Users
%   \begin{lstlisting}
%   % ls   #  現在のディレクトリにあるディレクトリやファイル表示
%   \end{lstlisting}
%   Shared   sonw00two
%   \begin{lstlisting}
%   % ls -a  #  不可視なものも含めたディレクトリやファイル表示
%      localized   Shared   sonw00two
%   % cd sonw00two  # このディレクトリのディレクトリ sonw00two へ移動
%   \end{lstlisting}\begin{lstlisting}
%   % mkdir tmp  # ディレクトリ tmp の作成
%   \end{lstlisting}
%   \begin{lstlisting}
%   % cd tmp  # ディレクトリ tmp へ移動
%   \end{lstlisting}
%   \begin{lstlisting}
%   % ps aux >ps.txt  # 全てのプロセスを ps.txt に書く
%   \end{lstlisting}
%   \begin{lstlisting}
%   % more ps.txt  # ps.txt  を表示．q で表示を閉じる
%   \end{lstlisting}
%   \begin{lstlisting}
%   % mv ps.txt ps1.txt  #  ps.txt の名前をps1.txt へ
%   \end{lstlisting}
%   \begin{lstlisting}
%   % cp ps1.txt ps2.txt  # ps1.txt  をps2.txt にコピー
%   \end{lstlisting}
%   \begin{lstlisting}
%   % rm *  # このディレクトリの任意のファイルを削除．*   はこのディレクトリの任意のファイル名を表すいわゆるワイルドカード．また，ファイルを全部消さないとディレクトリは削除できない．
%   \end{lstlisting}
%   \begin{lstlisting}z
%   % cd ..  
%   \end{lstlisting}
%   \begin{lstlisting}
%   % rmdir tmp  # ディレクトリ tmp の削除 > rm -R tmp で中のファイルを消さなくても削除可能ではあるが慣れるまでは確認してからの削除を勧める．
%
%   \end{lstlisting}
%\emptyline
%
%shellには全てをキーインしなくても，ファイル名などをある程度打ってタブキーを打つとあとの部分を補完してくれる機能がある．他の開発環境でも異なるキーかもしれないが大抵補間機能がついている．このファイルの最後の部分によく使うコマンドのリストを上げたので参考にしてほしい．

\section{標準入出力とパイプ}
最初難しく感じるのが標準入力・標準出力・標準エラー出力と
ファイルディスクリプタであると思うので解説する．
\begin{verbatim}
   % echo "example text"
   example text
\end{verbatim}
この例では"example text"標準入力で次の行が標準出力である．
\begin{verbatim}
   % echo "example text" > text
   %
\end{verbatim}
この例では標準出力がファイルtext に出力され，そのせいでターミナルには出力されない．\verb/>/をリダイレクトという．確かめてみると
\begin{verbatim}
   % cat text
   example text
\end{verbatim}
となり，ファイルtextの内容が標準入力となり，次の行が標準出力である．
\begin{verbatim}
   % cat test
   cat: test: No such file or directory
\end{verbatim}
この例では間違えた名前を打ったので，そのようなファイルはないと
標準エラー出力に表示された．つまり，標準エラー出力と標準出力は
ファオルトではターミナルの次の行である．

次にパイプについて説明する．
\begin{verbatim}
   % cat text | wc
       1       2      13
\end{verbatim}
この\verb/|/がパイプである．前半のコマンドcat textの標準出力が直接
wcの標準入力に送られその標準出力が次の行に出力されている．

入出力の管理はファイルディスクリプターによって制御されている．
これは比較的小さな非負整数で，標準入力は$0$,
標準出力は$1$, 標準エラー出力$2$となっていいる．
新しいデバイスを開いたりすると追加の非負整数が割り振られる．
これを使ってShell は細かな入出力の管理をしている．

\section{制御構文}
・for Loop : 

例：
for i in 1 2 3 4; do echo ch\$i ;done

for i in {1..9}; do echo ch\$i ;done
 
 >for i in {34..47}

comand line:
\begin{verbatim}
for> do echo 
'' > sample05\$i.py  
for> done
\end{verbatim}

・If 文:

If []
then

else

fi


\section{Shell Script}
Shell の操作をまとめてScriptを作ることが出来る．
これをShell Scriptという．まず
\begin{verbatim}
   #!/bin/sh
   
   echo "Hello World!"
\end{verbatim}
という内容を hello という名前で保存する．
そのディレクトリで
\begin{verbatim}
   % ls -l
   total 8
   -rw-r--r--  1 snow00two  staff  33 Apr 30 13:28 hello
   # このままだと実行可能でない
   %  chmod +x hello   # 実行可能に変換
   %  ./hello
   Hello World!
\end{verbatim}

\begin{Q}
   コマンドラインを使い，ディレクトリーを移動したり，
   ディレクトリーやファイルの作成や削除を行って下さい．
\end{Q}

% \newpage
%   \item
%      正規表現：パターンマッチングの表現．
%      どれでも良いのでよく使う言語の正規表現をマスターすることをお奨めする．
%      すると少しの変更で他の言語でも使えるようになる．

\begin{table}[h]
   \centering
   \begin{tabular}{ll}\noalign{\hrule height 1pt} 
   pwd&  現在のディレクトリー表示\\ \hline 
   ls&  現在のディレクトリーの中のファイル表示\\ \hline 
   cd dir1&  ディレクトリー dir1 へ移動\\ \hline 
   cd & Home ディレクトリーへ移動\\ \hline 
   cp file1 file2& file1 をfile2 にコピー\\ \hline 
   mv file1 file2&  file1 の名前をfile2 に変更\\ \hline 
   rm file1& file1 の削除 \\ \hline 
   touch file1 & file1 の変更時刻の更新，file1 が無ければ file1 を作成\\ \hline 
   mkdir dir1&  ディレクトリー dir1 を作る\\ \hline 
   rmdir dir1&  ディレクトリー dir1を削除\\ \hline 
   ps&  現在のプロセスを表示\\ \hline 
   kill pid &  プロセスid (自然数)pid のプロセスの終了\\ \hline 
   which cmd&  cmd が使えればそのコマンドの pathを表示\\ \hline 
   who&  誰が今ログインしてるかを表示\\ \hline 
   who am i&  自分がどのusername でログインしているか表示\\ \hline 
   date&  現在の日時を表示\\ \hline 
   cal &  現在の月のカレンダーを表示\\ \hline 
   chmod +x file1&   file1をコマンドとして実行可能にする\\ \hline 
   ./file1&   現在のディレクトリの file1をコマンドとして実行\\ \hline 
   chmod -x file1&   file1をコマンドとして実行不可能にする\\ \hline 
   echo "text" &  文字列 text を表示\\ \hline 
   cat file1 file2 &  file1, file2の内容を連結(concatenate)して表示\\ \hline 
   pr file1 &   file1 の内容をフォーマットして表示\\ \hline 
   tail -n  file1&  file1の最後の$n$行を表示\\ \hline 
   head -n file1 & file1の最初の$n$行を表示 \\ \hline 
   more file1&   file1 の内容を確認しながら表示\\       \noalign{\hrule height 1pt}
   \end{tabular}
   \smallemptyline
   \caption{コマンドの基本}
\end{table}

\begin{table}[htb]
   \centering
   \begin{tabular}{ll}\noalign{\hrule height 1pt} 
   ed file1&  ed でfile1を編集\\ \hline 
   vim file1&  vim でfile1を編集\\ \hline 
   code file1&  VSCode でfile1を編集\\ \hline 
   open -a apli file1& アプリ apli で file1を編集\\ \hline 
   sh&  subshell を開く\\ \hline 
   wc file1&  file1 の行数，語数，文字数を表示\\ \hline 
   cmp file1 file2&  file1 と file2 が違えば最初の違う行を表示\\ \hline 
   diff file1 file2&  file1 と file2 がどのように違うかを表示\\ \hline 
   grep hoge file1&  file1のhogeを含む行を表示\\ \hline 
   sort file1&   file1の行毎にソートする\\ \hline
   man cmd1& cmd1 のマニュアル表示\\ \hline
      brew install hoge & brew で tool hoge をインストール\\ \hline
      history & これまで使ったコマンドのリスト\\ \hline
      printenv  &全ての環境変数の表示\\ \hline
      echo \verb/$/ENV  &環境変数ENV(が定義されていれば)表示\\ \hline
      su hoge&  hoge というusernameのユーザに変わる\\ \hline
      su &  root に変わる．Macでは廃止\\ \hline
      sudo cmd1 & スーパーユーザー権限でcmd1 を実行\\
      \noalign{\hrule height 1pt}
   \end{tabular}
   \smallemptyline
   \caption{コマンドの基本2}
\end{table}


\begin{table}[htb]
   \centering
   \begin{tabular}{ll}\noalign{\hrule height 1pt} 
      control  D & end of file \\ \hline
      control  S & log out\\  \hline
      tab & filename 補間 \\ \hline
      \verb/*/ & 任意filename \\ \hline
      \verb/*/hoge & 最後がhoge の任意filename \\ \hline
      \verb/!!/ &直前のコマンドをもう一度 \\ \hline
       \verb/!/cmd1 & cmd1で始まる一番新しいコマンドをもう一度 \\ \hline
     cmd1\verb/;/ cmd2& コマンド cmd1と cmd2の並列 \\ \hline
      cmd1\verb/|/ cmd2& pipe:コマンドcmd1の標準出力を
      cmd2の標準入力へ \\ \hline
      \verb/>/ file1& file1へredirect \\ \hline 
      \verb/>>/ file1 & file1へ追加 \\ \hline 
     \verb/</ file1& file1からredirect \\ \hline
    \verb/./ & current directory \\ \hline
     \verb/../ & parent directory \\
     \noalign{\hrule height 1pt} 
   \end{tabular}
   \smallemptyline
   \caption{terminalとshellのメタキャラクター}
\end{table}

\begin{table}[htb]
   \centering
   \begin{tabular}{ll}\noalign{\hrule height 1pt} 
      control A& 行の先頭へ\\ \hline
      control B & 一文字前へ\\ \hline
      control E & 行の最後へ\\ \hline
      control F& 一文字後ろへ\\ \hline
      control N& 次の行へ\\ \hline
      control P & 前の行へ\\ \hline
      control D & 右を消去\\ \hline
      control K & カーソルの後を消去\\ \hline
      control H & 左を消去\\ \hline
      control V & 一画面移動\\ \hline
      control Y & 再入力\\ \noalign{\hrule height 1pt}
   \end{tabular}
   \smallemptyline
   \caption{Emacs由来}
\end{table}



%\end{document}


%\documentclass[a4paper,12pt]{amsart}
%\usepackage[dvipdfmx]{graphicx,xcolor}%graphics
%\usepackage{float}
%\usepackage{amsmath}
%\usepackage{verbatim}
%\usepackage{latexsym}
%\usepackage{amssymb}
%%\usepackage{emath}
%\usepackage{listings}
%\usepackage{ccicons}
%\usepackage{mflogo}
%%\usepackage{fancybox}%enclose by box
%\usepackage[dvipdfmx]{hyperref}
%\usepackage{verbatim}
%\hypersetup{% hyperrefオプションリスト
%   setpagesize=false,
%   bookmarksnumbered=true,%
%   bookmarksopen=true,%
%   colorlinks=true,%
%   linkcolor=red,
%   citecolor=blue
%} 
%
%\newtheorem*{Thm}{定理}
%\newtheorem*{Prop}{命題}
%\newtheorem*{exQ}{レポート}
%
%\theoremstyle{remark}
%\newtheorem*{ex}{例}
%\newtheorem*{Q}{演習}
%\newtheorem*{Quiz}{問題}
%
%\renewcommand{\figurename}{図}
%\renewcommand{\proofname}{証明}
%\renewcommand{\tablename}{表}
%
%\def\emptyline{\vspace{11pt}}
%\def\smallemptyline{\vspace{6pt}}
%
%\begin{document}
%\begin{flushright}
%   2026/08/01
%      
%   snow00two
%\end{flushright}
%
%\begin{center}
%   {\Large Git と GitHub}
%   \href{https://creativecommons.org/licenses/by-nc-nd/4.0/}{\ccbyncnd}
%\end{center}

\chapter{Git と GitHub}

\section{バージョン管理の必要性}
\textcolor{red}{バージョン管理}とはあるプロジェクトのファイルやディレクトリーを
いつでも元の状態に戻れるように状態を保ちながら
編集し続けることをいう．
例えば事務的なExcel書類 hoge.xlsmを

\begin{verbatim}
   hoge7_4.xlsm
\end{verbatim}

\noindent
のような形で最終変更日を付けて保存したりしたことがあると思う．
しかし，複数の人が同じファイルを扱う場合（自分一人の場合でも），
2つのファイル
\begin{verbatim}
   hoge7_6.xlsm
   hoge7_7.xlsm
\end{verbatim}
が\verb/hoge2026_7_4.xlsm/から
どの順番に生成されたか分からなくなったような経験のある人も多いだろう．

\begin{figure}[htb]
   \centering
   \includegraphics{~/Documents/GitHub/introInfo/appe/images/figure.1}
   \caption{どっち？}\label{}
\end{figure}
   
誰が何時手を入れて，
それがどのように統合されたのか履歴と責任を明らかにし，
そのサポートを行うのがバージョン管理ソフトの仕事である．
もちろん，使う側にもその使い方の理解が欠かせない．
\textcolor{red}{Git }はそのようなバージョン管理をするためのツールであり，
最近ではバージョン管理のデファクトスタンダードである．

Git ではブランチという枝の様な構造を使ってバージョン管理を
行っている．
新しいことを付け加えるときはまず幹に当たるmain ブランチから
新しいブランチを切ってそこで作業し，
一段落付いたらmain ブランチに統合する，あるいは他の人間にレビューしてもらってその後統合する
という方針で作業を行う．

\textcolor{red}{GitHub}というクラウドサーバーに
適宜ファイルを保存し，
ネット上の仲間と共同でバージョン管理をするのが
今の開発の主流である．

\section{バージョン管理ツールの歴史}
%\begin{table}[h]
%	\begin{tabular}{llll} 
%  		\noalign{\hrule height 1pt}
%		 	名前&  &  &  \\ \hline \hline
%		 	CSSV&  &  &  \\ \hline
%			CVS& &  &   \\ \hline
%			Subversion& &  &  \\ \hline
%			Git & &  &  \\ \hline
%  		\noalign{\hrule height 1pt}
%	\end{tabular}
%	\caption{バージョン管理ツールの歴史}
%\end{table}

\begin{itemize}
   \item
      SCCS 1972年IBM/System 370向けに開発された．ローカル型．
   \item
      RCS 1982年，ローカル型．
   \item
      CVS 1986年，クライアントサーバー型．
   \item
      Subversion  2000年，クライアントサーバー型．
   \item
      Bitkeeper  2000年，分散型．商用ソフトでLinux の開発でフリーで使われていた．
   \item
      Git 2005年，分散型．突然 Bitkeeper が使えなくなったリーナス・トーバルズが開発を始めた．
\end{itemize}


\section{Gitのインストール}
Mac やLinux だと既にgit がインストールされてるので，ターミナルで
\begin{verbatim}
   % git --version
   git version 2.54.0
\end{verbatim}
のようにすれば確認出来る．
ただし\verb/% /はプロンプトで入力する必要は無い．
必要に応じて次の行に出力も書いておく．
Windows だとインストールが必要．

Git の初期化：まずGitに自分のことを教える必要がある．
\begin{verbatim}
   % git config --global user.name "hoge"
   % git config --global user.mail "hoge@example.com"
\end{verbatim}
後のGitHubの登録後に本人確認のメールとか届くので，
本物のメールアドレスを登録すること．
細かい設定も似た感じで出来る．例えば先程言った幹に当たるブランチの名前を
main にするには
\begin{verbatim}
   % git config --global init.defaultBranch main
\end{verbatim}
とすれば良い（以前はmaster という名前だったが奴隷制を思い出させるということで
main のような名前で呼ばれることが多くなった）．なお
\begin{verbatim}
   % git config --list
\end{verbatim}
でちゃんと設定できたか確認出来る．

\section{Git の最低限の基本操作}
Gitではまずはコマンドラインを使って操作し，
使いながらGitにおけるバージョン管理の概念・使用法の
理解を深めていくことをお奨めする．
分かってきたらグラフィカルな方法に移って構わない．

\begin{itemize}
   \item
      Git はディレクトリ単位でバージョン管理するので，
      まずディレクトリを初期化する必要がある
\begin{verbatim}
   % mkdir hogehoge
   % cd hogehoge   
   % git init 
\end{verbatim}
      新しいファイルを作って最初のコミットを行う．
      最初は説明を省くが後で説明を加える．
\begin{verbatim}
   % echo "Hello Git!" > textA.txt
\end{verbatim}
\begin{verbatim}
   % git add textA.txt
\end{verbatim}
\begin{verbatim}
   % git commit -m "first commit"
\end{verbatim}

   \item
      git の状態を表示
\begin{verbatim}
   % git status 
   On branch main
   nothing to commit, working tree clean 
\end{verbatim}

   \item
      新しいことをするには，main ブランチでは作業を行わず，
      例えばbranchAというブランチを作り，そこに移って作業する．
\begin{verbatim}
   % git branch branchA
   % git switch branchA
\end{verbatim}
      この段階で書類をVSCodeなどで編集すると，
      このブランチ branchAの上で編集されるだけで，
      main ブランチの内容はもとのままに止まっているが，
      現在の編集内容も保存されていない．

      例えば適当にファイルを変更して
\begin{verbatim}
   % echo "add a comment" >>  textA.txt 
\end{verbatim}

   \item
      編集内容を保存するため，ファイル textA.txtをステージング
\begin{verbatim}
   % git add textA.txt 
\end{verbatim}

\noindent
      し，問題なければコミットして作業中の内容を保存．
\begin{verbatim}
   % git commit -m "edit textA"
\end{verbatim}
      \verb/"edit textA"/ がそのコミットのコメントで，
      後で確認出来るように必ず付け加えることになっている．
      もし git commitのみだとエディターが自動的に立ち上がって，
      入力を要求される．
      
      以後同じことを繰り返す．


\begin{figure}[htb]
   \centering
   \includegraphics{~/Documents/GitHub/introInfo/appe/images/figure.2}
   \caption{作業 branchとmain branch}\label{}
\end{figure}

   \item
      幾つものファイルを同時に編集しても構わないし，
      何度かブランチ branchA上でコミットを続けることも出来るので
      適宜部分的に保存しながら作業できる．
   
   \item
      branchAブランチの変更した内容をmainブランチに付け加えることを
      マージというが，
      特に一人で開発しているときはブランチ一個で作業し
      \textcolor{red}{ストレートフォワードマージ}といい，
      何も手を入れずそのままmain にマージすることが多いだろう．
      
      やり方を説明すると，一回メインブランチに移動して
\begin{verbatim}
   % git switch main
   % git merge branchA
\end{verbatim}
%これは次と同値
%\begin{verbatim}
   % git merge origin/main 
%\end{verbatim}


\begin{figure}[htb]
   \centering
   \includegraphics{~/Documents/GitHub/introInfo/appe/images/figure.3}
   \caption{ branchAと Main branch とコミット状況}\label{}
\end{figure}
      
      大事なことなので強調しておくと，
      Git ではmain で編集作業をすることは推奨されていない．
      新たな編集作業は新たなブランチで行い，
      切りの良いときmainに統合することを繰り返すことが
      奨励されている．

\end{itemize}
\begin{Q}
   以上のことをターミナルで行なってください．
\end{Q}


\section{Git のやや詳しい操作}
以下もう少し詳しいやり方を説明する：
\begin{itemize}
%\begin{verbatim}
%   % git switch -d
%\end{verbatim}
   
   \item
      ブランチの関係を見る
\begin{verbatim}
   % git log
\end{verbatim}

      \noindent
      情報が多すぎるので各項目を一行で表す
\begin{verbatim}
   % git log --oneline
   011784d (HEAD -> main, origin/main, testSite) add site data
   c2b3d4f first commit
\end{verbatim}
      の方が見やすいかもしれない．
      最初の数字がコミットハッシュ(の一部)で，この数字でコミットを制御している．
      関係をグラフで表すことも出来る
\begin{verbatim}
   % git log --graph --oneline
\end{verbatim}

   \item
      リポジトリーの全てのブランチ名の表示
\begin{verbatim}
   % git branch
\end{verbatim}

   \item
      コッミットハッシュc1f4522のブランチに移動
\begin{verbatim}
   % git switch --detach c1f4522
\end{verbatim}

   \item
      ブランチ\verb/<branchName>/の削除
\begin{verbatim}
   % git switch -d <branchName>
\end{verbatim}
      適宜不必要なブランチを消して整理することをお奨めする．

   \item
      ステージングの取消し
\begin{verbatim}
   % git rm --cached <fileName>   
\end{verbatim}

   \item
       全ての変更ファイルをステージング：
\begin{verbatim}
   % git add -A
\end{verbatim}

   \item
       カレントディレクトリの全ての変更ファイルをステージング：
\begin{verbatim}
   % git add .
\end{verbatim}

   \item
      コミットしたコメントを間違えたときの変更の仕方：
\begin{verbatim}
   % git commit --amend -m “amend text”
\end{verbatim}

   \item
      最後のマージとの変更点表示
\begin{verbatim}
   % git diff  
\end{verbatim}
      
   \item
      \verb/<fileName>/の変更点表示
\begin{verbatim}
   % git diff -- <fileName> 
\end{verbatim}
      %他にも色々な使い方があるのでマニュアル等で調べてください．

   \item
      途中を消すには
\begin{verbatim}
   % git merge --squash <branchName>
\end{verbatim}

   \item
      merge で衝突が起こった場合の解消
\begin{verbatim}
   % git merge --abort
\end{verbatim}
      で一回中止し，編集して衝突を除く．

   \item
      マージの他にrebaseというのもある 
\begin{verbatim}
   % git rebase
\end{verbatim}

   \item
      mainで間違えて作業したときの対応：
      まず変更点を保存して
\begin{verbatim}
   % git stash
\end{verbatim}
      ブランチを切って
\begin{verbatim}
   % git switch -c branch
\end{verbatim}
      そのbranch に変更点をアップ
 \begin{verbatim}
   % git stash pop
\end{verbatim}
\end{itemize}
まだ他にも沢山のコマンドがあるので，参考文献などで調べてください．


\section{GitHub}\label{GitHub}
ファイルサーバーGitHubを使うとGitの有用性が広がる．
ファイルのクラウドへの保存と
ネット上でのファイルの共有とかが重要である．
個人ではフリーで使用できる．

\begin{itemize}
   \item
      \href{https://github.co.jp/}{GitHub} を訪れアカウントを作成
   \item
      GitHub の自分のページにサインイン

   \item 
      GitHub の自分のアカウントにssh で公開鍵を保存：
      まず自分のパソコンのログインディレクトリーで
      以下のように公開鍵を作成する．
\begin{verbatim}
   % cd  # ホームディレクトリに移動
   % ssh-keygen # RSA で秘密鍵と公開鍵の作成
\end{verbatim}
      途中でパスフレーズを2回聞かれるので（空白でも作れるが）
      絶対入力するように．
      上手くいくと
\begin{verbatim}
   % ls -a
   ...
   .ssh 
   ...
\end{verbatim}
      上のように鍵が作られている．
      その後, .sshのなかの公開鍵をGitHubの自分のホームに
      登録．

\begin{enumerate}
   \item
      GitHub のHome の一番右上のアイコンのボタンをクリック(図~\ref{GitHubHome}だと矢印の部分，ファビコンと言うらしい)．
\begin{figure}[htb]
   \centering
   \includegraphics[width=12cm]{~/Documents/GitHub/introInfo/appe/images/GitHub.pdf}
   \caption{GitHubのHome}\label{GitHubHome}
\end{figure}

   \item
      左の Setting のボタンをクリック
   
   \item
      SSH and GPG Key のボタンをクリック
   
   \item
      SSH の New Key のボタンをクリック
   
   \item
      Title と Key を入力
   \end{enumerate}
   
   \item
      GitHubで Repositoryの New ボタンをクリック
      Repository Name に(適当な名前)fugaを記入．
      他は省略してCreate Repository のボタンをクリック
      するとリポジトリー が作成され，そこには以下のような指示がある：

\begin{verbatim}
   …or create a new repository on the command line

   echo "# fuga" >> README.md
   git init
   git add README.md
   git commit -m "first commit"
   git branch -M main
   git remote add origin git@github.com:snow00two/fuga.git
   git push -u origin main
\end{verbatim}

      \noindent
      これに従ってコマンドを実行する．以下やや詳しく解説：
     
      ローカルでディレクトリfugaを作成し，そこに移動し
      最低限必要なREADME.mdを作成
\begin{verbatim}
   % mkdir fuga
   % cd fuga
   % echo "# fuga" >> README.md
\end{verbatim}

      そのディレクトリの初期化：
\begin{verbatim}
   % git init
\end{verbatim}

      README.mdをステージングしてコミット
\begin{verbatim}
   % git add README.md
\end{verbatim}
\begin{verbatim}
   % git commit -m "first commit"
\end{verbatim}
      Main ブランチを切り設定
\begin{verbatim}
   % git branch -M main
   % git remote add origin git@github.com:snow00two/fuga.git
\end{verbatim}
      最初のプッシュ：
\begin{verbatim}
   % git push -u origin main
\end{verbatim}
      ここでパスフレーズを聞かれて，正しければGitHubにデータが送られる．
 
\begin{figure}[htb]
   \centering
   \includegraphics{~/Documents/GitHub/introInfo/appe/images/figure.4}
   \caption{GitHub, push, pull}\label{}
\end{figure}

   \item
      以後 main ブランチの新しいデータをGitHub のにアップするのは
\begin{verbatim}
   % git push 
\end{verbatim}
      で良い．

   \item
      GitHub からデータを取ってくるには
\begin{verbatim}
   % git pull 
\end{verbatim}
      これは以下と同じ
\begin{verbatim}
   % git fetch 
   % git merge 
\end{verbatim}
      一つづつ確認して作業を行いたい場合はこちらをお奨めする．
      
      なお最初にGitHub からデータを取ってくるには，例えば
 \begin{verbatim}
   % mkdir fractalApp
   % cd fractalApp
   % git clone https://github.com/snow00two/fractalApp.git
\end{verbatim}
      のようにすれば良い．    
      
\end{itemize}

\begin{Q}
   以上のことをGitHubとターミナルで行なってください．
\end{Q}

\section{VSCode と Git/GitHub}
Git/GitHub と VSCodeは相性が良いので，
VSCode で Git/GitHub を使うことをお奨めする．
\begin{itemize}
   \item
      VSCode で Git で管理されているディレクトリーを開くと
      自動でGit の管理を感知して
      ターミナルはそのディレクトリーをカレントディレクトリーとして
      開いてくれる．また．以前使ったコマンドのHistory の復元を行ってくれ，
      コマンドラインから Git やGitHubを管理し易い環境になっている．
      ファイルの編集自体はVSCodeの編集画面を使えば良い．
      
   \item
      VSCode の設定を GitHubで同期するように設定できる．
      
   \item
      GitHub のStudent Developer Pack\footnote{2026/4 から新規受付を停止していた．
      いろいろ変わるので詳しいことはネットで確かめてください．} に入っていると，
      VScode の拡張機能タブで「GitHub Copilot」を検索してインストールすると，
      使用料無しでVScodeから生成AI の GitHub Copilot が使えるようになり，
      生成AI にコーディングの補助を受けながら
      コーディングすることが出来て便利である．
\end{itemize}


\section{GitHub Pages}
GitHubではPagesというサービスを使って
自分のホームページを世界中に公開できる．

\begin{itemize}
   \item
      \S~\ref{GitHub} で作ったローカルのRepository  fuga に
      docsというディレクトリを作り，
      その中にホームページとして実現したい
      html, css, js のファイルを準備する．markdown でも十分で
\begin{verbatim}
   % git branch pages
   % git switch pages
   % mkdir src docs
   % cd src 
   % echo "# This is my test site." > index.md
   % pandoc -f markdown -t html -s -o ../docs/index.html index.md
   % git add -A
   % git commit -m "add site data"
   % git switch main
   % git merge pages
   % git push
\end{verbatim}
   \item
      GitHubのホームページのRepository fuga の上の段の
      Setting(見えないときはMoreの下に隠れている)を選び，
      左のコラムのPages を選択し，
      セレクターからDeploy from a branch を選び
      その下のBranchでセレクターから main と /docs 選ぶ．
   
   \item
      暫くするとSetting でPages の上の部分に

Your site is live at https://snow00two.github.io/fuga/
      
      \noindent
      のようなアドレスが現れると思う．
      そこをクリックするとあなたのホームページが表示される(図~\ref{pagesTest})．markdownについては
      \href{https://snow00two.github.io/docs/introductionTex.pdf}{latex と markdown}を参照．
 \begin{figure}[htb]
   \centering
   \includegraphics[width=11cm]{~/Documents/GitHub/introInfo/appe/images/pagesTest.pdf}
   \caption{作られたホームページ}\label{pagesTest}
\end{figure}     
      
      ただもう少し凝ったホームページを作りたい場合これでは不十分なので，
      \cite{M}や動画などを参考にしてHTMLやCSSやデザインの基礎を学んで欲しい．


\end{itemize}
\begin{Q}
   以上のことをGitHubで行なってください．
\end{Q}
\begin{exQ}
   皆さんの作ったアプリや動画や論文を世界に向けて発信
   するウェブサイトを作って公開し，
   そのアドレスを私の所に送って下さい．
\end{exQ}


\section{この先に}
多人数での開発のためにはPull request と Issueが必要になるが，
これについては参考文献などを調べてほしい．
また，マージするときに衝突が起こる場合があるので
その場合は手動でマージする必要がある．
そのやり方も参考文献などを調べてほしい．

また，Actions というGitHubの自動化のサービスもあるので調べて活用してほしい．


%\begin{thebibliography}{99}
%   \bibitem{S}  \href{https://backlog.com/ja/git-tutorial/}{サル先生のGit入門}
%
%   \bibitem{T}  高橋 あおい, つくって、壊して、直して学ぶ Git＆GitHub 入門, 翔泳社, 2025.
%   
%   \bibitem{W}  渡辺宙志, ゼロから学ぶGit/GitHub, 講談社, 2024.
%
%   \bibitem{G} 
%   ポンストール, ロリガー, 実用 Git 第3版, オライリー, 2024.
%   
%   \bibitem{M}
%   Mana, 1冊ですべて身につくHTML \verb/&/ CSSとWebデザイン入門講座［第2版］, SBクリエイティブ, 2024.
%\end{thebibliography}
%
%\end{document}


\backmatter


\begin{thebibliography}{99}
    \bibitem{Bo}
      Bool, {\em Laws of thought}, 1848
    \bibitem{BF} 
      ボズウェル, フィッチャー, リーダブルコード
   \bibitem{Ba} 
         ジェイミー・バートレット, 
         操られる民主主義: デジタル・テクノロジーはいかにして社会を破壊するか, 
         草思社, 2018

   \bibitem{Br} 
      ブレーザー, 失われてゆく我々の内なる細菌, みすず書店, 2015.

   \bibitem{C}    
      チャルディーニ：影響力の武器

   \bibitem{F}
      Feynman, Lectures on computation, Perseus, 1996, 
      日本語版：ファインマン計算機科学, 岩波書店, 1999.


   \bibitem{Fu}
      藤原進之介, 
      学校で習っていなくても読んで理解できる藤原進之介のゼロから始める情報 I, 
      KADOKAWA, 2023. 

   \bibitem{GM}
       Joe Gray, Jin Maeda,
      人を動かすハッカーの技術：ソーシャルエンジニアリングの実践と防御 

   \bibitem{history}
      ゴールドスタイン, 復刊計算機の歴史, 共立, 2016.

   \bibitem{Ho} 
      堀田 良之, 加群十話: 代数学入門 (すうがくぶっくす 3), 朝倉書店, 1988.
 
   \bibitem{Han} 
      アンデシュ・ハンセン, スマホ脳, 新潮新書, 2020.
   \bibitem{H} 
      ユバル・ノラ・ハラリ, サピエンス全史
      
    \bibitem{HB}
      NISC,
      \href{https://security-portal.cyber.go.jp/guidance/handbook.html}
      {インターネット安全・安心ハンドブック} 
   
   \bibitem{Ha} 
      林 裕之, 「隠れ油」という大問題, 三五館, 2017.
      
   \bibitem{Hi}
      平田 賀一, 
      ITエンジニアのためのMarkdown実践入門 ── 生成AI時代の"伝わる"書き方, 
      技術評論社, 2026.

   \bibitem{I} 
      石井光太, ルポ スマホ育児が子どもを壊す, 新潮社, 2024.
   
   \bibitem{interface}
      Apple Computer, 
      Apple Human Interface Guidelines: 
      The Apple Desktop Interface, Addison-Wesley, 1987,
      日本語版, 1989.


  \bibitem{IT}
      渡辺さき［令和8年度］ITパスポート 超効率の教科書＋よく出る問題集, インプレス, 2026.
    \bibitem{IPA}
       情報処理推進機構, 
       \href{https://www.ipa.go.jp/security/vuln/websecurity/about.html}{安全なウェブサイトの作り方}

    \bibitem{IPA2}
       情報処理推進機構, 
      \href{https://www.meti.go.jp/policy/netsecurity/downloadfiles/guide_v3.0.pdf}{サイバーセキュリティ経営ガイドライン}

    \bibitem{IPA3}
       情報処理推進機構, 
     \href{https://www.ipa.go.jp/security/guide/sme/ug65p90000019cbk-att/sme_guideline_v4.0.pdf}{中小企業の情報セキュリティガイドライン}
    \bibitem{JPCERT} 
       一 般社団法人 JPCERT コーディネーションセンター, 
       \href{https://www.jpcert.or.jp/csirt_material/files/manual_ver1.0_20211130.pdf}{インシデントハンドリングマニュアル}

  \bibitem{Jobs}
      アイザックソン, スティーブ・ジョブスI, II, 講談社, 2011.

   \bibitem{K}  
      カーニハン, 
      教養としてのコンピューターサイエンス講義 第二版, 日経BP, 2020.
   \bibitem{Ka}
      木村秀紀, Fourier-Laplace 解析, 岩波書店, 1993.
 \bibitem{Kabuki}
      小学館DVDBook, 歌舞伎のいき, 第2巻「時代物・荒事編」, 2008. 
     \bibitem{UNIX}
      カーニハン,
      カーニハンのUNIX回顧録, 丸善, 2025.
    \bibitem{KP}
      Kernighan, Pike, 
      Unix programming environment, Pearson, 1984
      (訳 UNIX プログラミング環境, KADOKAWA, 2017).
   \bibitem{KR}
      Kernighan, Ritchie, 
      The C programming  Language--2nd edition, 
      Prentice Hall,1989. 
      (訳プログラミング言語C 第2版 ANSI規格準拠, 共立出版, 1989.)

    \bibitem{Kay}
      アラン・ケイ, アラン・ケイ, ASCII, 1992.
   \bibitem{Ko} Koblitz, A course in Number theory and Cryptography, GTM,  
   \href{https://link.springer.com/search?package=11649&facet-content-type=%22Book%22}{ダウンロード}

   \bibitem{Ko}
      小森 裕介,［ 改訂新版］プロになるためのWeb技術入門, 
      技術評論社, 2024.      
 \bibitem{K} 
      川本 晃司, スマホ失明, かんき出版, 2022.
   \bibitem{KS}
       川島 隆太, 榊 浩平, スマホはどこまで脳を壊すか, 朝日新書,  2023.
 
    \bibitem{Kum}
   熊野 整, 外資系投資銀行のExcel仕事術, ダイヤモンド社, 2015.
 
    \bibitem{Kud}
   工藤 喜美枝, 入門!Excel VBAクイックリファレンス 改訂版, ムイスリ出版, 2018.  
 
   \bibitem{KL} 
      ヴァイバー・クリガン=リード, 
      サピエンス異変──新たな時代「人新世」の衝撃, 
      飛鳥新社, 2018.
  \bibitem{PK} 
      カール・ポッパー, 科学的発見の論理 上下, 恒星社厚生閣, 1971.
   
   
   \bibitem{L} 
      ダニエル・リーバーマン, 人体六〇〇万年史 
      上・下──科学が明かす進化・健康・疾病, 早川書房, 2015.
   
   \bibitem{LL} 
      ダニエル・リーバーマン, マイケル・ロング , 
      もっと! : 愛と創造，支配と進歩をもたらすドーパミンの最新脳科学, 
   インターシフト (合同出版), 2020.

   \bibitem{LT} 
      ティム・バーナーズ=リー, Webの創成 — 
      World Wide Webはいかにして生まれどこに向かうのか, 
      毎日コミュニケーションズ, 2001.

   \bibitem{LLM}
   Berryman, Zigler, LLMのプロンプトエンジニアリング, オライリージャパン, 2025.


   \bibitem{Mn} Minsky, Computation: Finite and Infinite Machines, 1967.
   
   \bibitem{Mu}三浦伸夫, 数学の歴史, 放送大学教育振興会, 2019.

   \bibitem{M} 
   マロット, 行動分析学入門, 産業図書, 1998.

   
   \bibitem{M1} 
      物江 潤, デマ・陰謀論・カルト (新潮新書) , 新潮社, 2022.
   
   \bibitem{M2} 
      物江 潤, デジタル教育という幻想 GIGAスクール構想の過ち, 
      平凡社, 2023.

   \bibitem{P5}
      松田晃一, p5.jsプログラミングガイド  改訂版, 2021, カットシステム
 
     \bibitem{N} 
    一般社団法人 日本シーサート協議会シーサート人材ワーキンググループ, 
    改訂新版 セキュリティエンジニアの教科書 一般社団法人 日本シーサート協議会シーサート人材ワーキンググループ  
  
   \bibitem{NKO} 
      長迫, 小谷, 大澤, SNS時代の戦略兵器 陰謀論 民主主義をむしばむ認知戦の脅威,
    ウェッジ, 2025.
 
      
    \bibitem{N}
      NISC, インターネット安全・安心ハンドブック
   \href{https://security-portal.cyber.go.jp/guidance/handbook.html}{ダウンロード}
    \bibitem{O} 
       小川和久, 総理、国防も安全も穴だらけ！国民を守れない国・ニッポン


   \bibitem{Orwell}
      H.G.~Orwell, Nineteen Eighty-Four, 1949.

    \bibitem{OM}
       奥村 晴彦, 森本 尚之,
       ［改訂第5版］基礎からわかる情報リテラシー, 技術評論社, 2023. 
    \bibitem{P} 
      イーライ・パリサー, 閉じこもるインターネット——
      グーグル・パーソナライズ・民主主義 , 早川書房, 2012.
      改題：フィルターバブル──インターネットが隠していること
      (ハヤカワ文庫NF), 早川書房, 2016.

   \bibitem{PC}
      キャサリン・プライス,  
      スマホ断ち 30日でスマホ依存から抜け出す方法, 
      角川新書, 2024.
      
   \bibitem{tools}
      ラインゴールド, 思考のための道具, パーソナルメディア, 1987.
         
    \bibitem{R}   
       ガー・レイノルズ, プレゼンテーションZen 第3版, 丸善出版, 2021.
   \bibitem{Sa}
      坂村健, 高等学校 情報 I , 数研出版, 2021.   

   \bibitem{ST} Silverman, Tate,
   Rational Points on Elliptic Curves, Springer, 2015 
%   \href{https://link.springer.com/search?package=11649&facet-content-type=%22Book%22}{ダウンロード}
    \bibitem{SS}
    瀬戸, 斎藤, 徹底攻略情報セキュリティマネージメント, インプレス   
    \bibitem{Ta} 
      高木貞治, 初等整数論講義, 共立出版, 1971.
    \bibitem{TH} 
      達人プログラマー

    \bibitem{TK}
    高橋 佑磨, 片山 なつ, 
    伝わるデザインの基本 増補改訂3版 
    よい資料を作るためのレイアウトのルール, 技術評論社, 2021.
   
   
    \bibitem{T}
    田中靖浩, 会計の世界史, 日本経済新聞出版, 2018.

   \bibitem{TM}  
     井上 直也, 村山 公保, 
     マスタリングTCP/IP—入門編—(第6版), オーム社, 2019.
  

    \bibitem{TuringM}
      アンドルー・ホッジス, エニグマ アラン・チューリング伝 上下,
      勁草書房, 2015.
      映画化：イミテーション・ゲーム エニグマと天才数学者の秘密,
      モルテン・ティルドゥム監督, 2014.

      
    \bibitem{Turing}
       A.~Turing, On Computable Numbers, with an Application to the Entscheidungsproblem, Proceedings of the London Mathematical Society 42, 1936.
       

   \bibitem{Uts} 
   宇津木龍一, シャンプーをやめると髪が増える, KADOKAWA, 2013.
   
    \bibitem{Uti}
       内田 広由紀, 7日間でマスターする 配色基礎講座 
       (デザイン ビギナー シリーズ), 2000
 \bibitem{W}
     クジラ飛行机, プログラミング言語大全, 技術評論社, 2025.

    \bibitem{Wi}
       Robin Williams,  ノンデザイナーズ・デザインブック [第4版], 
       マイナビ出版, 2016
   \bibitem{Y}
      矢野 啓介, 
      [改訂新版］プログラマのための文字コード技術入門, 技術評論社, 2018.    
  \bibitem{Ya}
   山森丈範，[改訂第4版]シェルスクリプト 基本レファレンス, 技術評論社, 2024.

   \bibitem{Yamakawa}
      英文詳説世界史, 山川出版, 2019. 
    \bibitem{Yag}
      山口泰, 東京大学教養学部テキスト情報 第三版, 
      東京大学出版会, 2025.
     \bibitem{Yo}
      吉田 拳, たった1日で即戦力になるExcelの教科書, 技術評論社, 2014.
   \bibitem{EV}
   国本 温子, 緑川 吉行, できる大事典 Excel VBA 2019/2016/2013\& Microsoft 365, インプレス, 2021.

   \bibitem{EL}
   尾崎裕子, 日花弘子, できる大事典 Excel 2016, インプレス, 2021.
   
   \bibitem{EF}
   羽山博, 吉川明広, できる大事典 Excel 関数 2016/2013/2010, インプレス, 2017.
      
    \bibitem{YM}
    シェルスクリプト基本レファレンス

   \bibitem{Yuki}
      結城浩, 暗号技術入門 第３版, SBクリエイティブ, 2015.
   \bibitem{YY}
      山田祥寛, 
      改訂3版JavaScript本格入門, 2023, 技術評論社

%   \bibitem{B}Bool, {\em Laws of thought}, 1848
%    \bibitem{Ko} Koblitz, A course in Number theory and Cryptography, GTM,  \href{https://link.springer.com/search?package=11649&facet-content-type=%22Book%22}{ダウンロード}
  
   \bibitem{TeXShop}
   TeXShop のメニューバーの Helpメニューにある説明文書と動画
   
   \bibitem{texlive}
   /usr/local/texlive/2024/texmf-dist/doc/latex の中のドキュメントファイル:  
    User's guide for amsmath package, Using the amsthm package, 
    XY-pic User's Guide, XY-pic Reference Manual,
    The \MP\, system, A users manual for \MP など
   
   \bibitem{OK}
   奥村晴彦, 黒木裕介, 
   ［改訂第9版］\LaTeX  美文書作成入門, 技術評論社, 2023.
   
   \bibitem{LCs}
   {\em The  \LaTeX \, Companions Third Revised Boxed Set}, Addison-Wesley, 2007.
   
   \bibitem{Knuth} 
   Donald E.~Knuth, {\em The \MF book},  Addison-Wesley,1986.
   
   \bibitem{Saru}  \href{https://backlog.com/ja/git-tutorial/}{サル先生のGit入門}

   \bibitem{Takahashi}  
   高橋 あおい, つくって、壊して、直して学ぶ Git＆GitHub 入門, 翔泳社, 2025.
   
   \bibitem{Wa}  渡辺宙志, ゼロから学ぶGit/GitHub, 講談社, 2024.

   \bibitem{PL} 
   ポンストール, ロリガー, 実用 Git 第3版, オライリー, 2024.
   
   \bibitem{M}
   Mana, 1冊ですべて身につくHTML \verb/&/ CSSとWebデザイン入門講座［第2版］, SBクリエイティブ, 2024.


   

\end{thebibliography}

\printindex

\end{document}

